\chapter{外场、反冲与量子交换}\label{chap:external-pinem}

一个电子穿过受激光照射的纳米结构时，近场可以使它获得或失去能量。若光主要含有角频率 $\omega$，出射能谱会在入射中心能量 $E_0$ 两侧出现间距约为 $\hbar\omega$ 的峰。这类利用光学近场调制并读取电子的实验称为\emph{光子诱导近场电子显微}（photon-induced near-field electron microscopy, PINEM）。用整数 $\ell$ 给这些能量边带编号，$\ell>0$ 表示电子净获得能量，$\ell<0$ 表示净损失能量。上一章已经建立有限时间内的电子--环境跃迁理论，现在用它来求各条边带的振幅和相位。

先考虑光场可由给定经典电磁场描述的情形。此时需要求电子在外场中的波函数，近场的空间分布和光学相位都进入散射振幅。随后再保留光场的产生、湮灭算符，考察光子数涨落以及电子与光的纠缠；这部分通常称为\emph{量子 PINEM}（quantum PINEM, QPINEM）。两种描述的联系由相干光的经典极限给出，而电子反冲和有限脉冲的影响需要在各自模型中检查。\cite{Feist2015Nature,Wang2020Cavity,Dahan2021Science,Ruimy2025Perspective}

\section{外场散射与相位约化}

给定入射波包和外场四势 $A^\mu(x)$，外场散射要求求出电子离开相互作用区时的振幅。我们从 Dirac 方程的传播子开始。当入射动量集中在中心值附近、外场也不明显改变电子前进方向时，可以把波函数分成快速振荡的载波与缓慢变化的包络。再沿电子运动方向求包络的相位，就能把外场作用和能谱边带联系起来。

每一步简化都有自己的条件。正能投影舍去不能真实到达的负能通道；窄带光场选择交换能量；当不同电子通道之间的反冲差别足够小时，最后得到近似均匀的能量梯。下面的关系概括了所用模型之间的联系：
\begin{equation}
\boxed{
\text{多模连续谱}
\longrightarrow
\text{正能窄带电子}
\longrightarrow
\text{单频/单模通道}
\longrightarrow
\text{真实反冲格}
\longrightarrow
\text{均匀能量梯}.}
\label{eq:pinem_reduction_spine}
\end{equation}
这些箭头包含对工作能区和场谱的选取，以及对反冲和耦合差别的近似。常见的 Bessel 边带公式将在均匀梯模型中导出；若不同边的失谐或耦合已明显不同，就需要保留各通道的独立演化方程。

先把共振条件写清楚。一次交换若向电子传递机械动量 $\hbar\boldsymbol\kappa$，电子的能量就从 $E(\mathbf p)$ 变为 $E(\mathbf p+\hbar\boldsymbol\kappa)$。设环境同时提供能量 $\hbar\Omega$，两者之差除以 $\hbar$，定义为\emph{失谐}。由于这里采用真实相对论色散，也称为质量壳失谐：
\begin{equation}
\boxed{
\Delta_{\rm ms}(\mathbf p;\boldsymbol\kappa,\Omega)
\equiv
\frac{
E(\mathbf p+\hbar\boldsymbol\kappa)-E(\mathbf p)
}{\hbar}
-\Omega .}
\label{eq:free_electron_exact_mass_shell_mismatch}
\end{equation}
其中 $\boldsymbol\kappa$ 的单位为 ${\rm m^{-1}}$，$\Omega$ 与 $\Delta_{\rm ms}$ 的单位为 ${\rm s^{-1}}$；$\Omega$ 可取正负，分别表示向电子提供或从电子取走能量。$\Delta_{\rm ms}=0$ 表示能量恰好匹配。根据上一章的有限时间跃迁公式，作用时间为 $T$ 时，
满足 $|\Delta_{\rm ms}|T\lesssim1$ 的通道位于共振主峰附近。对重复的单一交换步长，
\[
\mathbf p_\ell=\mathbf p_0+\ell\hbar\boldsymbol\kappa,
\qquad
\Delta_\ell
=
\Delta_{\rm ms}(\mathbf p_\ell;\boldsymbol\kappa,\omega),
\]
所以“低反冲”真正要求的是实际占据窗口内 $\Delta_\ell$ 和逐边矩阵元的变化小于动力学分辨尺度；它不等价于只要求 $\hbar\omega/E_0$ 很小。反过来，当不同 $\Delta_\ell$ 可被有限时间分辨时，真实反冲会选择少数共振边，平移对称能量梯必须退出。这一质量壳失谐把 PINEM、有限反冲格和少态共振统一成同一个运动学问题。

为计算实际电子的速度和反冲，沿用中心动量 $p_0$，相应总能量、Lorentz 因子与群速度为
\begin{equation*}
 p_0,\qquad E_0=\sqrt{\me^2c^4+c^2p_0^2},\qquad
 \gamma_0=\frac{E_0}{\me c^2},\qquad
 v_0=\frac{c^2p_0}{E_0},
\end{equation*}
纵向二阶色散质量固定写为 $m_\parallel=\gamma_0^3\me$。对选中的交换角频率 $\omega$，理想能量梯用 $|\ell\rangle$ 表示，$\ell\in\mathbb Z$ 是电子相对于中心通道交换的净量子数；真实通道能量始终写 $\mathcal E_\ell=E(p_\ell)$，只有在均匀梯近似成立后才写 $E_\ell^{(0)}=E_0+\ell\hbar\omega$。经典场作用在整个路径上积累的无量纲耦合记为 $\beta$。它将由沿轨迹的场积分确定；$|\beta|$ 控制边带分布的宽度，$\arg\beta$ 控制各边带振幅的相位。光场量子化后的耦合与算符将在量子交换一节引入。

为区别时间域实场与单频复振幅，实时间域电场记为 $\mathbf E(\mathbf r,t)$，单频正频分量记为
\begin{equation*}
 \mathbf E^{(+)}(\mathbf r,t)=\mathbf E_\omega(\mathbf r)\ee^{-\ii\omega t},
 \qquad
 E_{\omega,z}(\mathbf r)\equiv\hat{\mathbf z}\cdot\mathbf E_\omega(\mathbf r).
\end{equation*}
$E_{\omega,z}$ 的单位为 $\mathrm{V\,m^{-1}}$。实场由正频分量及其共轭相加，所以纵向实场峰值为 $2|E_{\omega,z}|$。它与单位为能量的通道能量 $\mathcal E_\ell$ 是不同的量。对一般非单色信号，Fourier 变换仍按全书约定用波浪号表示。

离散通道能否近似成均匀能量梯，首先由电子色散曲率决定。若单次交换对应的纵向动量步长为 $\hbar \kappa_z$，二阶色散自然给出反冲频率与有限时间累积相位
\begin{equation}
\boxed{
\omega_{\rm r}\equiv\frac{\hbar \kappa_z^2}{2m_\parallel},\qquad
\varphi_{\rm r}\equiv\omega_{\rm r}T,}
\label{eq:pinem_recoil_frequency_master}
\end{equation}
其中 $T$ 是实际作用时间。$\omega_{\rm r}$ 是相邻动量步长经纵向色散曲率产生的反冲频率尺度，$\varphi_{\rm r}$ 是作用窗口内累计的反冲相位。同一 $\omega_{\rm r}$ 与 $\varphi_{\rm r}$ 同时控制有限反冲格、自由漂移相位和实验反演中的反冲分辨，因此这些对象使用同一组定义。

电子--量子场联合哈密顿量的经典场极限由\eqnrefs{eq:multimode_coherent_classical_scaling} 给出：在保持平均电磁场有限时令相对量子涨落趋于零，光场对电子便退化为给定的经典背景。该极限只约束光场的相对量子涨落；电子动量宽度、自旋、反冲、场频谱和作用强度仍由电子侧的独立控制参数决定，因此起点仍是完整外场 Dirac 传播子。

\subsection{外场 Dirac 传播子}

给定外场后，电子波函数满足含电磁耦合的 Dirac 方程。这里把光场视为已知函数，电子的自旋和相对论运动仍由四分量波函数描述。采用 SI 四势
\begin{equation*}
A^\mu(x)\equiv\left(\frac{\phi(\mathbf r,t)}{c},\mathbf A(\mathbf r,t)\right)
\end{equation*}
参数化外场，其中 $\phi$ 与 $\mathbf A$ 分别是标势和矢势；电子电荷为 $\qe=-e$。最小耦合 Dirac 哈密顿量由\eqnrefs{eq:dirac-hamiltonian-em} 定义。给定这一哈密顿量，在共同稠密定义域与通常保证非自治演化存在的连续性条件下，定义唯一幺正传播子 $\hat U_A(t_f,t_i)$，满足
\begin{equation}
\boxed{
\ii\hbar\partial_{t_f}\hat U_A(t_f,t_i)
=\hat H_D[A](t_f)\hat U_A(t_f,t_i),
\qquad
\hat U_A(t_i,t_i)=\mathbb I.}
\label{eq:exact_external_field_evolution}
\end{equation}
给定入射态后，作用这个传播子就得到出射态；后面的相位调制公式将由它在窄波包条件下的近似求得。

为了区分自由传播与外场作用，写
\begin{equation*}
\hat H_0=c\boldsymbol\alpha\cdot\hat{\mathbf p}+\beta_D\me c^2,
\qquad
\hat V_A(t)=\qe\phi(\hat{\mathbf r},t)-\qe c\boldsymbol\alpha\cdot\mathbf A(\hat{\mathbf r},t),
\end{equation*}
并以 $t_i$ 为基准定义相互作用绘景中的作用算符
\begin{equation*}
\hat V_I(t)
\equiv
\ee^{+\ii\hat H_0(t-t_i)/\hbar}
\hat V_A(t)
\ee^{-\ii\hat H_0(t-t_i)/\hbar}.
\end{equation*}
Dyson 时间排序指数已经由 Volterra 积分方程的逐次迭代确定。把当前 Dirac 外场 Hamiltonian 代入这一一般传播子结构，严格得到
\begin{equation}
\boxed{
\hat U_A(t_f,t_i)
=\ee^{-\ii\hat H_0(t_f-t_i)/\hbar}\hat U_I(t_f,t_i),
\qquad
\hat U_I
=\mathcal T\exp\!\left[-\frac{\ii}{\hbar}
\int_{t_i}^{t_f}\dd t\,\hat V_I(t)\right].}
\label{eq:exact_external_field_interaction_picture}
\end{equation}
时间排序指数定义完整的有限时传播子；弱作用或短时模型由前文 Dyson/Magnus 展开的截断阶数与相应误差控制量决定。

在坐标表象中定义 Dirac 传播核
\begin{equation}
\boxed{
\Psi_f(\mathbf r_f)
=\int\dd^3 r_i\,
K_A(\mathbf r_f,t_f;\mathbf r_i,t_i)\Psi_i(\mathbf r_i),
\qquad
K_A\equiv\langle\mathbf r_f|\hat U_A|\mathbf r_i\rangle .}
\label{eq:exact_external_field_kernel}
\end{equation}
如果入射与出射区域外场已经消失，则可用自由 Dirac 正能态定义正能散射块
\begin{equation}
\boxed{
\mathcal S^{++}_{s_fs_i}(\mathbf p_f,\mathbf p_i;t_f,t_i)
=\langle\mathbf p_f,s_f,+|\hat U_A|\mathbf p_i,s_i,+\rangle .}
\label{eq:exact_positive_energy_external_scattering_kernel}
\end{equation}
任意初始动量--自旋振幅因此变成
\begin{equation}
\boxed{
a^{(f)}_{s_f}(\mathbf p_f)
=\sum_{s_i}\int\dd^3 p_i\,
\mathcal S^{++}_{s_fs_i}(\mathbf p_f,\mathbf p_i)
a^{(i)}_{s_i}(\mathbf p_i).}
\label{eq:exact_external_field_wavepacket_scattering}
\end{equation}
$\mathcal S^{++}$ 只是完整幺正传播子的一个子块，一般并不单独幺正。对正能初态，概率缺口恰好等于传播到负能子块的概率；若这一泄漏不可忽略，就必须保留完整 $\mathcal S^{\sigma_f\sigma_i}$，而不能把缺口解释成唯象损耗。

同一结构同时固定轨迹相位的规范变换性质。采用
\begin{equation}
\boxed{
A_\mu' = A_\mu-\partial_\mu\chi,
\qquad
\Psi' = \ee^{\ii\qe\chi/\hbar}\Psi.}
\label{eq:external_field_gauge_transformation}
\end{equation}
完整传播核满足端点协变律
\begin{equation}
\boxed{
K_{A-\partial\chi}(x_f,x_i)
=\ee^{\ii\qe\chi(x_f)/\hbar}
K_A(x_f,x_i)
\ee^{-\ii\qe\chi(x_i)/\hbar}.}
\label{eq:exact_dirac_kernel_gauge_covariance}
\end{equation}
\eqnrefs{eq:exact_dirac_kernel_gauge_covariance} 已固定传播核在规范变换下的端点相位。对传播核做正能、慢包络或程函约化时，这个端点变换律必须由约化核继承，因此 Wilson 线与程函相位的规范因子由母传播核的协变性唯一决定。任何降维都必须同时记录被投影的 Hilbert 空间扇区、被舍弃的谱相位以及相应无量纲控制参数。


端点相位律要求任何由同一外场 Dirac 动力学约化得到的传播核继承相同的规范协变结构；Wilson 线与程函近似中的端点相位正是这一要求的具体表现。

\begin{derivation}[从\eqnrefs{eq:exact_positive_energy_external_scattering_kernel}到\eqnrefs{eq:exact_external_field_wavepacket_scattering}]
因此\eqnrefs{eq:exact_external_field_wavepacket_scattering} 只需由完备关系投影得到； 忽略负能扇区则要求相应泄漏概率可忽略。 自由 Dirac 哈密顿量的广义本征态按正、负能两支完备。 采用连续归一化，先写

\begin{equation*}
\mathbb I
=\sum_s\int\dd^3 p\,
\Bigl(
|\mathbf p,s,+\rangle\langle\mathbf p,s,+|
+|\mathbf p,s,-\rangle\langle\mathbf p,s,-|
\Bigr)
\equiv P_++P_-.
\end{equation*}
一个纯正能入射波包因此是
\begin{equation*}
|\Psi_i\rangle
=\sum_{s_i}\int\dd^3 p_i\,
 a^{(i)}_{s_i}(\mathbf p_i)|\mathbf p_i,s_i,+\rangle,
\qquad
\sum_{s_i}\int\dd^3 p_i\,
|a^{(i)}_{s_i}(\mathbf p_i)|^2=1.
\end{equation*}
完整外场传播后插入末态完备关系，并投影到正能支：
\begin{align*}
P_+|\Psi_f\rangle
&=P_+U_AP_+|\Psi_i\rangle\\
&=\sum_{s_f}\int\dd^3 p_f\,
|\mathbf p_f,s_f,+\rangle
\langle\mathbf p_f,s_f,+|U_AP_+|\Psi_i\rangle.
\end{align*}
把入射波包展开代回最后一个矩阵元，得到
\begin{align*}
\langle\mathbf p_f,s_f,+|U_A|\Psi_i\rangle
&=\sum_{s_i}\int\dd^3 p_i\,
\langle\mathbf p_f,s_f,+|U_A|\mathbf p_i,s_i,+\rangle
 a^{(i)}_{s_i}(\mathbf p_i)\\
&=\sum_{s_i}\int\dd^3 p_i\,
\mathcal S^{++}_{s_fs_i}(\mathbf p_f,\mathbf p_i)
 a^{(i)}_{s_i}(\mathbf p_i),
\end{align*}
这就是\eqnrefs{eq:exact_external_field_wavepacket_scattering}。积分测度、Dirac $\delta$ 归一化与自旋求和均由自由态基底固定，因此不会额外产生唯象态密度因子。再说明为什么 $S^{++}$ 一般不单独幺正。完整传播子满足 $U_A^\dagger U_A=\mathbb I$，在正能入射子空间两侧夹上 $P_+$：
\begin{align*}
P_+
&=P_+U_A^\dagger(P_++P_-)U_AP_+\\
&=(P_+U_A^\dagger P_+)(P_+U_AP_+)
 +(P_+U_A^\dagger P_-)(P_-U_AP_+).
\end{align*}
于是作为正能子空间上的算符恒等式，
\begin{equation*}
(S^{++})^\dagger S^{++}+(S^{-+})^\dagger S^{-+}=I_+.
\end{equation*}
对任意归一化正能初态取期望值，就得到
\begin{equation*}
\|S^{++}\psi\|^2
=1-\|S^{-+}\psi\|^2.
\end{equation*}
\end{derivation}
正能散射概率的缺口严格等于负能扇区泄漏；只有当 $P_-U_AP_+$ 在实际波包上可忽略时，$S^{++}$ 才能作为近似幺正的有效电子散射算符。

\begin{derivation}[从\eqnrefs{eq:external_field_gauge_transformation}到\eqnrefs{eq:exact_dirac_kernel_gauge_covariance}]

先在算符层面证明，而不是直接比较两个积分核。外场 Dirac 方程写成
\[
\ii\hbar\partial_t|\Psi(t)\rangle
=
\hat H[A](t)|\Psi(t)\rangle .
\]
规范变换取正文约定
\[
A_\mu'(x)=A_\mu(x)-\partial_\mu\chi(x),
\qquad
|\Psi'(t)\rangle
=
\hat G_\chi(t)|\Psi(t)\rangle,
\]
其中在坐标表象
\[
\langle\mathbf r|\hat G_\chi(t)|\Psi\rangle
=
\ee^{+\ii\qe\chi(\mathbf r,t)/\hbar}
\Psi(\mathbf r,t).
\]
对变换后的态求时间导数：
\begin{align*}
\ii\hbar\partial_t|\Psi'\rangle
&=
\ii\hbar(\partial_t\hat G_\chi)|\Psi\rangle
+\hat G_\chi\hat H[A]|\Psi\rangle\\
&=
\left[
\hat G_\chi\hat H[A]\hat G_\chi^\dagger
+\ii\hbar(\partial_t\hat G_\chi)\hat G_\chi^\dagger
\right]
|\Psi'\rangle.
\end{align*}
因此变换后的 Hamiltonian 必须满足
\[
\hat H[A']
=
\hat G_\chi\hat H[A]\hat G_\chi^\dagger
+\ii\hbar(\partial_t\hat G_\chi)\hat G_\chi^\dagger.
\]
第二项正是时间依赖幺正变换中必需的附加项；若只写
\(\hat H'\!=\hat G\hat H\hat G^\dagger\)，会漏掉标势的规范变化。

现在令 \(\hat U_A(t,t_i)\) 是原 Hamiltonian 的时间序传播算符，
\[
\ii\hbar\partial_t\hat U_A(t,t_i)
=
\hat H[A](t)\hat U_A(t,t_i),
\qquad
\hat U_A(t_i,t_i)=I.
\]
定义
\[
\widetilde U(t,t_i)
\equiv
\hat G_\chi(t)\hat U_A(t,t_i)\hat G_\chi^\dagger(t_i).
\]
对 \(t\) 求导，
\begin{align*}
\ii\hbar\partial_t\widetilde U
&=
\left[
\ii\hbar(\partial_t\hat G_\chi)\hat G_\chi^\dagger
+\hat G_\chi\hat H[A]\hat G_\chi^\dagger
\right]\widetilde U\\
&=
\hat H[A'](t)\widetilde U.
\end{align*}
同时
\[
\widetilde U(t_i,t_i)=I.
\]
由传播子初值问题的唯一性，
\[
\boxed{
\hat U_{A-\partial\chi}(t_f,t_i)
=
\hat G_\chi(t_f)
\hat U_A(t_f,t_i)
\hat G_\chi^\dagger(t_i).
}
\]

最后取坐标表象。定义
\[
K_A(x_f,x_i)
\equiv
\langle\mathbf r_f|
\hat U_A(t_f,t_i)
|\mathbf r_i\rangle .
\]
利用 \(\hat G_\chi\) 在坐标表象中为乘法算符，
\begin{align*}
K_{A-\partial\chi}(x_f,x_i)
&=
\langle\mathbf r_f|
\hat G_\chi(t_f)\hat U_A
\hat G_\chi^\dagger(t_i)
|\mathbf r_i\rangle\\
&=
\ee^{+\ii\qe\chi(x_f)/\hbar}
K_A(x_f,x_i)
\ee^{-\ii\qe\chi(x_i)/\hbar}.
\end{align*}
因此
\[
\boxed{
K_{A-\partial\chi}(x_f,x_i)
=
\ee^{\ii\qe[\chi(x_f)-\chi(x_i)]/\hbar}
K_A(x_f,x_i),
}
\]
即正文\eqnrefs{eq:exact_dirac_kernel_gauge_covariance}。
\end{derivation}

\subsection{正能包络与程函近似}

对于动量集中在 $\mathbf p_0$ 附近的入射波包，大部分空间振荡来自中心平面波。把这一已知振荡因子提出后，剩下的包络变化较慢。若外场主要改变沿途相位，而对前进方向和自旋的影响较小，包络方程可以进一步简化。这种沿近似轨迹积累相位的方法称为\emph{程函近似}（eikonal approximation）。下面从 Dirac 方程检验它的条件，并求出包络满足的方程。

这里的\emph{程函近似}（eikonal approximation）是一种高能、慢变的输运约化：把完整波动方程中快速变化的自由载波显式提出以后，只保留慢包络沿中心群速度轨迹的一阶输运；外场的主要作用因而表现为沿该轨迹积累的相位。它保留相对论色散本身；受控忽略的是自旋子随动量的显著变化、强横向偏转、高阶色散以及正负能扇区混合。相应控制参数先逐一定义，再由传播子误差界判断这一约化在给定电子束和作用长度下是否成立。

这个慢变量就是去掉快速自由载波后的电子包络。设入射波包以 $\mathbf p_0$ 为中心。自由色散关系
\[
E(\mathbf p)=\sqrt{\me^2c^4+c^2\mathbf p^2}
\]
在 $\mathbf p_0$ 附近的一阶色散梯度给出群速度，因此决定包络中心的平移；二阶及更高阶导数分别产生波包展宽和更高阶畸变。在外场引起的中心动量偏移与自旋旋转均由相应小参数控制时，最低阶输运方程沿群速度轨迹积分，得到经典轨迹上的相位累积。

对窄正能波包，以中心动量 $\mathbf p_0$ 的正能自由旋量作为局域载波基底，并把快速自由相位与慢包络分开：
\begin{equation*}
\Psi(\mathbf r,t)
=u_0\,\fenv(\mathbf r,t)
\exp\!\left[\frac{\ii}{\hbar}(\mathbf p_0\cdot\mathbf r-E_0t)\right]
+\delta\Psi,
\end{equation*}
其中 $E_0=\sqrt{\me^2c^4+c^2p_0^2}$，$u_0$ 是归一化正能 Dirac 自旋子，$\fenv$ 是电子包络，$\delta\Psi$ 收集自旋子动量依赖、色散高阶项和外场诱导的负能分量。自由正能流给出群速度
\begin{equation*}
\mathbf v_0=u_0^\dagger c\boldsymbol\alpha u_0
=\nabla_{\mathbf p}E|_{\mathbf p_0}
=\frac{c^2\mathbf p_0}{E_0}.
\end{equation*}

固定轨迹程函近似只在四类变化都慢时成立：正能自旋子在电子带宽内近似不变、横向偏转小、总动量改变量小、背景和包络都缓于 de Broglie 尺度。令 $\Delta p$ 是初始动量带宽，$\Delta p_{\rm int}$ 是相互作用诱导的最大动量变化，$L_A,L_f$ 分别是外场和包络的最短变化尺度。再令 $u_s(\mathbf p)$ 为工作偏振的归一化正能旋量，并定义与任意动量依赖相位无关的秩一投影 $\Pi_s(\mathbf p)\equiv u_s(\mathbf p)u_s^\dagger(\mathbf p)$。于是旋量基底变化可以严格量化为
\begin{align*}
\epsilon_u&\equiv
\sup_{|\mathbf p-\mathbf p_0|\le\Delta p}
\|\Pi_s(\mathbf p)-\Pi_s(\mathbf p_0)\|,
&\epsilon_\perp&\equiv\frac{|\Delta\mathbf p_{\perp,{\rm int}}|}{p_0},\\
\epsilon_p&\equiv\frac{\Delta p_{\rm int}}{p_0},
&\epsilon_{\rm WKB}&\equiv\max\!\left(\frac1{\kappa_{\rm e0}L_A},\frac1{\kappa_{\rm e0}L_f}\right).
\end{align*}
其中 $\kappa_{\rm e0}=p_0/\hbar$。方便的局域充分条件是
\begin{equation}
\boxed{
\epsilon_{\rm eik}
\equiv\max(\epsilon_u,\epsilon_\perp,\epsilon_p,\epsilon_{\rm WKB})\ll1.}
\label{eq:general_eikonal_control_parameter}
\end{equation}
这还必须配合正负能扇区之间的实际矩阵元远小于约 $2\me c^2$ 的谱隙；“电子很快”本身不能替代这一检查。

更可靠的判据是传播子误差。令 $P_+$、$P_-$ 是自由 Dirac 正负能投影，$P_{\mathcal W}\le P_+$ 选取实际电子窄带；以线性化正能包络生成元和自由负能块构造块对角近似哈密顿量$H_{\rm app}$，并定义
\begin{equation*}
R_{\rm eik}(t)=H_D[A](t)-H_{\rm app}(t).
\end{equation*}
则完整 Dirac 传播子 $U_A$ 与近似传播子 $U_{\rm app}$ 满足
\begin{equation}
\boxed{
\|[U_A-U_{\rm app}]P_{\mathcal W}\|
\le
\epsilon_{\rm eik}^{\rm op}
\equiv\frac1\hbar\int_{t_i}^{t_f}\dd t\,
\|R_{\rm eik}(t)U_{\rm app}(t,t_i)P_{\mathcal W}\|.}
\label{eq:eikonal_duhamel_error_bound}
\end{equation}
因此真正需要小的是累计余项 $\epsilon_{\rm eik}^{\rm op}$；\eqnrefs{eq:general_eikonal_control_parameter} 只是对它的局域尺度估计。

\begin{derivation}[从\eqnrefs{eq:duhamel_master_identity}到\eqnrefs{eq:eikonal_duhamel_error_bound}]
从完整最小耦合 Dirac 方程出发，先把快速自由载波与慢包络分开。取
\[
\Psi(\mathbf r,t)
=u_0\,\fenv(\mathbf r,t)
\exp\!\left[\frac{\ii}{\hbar}(\mathbf p_0\cdot\mathbf r-E_0t)\right]
+\delta\Psi,
\]
其中自由正能旋量满足
\[
\left(c\boldsymbol\alpha\cdot\mathbf p_0+\beta_D\me c^2\right)u_0=E_0u_0,
\qquad u_0^\dagger u_0=1.
\]
对载波求导时
\begin{align*}
\ii\hbar\partial_t(\fenv\ee^{\ii(\mathbf p_0\cdot\mathbf r-E_0t)/\hbar})
&=\left(E_0\fenv+\ii\hbar\partial_t\fenv\right)
\ee^{\ii(\mathbf p_0\cdot\mathbf r-E_0t)/\hbar},\\
\hat{\mathbf p}(\fenv\ee^{\ii(\mathbf p_0\cdot\mathbf r-E_0t)/\hbar})
&=\left(\mathbf p_0\fenv-\ii\hbar\nabla\fenv\right)
\ee^{\ii(\mathbf p_0\cdot\mathbf r-E_0t)/\hbar}.
\end{align*}
代入
\[
\ii\hbar\partial_t\Psi=
\left[c\boldsymbol\alpha\cdot(\hat{\mathbf p}-\qe\mathbf A)
+\beta_D\me c^2+\qe\phi\right]\Psi
\]
以后，自由本征方程把所有与 $E_0$ 和 $\mathbf p_0$ 成正比的载波项精确抵消，得到
\[
\ii\hbar u_0\partial_t\fenv
=-\ii\hbar c\boldsymbol\alpha\cdot(\nabla\fenv)u_0
+\qe(\phi-c\boldsymbol\alpha\cdot\mathbf A)u_0\fenv
+\mathcal R_D.
\]
$\mathcal R_D$ 收集没有进入固定旋量、固定中心轨迹一阶输运方程的项，包括旋量随动量的变化、横向偏转、高阶色散与正负能混合。左乘 $u_0^\dagger$，并使用自由 Dirac 流恒等式
\[
u_0^\dagger c\boldsymbol\alpha u_0=\mathbf v_0,
\qquad
\mathbf v_0=\nabla_{\mathbf p}E(\mathbf p_0),
\]
可得
\[
\ii\hbar(\partial_t+\mathbf v_0\cdot\nabla)\fenv
=\qe(\phi-\mathbf v_0\cdot\mathbf A)\fenv
+u_0^\dagger\mathcal R_D.
\]
把最后一项删除所定义的近似生成元记为 $H_{\rm app}(t)$，完整 Dirac 生成元与它的差为
\[
R_{\rm eik}(t)=H_D[A](t)-H_{\rm app}(t).
\]
令 $U_A$ 与 $U_{\rm app}$ 分别满足
\begin{align*}
\ii\hbar\partial_tU_A(t,t_i)&=H_D[A](t)U_A(t,t_i),\\
\ii\hbar\partial_tU_{\rm app}(t,t_i)&=H_{\rm app}(t)U_{\rm app}(t,t_i).
\end{align*}
为直接比较两个传播子，定义
\[
F(t)=U_A(t_f,t)U_{\rm app}(t,t_i).
\]
反向传播子的导数满足
\[
\partial_tU_A(t_f,t)=\frac{\ii}{\hbar}U_A(t_f,t)H_D[A](t),
\]
所以乘积求导给
\begin{align*}
\dot F(t)
&=\frac{\ii}{\hbar}U_A(t_f,t)H_D[A](t)U_{\rm app}(t,t_i)
-\frac{\ii}{\hbar}U_A(t_f,t)H_{\rm app}(t)U_{\rm app}(t,t_i)\\
&=\frac{\ii}{\hbar}U_A(t_f,t)R_{\rm eik}(t)U_{\rm app}(t,t_i).
\end{align*}
端点分别为
\[
F(t_f)=U_{\rm app}(t_f,t_i),
\qquad
F(t_i)=U_A(t_f,t_i).
\]
因此
\[
U_A(t_f,t_i)-U_{\rm app}(t_f,t_i)
=-\frac{\ii}{\hbar}\int_{t_i}^{t_f}\dd t\,
U_A(t_f,t)R_{\rm eik}(t)U_{\rm app}(t,t_i).
\]
右乘实际工作窗投影 $P_{\mathcal W}$，取算符范数并使用
$\|U_A(t_f,t)\|=1$ 与范数次乘性，得到
\begin{align*}
\|[U_A-U_{\rm app}]P_{\mathcal W}\|
&\le\frac1\hbar\int_{t_i}^{t_f}\dd t\,
\|U_A(t_f,t)R_{\rm eik}(t)U_{\rm app}(t,t_i)P_{\mathcal W}\|\\
&\le\frac1\hbar\int_{t_i}^{t_f}\dd t\,
\|R_{\rm eik}(t)U_{\rm app}(t,t_i)P_{\mathcal W}\|,
\end{align*}
即正文\eqnrefs{eq:eikonal_duhamel_error_bound}。
\end{derivation}
\subsection{轨迹相位与频谱卷积}

在程函近似误差受控以后，正能包络沿经典特征线 $\mathbf r_e(t)=\mathbf r_{e0}+\mathbf v_0t$ 满足一阶标量方程
\begin{equation}
\boxed{
\ii\hbar\frac{\dd \fenv}{\dd t}
=\qe[\phi(\mathbf r_e,t)-\mathbf v_0\cdot\mathbf A(\mathbf r_e,t)]\fenv.}
\label{eq:general_gauge_covariant_eikonal}
\end{equation}
它的解是沿电子世界线的阿贝尔 Wilson 相位
\begin{equation}
\boxed{
\frac{f_{\mathrm e,f}}{f_{\mathrm e,i}}
=\exp\!\left[-\frac{\ii \qe}{\hbar}
\int_{\mathcal C}\dd t\,(\phi-\mathbf v_0\cdot\mathbf A)\right].}
\label{eq:general_eikonal_wilson_phase}
\end{equation}
其中 $\mathcal C$ 表示实际电子轨迹。规范变换只改变端点相位；入射、出射基底作相同规范变换后，散射概率保持不变。

令 $\zeta$ 标记电子切片，定义轨迹相位 $\varphi_{\rm eik}(\zeta)$ 和乘法算符 $\mathcal U_{\rm eik}(\zeta)=\ee^{\ii\varphi_{\rm eik}(\zeta)}$。坐标空间的纯相位乘法在动量空间成为卷积：
\begin{equation*}
\widetilde f_{\mathrm e,f}(\kappa)=\int\dd \kappa'\,\mathcal K_{\rm eik}(\kappa-\kappa')\widetilde f_{\mathrm e,i}(\kappa'),
\qquad
\mathcal K_{\rm eik}(\Delta\kappa)=\frac1{2\pi}\int\dd\zeta\,\ee^{-\ii\Delta\kappa\zeta}\mathcal U_{\rm eik}(\zeta).
\end{equation*}
若 $\mathcal U_{\rm eik}$ 对纵向坐标具有周期 $2\pi/\kappa_{\rm g}$，写傅里叶系数 $A_\ell$ 后得到离散平移律
\begin{equation}
\boxed{
\widetilde f_{\mathrm e,f}(\kappa)=\sum_{\ell\in\mathbb Z}A_\ell\widetilde f_{\mathrm e,i}(\kappa-\ell\kappa_{\rm g}).}
\label{eq:general_periodic_eikonal_discrete_convolution}
\end{equation}
周期近场因此不是“产生新的电子能级”，而是其傅里叶谐波选择了一组相差 $\hbar\kappa_{\rm g}$ 的自由电子动量副本。

要把这些副本解释成互斥边带概率，还必须检查它们近似正交。对归一化入射谱定义
\begin{equation}
\boxed{
\epsilon_{\rm ov}
\equiv\sup_{\ell\neq m}
\frac{\left|\int\frac{\dd \kappa}{2\pi}\,
\widetilde f_{\mathrm e,i}^*(\kappa-\ell\kappa_{\rm g})\widetilde f_{\mathrm e,i}(\kappa-m\kappa_{\rm g})\right|}
{\int\frac{\dd \kappa}{2\pi}|\widetilde f_{\mathrm e,i}(\kappa)|^2}\ll1.}
\label{eq:translated_packet_overlap_condition}
\end{equation}
只有在这一条件成立时，$|A_\ell|^2$ 才可直接解释为分辨的边带权重；否则不同平移副本之间的交叉项就是可观测相干。

周期纯相位散射还具有一个比 Bessel 归一化更一般的幺正约束。把
\[
\mathcal U_{\rm eik}(\zeta)
=\sum_{n\in\mathbb Z}A_n\ee^{+\ii n\kappa_{\rm g}\zeta}
\]
看作单位圆上的 Fourier 级数，并在理想离散通道空间 $\ell^2(\mathbb Z)$ 上定义 Toeplitz 散射矩阵
\begin{equation}
\boxed{
S_{\ell m}=A_{\ell-m}.}
\label{eq:periodic_phase_toeplitz_matrix}
\end{equation}
由于纯相位因子逐点满足 $|\mathcal U_{\rm eik}(\zeta)|=1$，其 Fourier 系数并非任意，而必须满足
\begin{equation}
\boxed{
\sum_{n\in\mathbb Z}A_{n+r}A_n^*=\delta_{r0},
\qquad r\in\mathbb Z.}
\label{eq:periodic_phase_fourier_unitarity_identity}
\end{equation}
这正是 $S^\dagger S=SS^\dagger=I$ 的矩阵元条件。特别地，$r=0$ 给出
$\sum_n|A_n|^2=1$；$r\neq0$ 则给出不同列之间的正交性。因此任意周期纯相位调制都定义一个严格幺正的双边通道散射，而 Bessel 恒等式
$\sum_nJ_n^2(x)=1$ 只是单谐波相位的特例。若实际平移波包不满足\eqnrefs{eq:translated_packet_overlap_condition}，不能把 $|A_n|^2$ 当成互斥概率，但完整连续波函数的范数仍由这个母幺正性严格保持。

对单色正频场，轨迹相位还可以改写为电场的同步 Fourier 投影。写成
$\mathbf E^{(+)}(\mathbf r,t)=\mathbf E_\omega(\mathbf r)\ee^{-\ii\omega t}$，并定义
$\mathbf E_\omega=-\nabla\phi_\omega+\ii\omega\mathbf A_\omega$。这样，轨迹上的势组合可以分成纵向电场积分与端点势差。散射边界条件要求电子进入和离开作用区时外场趋于零，端点项随之消失。因此 PINEM 累积耦合统一由沿电子轨迹取样的纵向单频复幅 $E_{\omega,z}$ 表示，其空间依赖同时给出后续所需的场包络。

实际连续电子波包若不同平移副本重叠，则离散通道基不再严格正交，应回到正文\eqnrefs{eq:general_periodic_eikonal_discrete_convolution} 在振幅层保留交叉项；连续 Hilbert 空间中的母传播子仍保持幺正。

\begin{derivation}[从\eqnrefs{eq:general_eikonal_wilson_phase}到\eqnrefs{eq:general_periodic_eikonal_discrete_convolution}]
其一般系数给出正文\eqnrefs{eq:general_periodic_eikonal_discrete_convolution}，而单色正弦相位的 Bessel 系数只是这一 Fourier 系数的特例。 沿固定特征线 $\mathbf r_e(t)=\mathbf r_{e0}+\mathbf v_0t$，链式法则给

\begin{equation*}
\frac{\dd}{\dd t}\fenv[\mathbf r_e(t),t]
=(\partial_t+\mathbf v_0\cdot\nabla)\fenv.
\end{equation*}
因此正能包络方程化为一阶常微分方程
\begin{equation*}
\frac{1}{\fenv}\frac{\dd \fenv}{\dd t}
=-\frac{\ii \qe}{\hbar}
[\phi(\mathbf r_e,t)-\mathbf v_0\cdot\mathbf A(\mathbf r_e,t)].
\end{equation*}
把两边对时间积分，左边是对数的全微分：
\begin{align*}
\int_{t_i}^{t_f}\dd t\,\frac1\fenv\frac{\dd \fenv}{\dd t}
&=\ln\frac{\fenv(t_f)}{\fenv(t_i)},\\
&=-\frac{\ii \qe}{\hbar}\int_{t_i}^{t_f}\dd t\,
[\phi(\mathbf r_e,t)-\mathbf v_0\cdot\mathbf A(\mathbf r_e,t)].
\end{align*}
指数化后得到沿电子世界线的 Wilson 相位因子
\begin{equation*}
\frac{\fenv(t_f)}{\fenv(t_i)}
=\exp\!\left\{-\frac{\ii \qe}{\hbar}\int_{t_i}^{t_f}\dd t\,
[\phi-\mathbf v_0\cdot\mathbf A]\right\}.
\end{equation*}
若作规范变换
$A_\mu\to A_\mu-\partial_\mu\chi$，指数中新增
\begin{align*}
-\frac{\ii \qe}{\hbar}
\int_{t_i}^{t_f}\dd t\,
[-\partial_t\chi-\mathbf v_0\cdot\nabla\chi]
&=\frac{\ii \qe}{\hbar}
\int_{t_i}^{t_f}\dd t\,\frac{\dd}{\dd t}\chi[\mathbf r_e(t),t]\\
&=\frac{\ii \qe}{\hbar}[\chi(x_f)-\chi(x_i)],
\end{align*}
正好重现完整 Dirac 传播核的端点协变因子。对单色正频分量写
\begin{equation*}
\phi^{(+)}=\phi_\omega(\mathbf r)\ee^{-\ii\omega t},
\qquad
\mathbf A^{(+)}=\mathbf A_\omega(\mathbf r)\ee^{-\ii\omega t}.
\end{equation*}
沿直线轨迹取 $z=v_0t$。利用
\begin{equation*}
E_{\omega,z}=-\partial_z\phi_\omega+\ii\omega A_{\omega z}
\end{equation*}
以及
\begin{equation*}
\frac{\dd}{\dd z}
\left[\phi_\omega(z)\ee^{-\ii\omega z/v_0}\right]
=(\partial_z\phi_\omega)\ee^{-\ii\omega z/v_0}
-\frac{\ii\omega}{v_0}\phi_\omega\ee^{-\ii\omega z/v_0},
\end{equation*}
可把势组合整理成
\begin{align*}
\int\dd z\,
\left(\frac{\phi_\omega}{v_0}-A_{\omega z}\right)
\ee^{-\ii\omega z/v_0}
&=\frac{1}{\ii\omega}
\int\dd z\,E_{\omega,z}(z)
\ee^{-\ii\omega z/v_0}\\
&\quad
-\frac{1}{\ii\omega}
\left[\phi_\omega(z)\ee^{-\ii\omega z/v_0}\right]_{-\infty}^{+\infty}.
\end{align*}
对散射问题，作用区外场消失，端点项为零；这就说明同步纵向复幅 $E_{\omega,z}$ 为何给出规范无关的实际耦合来源。最后推导周期相位的动量卷积。若 $\mathcal U_{\rm eik}(\zeta)=\ee^{\ii\varphi_{\rm eik}(\zeta)}$，则
\begin{align*}
\widetilde f_{\mathrm e,f}(\kappa)
&=\int\dd\zeta\,\ee^{-\ii \kappa\zeta}
\mathcal U_{\rm eik}(\zeta)f_{\mathrm e,i}(\zeta)\\
&=\int\frac{\dd \kappa'}{2\pi}
\left[\int\dd\zeta\,
 \ee^{-\ii(\kappa-\kappa')\zeta}\mathcal U_{\rm eik}(\zeta)\right]
\widetilde f_{\mathrm e,i}(\kappa').
\end{align*}
若 $\mathcal U_{\rm eik}$ 的周期为 $2\pi/\kappa_{\rm g}$，展开
\begin{equation*}
\mathcal U_{\rm eik}(\zeta)=\sum_{\ell\in\mathbb Z}A_\ell \ee^{\ii\ell\kappa_{\rm g}\zeta},
\qquad
A_\ell=\frac{\kappa_{\rm g}}{2\pi}
\int_0^{2\pi/\kappa_{\rm g}}\dd\zeta\,
\mathcal U_{\rm eik}(\zeta)\ee^{-\ii\ell\kappa_{\rm g}\zeta}.
\end{equation*}
代回 Fourier 积分后出现
$2\pi\delta(\kappa-\kappa'-\ell\kappa_{\rm g})$，执行 $\kappa'$ 积分便得到\eqnrefs{eq:general_periodic_eikonal_discrete_convolution}。由于 $|\mathcal U_{\rm eik}|=1$，Parseval 还给
\begin{equation*}
\sum_\ell|A_\ell|^2=1,
\end{equation*}
\end{derivation}
在平移副本互相正交时，Parseval 等式说明这些 Fourier 权重组成归一化的边带概率分布。

\begin{derivation}[从\eqnrefs{eq:general_periodic_eikonal_discrete_convolution}到\eqnrefs{eq:periodic_phase_fourier_unitarity_identity}]
令无量纲周期相位变量为 $\theta=\kappa_{\rm g}\zeta$，并把一个周期上的纯相位函数写成
\[
U(\theta)=\ee^{\ii\varphi(\theta)}
=\sum_{n\in\mathbb Z}A_n\ee^{+\ii n\theta},
\qquad
A_n=\frac1{2\pi}\int_{-\pi}^{\pi}\dd\theta\,
U(\theta)\ee^{-\ii n\theta}.
\]
因为 $\varphi(\theta)$ 为实函数，逐点有
\[
U(\theta)U^*(\theta)=1.
\]
把两边都展开成 Fourier 级数：
\begin{align*}
1
&=\left(\sum_nA_n\ee^{+\ii n\theta}\right)
  \left(\sum_mA_m^*\ee^{-\ii m\theta}\right)\\
&=\sum_{n,m}A_nA_m^*\ee^{+\ii(n-m)\theta}.
\end{align*}
令 $r=n-m$，即 $n=m+r$，则
\[
1
=\sum_{r\in\mathbb Z}
\left(\sum_mA_{m+r}A_m^*\right)
\ee^{+\ii r\theta}.
\]
常数函数 $1$ 的 Fourier 系数只有零阶非零，因此逐阶比较得到
\[
\sum_mA_{m+r}A_m^*=\delta_{r0},
\]
即正文\eqnrefs{eq:periodic_phase_fourier_unitarity_identity}。

现在验证它与离散通道散射矩阵的幺正性完全等价。由
$S_{\ell m}=A_{\ell-m}$，有
\begin{align*}
(S^\dagger S)_{mn}
&=\sum_\ell S_{\ell m}^*S_{\ell n}\\
&=\sum_\ell A_{\ell-m}^*A_{\ell-n}.
\end{align*}
令 $q=\ell-m$，则 $\ell-n=q+m-n$，于是
\[
(S^\dagger S)_{mn}
=\sum_q A_q^*A_{q+m-n}
=\delta_{mn}.
\]
同理 $SS^\dagger=I$。因此任意实周期相位都生成一个双边 Toeplitz 幺正算符；这一结论与相位是否正弦、驱动是否弱、边带数是否有限都无关。

若入射态在离散通道基中为 $|\psi_i\rangle=\sum_m c_m|m\rangle$，输出为
\[
c_\ell'=
\sum_mA_{\ell-m}c_m.
\]
幺正性立即给
\[
\sum_\ell|c_\ell'|^2
=\sum_m|c_m|^2.
\]
对于单通道入射 $c_m=\delta_{m0}$，进一步得到 $P_\ell=|A_\ell|^2$。
\end{derivation}
\subsection{低反冲平移极限}

现在让沿轨迹的场主要以角频率 $\omega$ 振荡。相位调制随电子到达时刻周期变化，它的 Fourier 系数就是能量边带振幅。为了使用同一个系数描述每次交换，需要各条边的耦合近似相同，且反冲造成的相对相位在作用时间内足够小。保留光脉冲的有限包络，把正频纵向场写成
\begin{equation*}
\hat{\mathbf z}\cdot\mathbf E^{(+)}(z,t)
=E_{\omega,z}(z)f(t-\tau)\ee^{-\ii\omega t},
\end{equation*}
其中 $E_{\omega,z}(z)$ 是载频 $\omega$ 的纵向复场幅，$f$ 是无量纲慢变包络，$\tau$ 是场脉冲相对于电子的时间延迟。沿中心速度 $v_0$ 的电子轨迹积累的无量纲单频耦合定义为
\begin{equation}
\boxed{
\beta(\zeta,\tau)
=\frac{\qe}{\hbar\omega}
\int_{-\infty}^{\infty}\dd z\,
E_{\omega,z}(z)
 f\!\left(\frac{z-\zeta}{v_0}-\tau\right)
\ee^{-\ii\omega z/v_0}.}
\label{eq:pinem_standard_beta_pulsed}
\end{equation}
$\zeta$ 标记电子包络内部的纵向切片。$|\beta|$ 衡量沿轨迹相干累积的交换强度，$\arg\beta$ 则保存该切片相对于光学载波的相位。这个积分只使用单频纵向场与给定电子轨迹；代入低反冲平移模型后便得到 PINEM 参数化。

\begin{exbox}[PINEM 低反冲平移极限]
对准直电子，若采用辐射规范并在一维纵向几何中把 $A^2$、横向偏转作为高阶项，常用充分条件是
\begin{equation}
\boxed{
\frac{|\qe A_z|}{p_0}\ll1,
\qquad
\frac{\Delta p_\perp}{p_0}\ll1.}
\label{eq:pinem_ap_small_parameters}
\end{equation}
但高速电子的主方程仍来自 Dirac 正能投影，而不是非相对论 Schr\"odinger 方程。对 $\phi=0$、$\mathbf v_0=v_0\hat z$，
\begin{equation}
\boxed{
\ii\hbar(\partial_t+v_0\partial_z)\fenv
=-\qe v_0A_z(z,t)\fenv+O(\epsilon_{\rm eik}).}
\label{eq:dirac_projected_pinem_envelope}
\end{equation}

正文\eqnrefs{eq:pinem_standard_beta_pulsed} 已把有限脉冲与纵向场几何压缩为切片依赖的复耦合 $\beta(\zeta,\tau)$。要进一步得到单一 Bessel 边带序列，还需要电子包络内部不同切片看到近似相同的 $\beta$。
令入射电子纵向包络的显著支持宽度为 $\Delta\zeta_e$。对时间包络 $f(t)$ 定义变化时间
\begin{equation*}
 \tau_f\equiv\frac{\|f\|_\infty}{\|\dot f\|_\infty},
 \qquad
 \epsilon_{\rm env}\equiv\frac{\Delta\zeta_e}{v_0\tau_f}\ll1,
\end{equation*}
其中 $\|\dot f\|_\infty=0$ 时取 $\tau_f=\infty$。由均值定理，电子包络不同纵向位置采样到的驱动差异相对主幅至多为 $O(\epsilon_{\rm env})$。因此在这一受控极限中可写 $\beta(\zeta,\tau)=\beta_0(\tau)[1+O(\epsilon_{\rm env})]$，输出只获得周期相位
\begin{equation*}
\frac{\fenv^{(f)}(\zeta)}{\fenv^{(i)}(\zeta)}
=\exp[\beta \ee^{\ii\omega\zeta/v_0}-\beta^*\ee^{-\ii\omega\zeta/v_0}].
\end{equation*}
由 Jacobi--Anger 展开，分辨边带的概率为
\begin{equation*}
P_\ell=J_\ell^2(2|\beta|).
\end{equation*}
这一 Bessel 分布同时要求 $\epsilon_{\rm env}\ll1$、平移对称误差受控以及能量梯边界概率足够小；任一条件失效时都应恢复位置依赖耦合或有限反冲动力学。
\end{exbox}
% -----------------------------------------------------------------------------
% -----------------------------------------------------------------------------

\begin{exbox}[200\,keV 电子穿过 800\,nm 高斯近场]
\eqnrefs{eq:pinem_standard_beta_pulsed} 已经把经典 PINEM 的全部场依赖压缩成沿电子轨迹的同步 Fourier 投影。这个积分衡量的是电子轨迹对同步 Fourier 分量的选择。取一个最简单但非平凡的纵向近场包络
\begin{equation*}
 E_{\omega,z}(z)=E_{\omega,0}\exp\!\left(-\frac{z^2}{2L^2}\right),
 \qquad f(t)=1,
\end{equation*}
其中 $E_{\omega,0}$ 是正频电场复振幅的峰值（实场余弦峰值为其模长的两倍），$L$ 是近场沿轨迹的均方根空间宽度。电子初始动能取
\begin{equation*}
 K_0=200\,\mathrm{keV},
\end{equation*}
光学波长、近场峰值和空间宽度分别取
\begin{equation*}
 \lambda=800\,\mathrm{nm},\qquad
 E_{\omega,0}=20\,\mathrm{MV\,m^{-1}},\qquad
 L=100\,\mathrm{nm}.
\end{equation*}

先由自由电子相对论运动学得到
\begin{align*}
 \gamma_0
 &=1+\frac{K_0}{\me c^2}
 =1.39139,\\
 \frac{v_0}{c}
 &=\sqrt{1-\gamma_0^{-2}}
 =0.69531,\\
 v_0&=2.0845\times10^8\,\mathrm{m\,s^{-1}}.
\end{align*}
光学角频率和单光子能量则为
\begin{align*}
 \omega&=\frac{2\pi c}{\lambda}
 =2.3546\times10^{15}\,\mathrm{s^{-1}},\\
 \hbar\omega&=1.5498\,\mathrm{eV}.
\end{align*}
在这个例子中，\eqnrefs{eq:pinem_standard_beta_pulsed} 变成一个标准 Gaussian Fourier 积分：
\begin{align}
 \beta
 &=\frac{\qe E_{\omega,0}}{\hbar\omega}
 \int_{-\infty}^{\infty}\dd z\,
 \exp\!\left(-\frac{z^2}{2L^2}\right)
 \exp\!\left(-\ii\frac{\omega z}{v_0}\right)\notag\\
 &=\frac{\qe E_{\omega,0}\sqrt{2\pi}L}{\hbar\omega}
 \exp\!\left[-\frac12\left(\frac{\omega L}{v_0}\right)^2\right].
 \label{eq:pinem-gaussian-beta-dp65}
\end{align}
这里
\begin{equation*}
 \varphi_{\rm pm}\equiv\frac{\omega L}{v_0}=1.1296
\end{equation*}
是电子飞越近场包络时累计的光学相位尺度，并由\eqnrefs{eq:pinem_standard_beta_pulsed} 中 Fourier 相位的无量纲组合直接确定。数值代入后
\begin{equation}
 \boxed{|\beta|=1.7092,\qquad 2|\beta|=3.4183.}
 \label{eq:pinem-gaussian-beta-number-dp65}
\end{equation}

于是理想低反冲边带概率不需要再解微分方程，直接是
\begin{equation*}
 P_\ell=J_\ell^2(3.4183).
\end{equation*}
几个主要边带为
\begin{center}
\begin{tabular}{c|cccccc}
 $|\ell|$ & 0 & 1 & 2 & 3 & 4 & 5\\\hline
 $P_\ell$ & 0.1351 & 0.02944 & 0.2189 & 0.1413 & 0.03685 & 0.00538
\end{tabular}
\end{center}
这里表中 $|\ell|\ge1$ 的数值是\emph{单侧}边带概率；因为当前纯相位调制满足 $P_{+\ell}=P_{-\ell}$，总概率仍按两侧共同归一化。输出能量峰位为
\begin{equation*}
 K_\ell\simeq K_0+\ell\hbar\omega
 =200\,\mathrm{keV}+1.5498\ell\,\mathrm{eV}.
\end{equation*}
Bessel 权重并不随 $|\ell|$ 单调衰减；例如这里 $|\ell|=2$ 比 $|\ell|=1$ 更强。这不是异常，而是不同量子交换路径在同一整数能量阶上的相干衍射结果。

利用恒等式
\begin{equation*}
 \sum_{\ell=-\infty}^{\infty}\ell^2J_\ell^2(x)=\frac{x^2}{2},
\end{equation*}
还可以不逐项求和就得到电子边带阶数的均方根宽度
\begin{align*}
 \sqrt{\langle\ell^2\rangle}
 &=\sqrt{2}\,|\beta|=2.417,\\
 \sigma_E
 &=\hbar\omega\sqrt{\langle\ell^2\rangle}
 =3.75\,\mathrm{eV}.
\end{align*}
这给出了一个很有用的尺度判断：$|\beta|$ 不是“平均吸收光子数”，而是控制相干能量衍射宽度的脉冲面积。

\begin{center}
\includegraphics[width=0.78\textwidth]{pinem_gaussian_example_dp65.png}\\[-0.2em]
{\small 计算图：200\,keV、800\,nm 高斯近场的离散边带概率。每一根竖线相隔 $\hbar\omega=1.5498\,\mathrm{eV}$，全部权重由同一个 $|\beta|=1.7092$ 决定。}
\end{center}

最后必须检查这组参数是否真的允许使用 Bessel 极限。局域最小耦合小参数为
\begin{equation*}
 \frac{|\qe|E_{\omega,0}/\omega}{p_0}
 =5.15\times10^{-6},
 \qquad p_0=\gamma_0\me v_0,
\end{equation*}
说明单点场强对载波动量只是微扰；但沿整个作用区相干积分后 $|\beta|$ 已经达到 $O(1)$。这两个量衡量的是不同问题，不能用“局域相互作用很弱”推断“最终只会出现一阶边带”。

再检查反冲。同步纵向波数为
\begin{equation*}
 \kappa_z=\frac{\omega}{v_0}=1.1296\times10^7\,\mathrm{m^{-1}},
\end{equation*}
纵向相对论有效质量为 $m_\parallel=\gamma_0^3\me=2.6937\me$，所以
\begin{align*}
 \omega_{\rm r}
 &=\frac{\hbar\kappa_z^2}{2m_\parallel}
 =2.742\times10^9\,\mathrm{s^{-1}},\\
 \hbar\omega_{\rm r}
 &=1.80\,\mu\mathrm{eV}.
\end{align*}
若用 $T_{\rm int}=2L/v_0=0.959\,\mathrm{fs}$ 作为作用时间尺度，并取实际显著占据窗口 $|\ell|\le5$，则\eqnrefs{eq:pinem_recoil_phase_condition} 给
\begin{equation}
 \boxed{
 \epsilon_{\rm rec}^{(T)}
 =(2\times5+1)\omega_{\rm r}T_{\rm int}
 =2.89\times10^{-5}.}
 \label{eq:pinem-gaussian-recoil-number-dp65}
\end{equation}
因此这组参数处在非常干净的低反冲区。反过来，在相同电子能量和光学频率下，如果周期结构把相互作用延长到 $T=1\,\mathrm{ps}$，同一参数已增至约 $3.0\times10^{-2}$；到 $10\,\mathrm{ps}$ 则约为 $0.30$。所以反冲是否重要取决于\emph{累计相位可分辨性}，而不只取决于单次交换的 $\hbar\omega/E_0$。
\end{exbox}



真实自由电子只有正能谱下界，所以理想双边格 $\ell\in\mathbb Z$ 只是局部坐标近似。若 $\Pi_{\rm edge}$ 投影到最低若干可访问格点，定义
\begin{equation}
\boxed{
P_{\rm edge}
=\sup_t\operatorname{Tr}[\rho_e(t)\Pi_{\rm edge}]\ll1.}
\label{eq:bilateral_ladder_edge_condition}
\end{equation}
若累计交换强度为 $\Lambda_I=\int\dd t\,|\upsilon_I(t)|$，则把真实半无限平移算符延拓成双边幺正平移算符的传播误差满足
\begin{equation}
\boxed{
\|[U_b-U_s]|\psi_i\rangle\|
\le\Lambda_I\sqrt{P_{\rm edge}},
\qquad
\epsilon_{\rm edge}^{\rm dyn}=\Lambda_I\sqrt{P_{\rm edge}}\ll1.}
\label{eq:bilateral_ladder_dynamical_error}
\end{equation}

在边界和逐边反冲都不可分辨时，引入双边平移算符 $\hat b|\ell\rangle=|\ell-1\rangle$。由于 $[\hat b,\hat b^\dagger]=0$，任意时间依赖的平移对称生成元在不同时刻也对易，因此传播子只依赖复脉冲面积
\begin{equation}
\boxed{
S_{\rm tr}(\mathcal Z)
=\exp(\mathcal Z \hat b^\dagger-\mathcal Z^*\hat b).}
\label{eq:universal_translation_propagator}
\end{equation}
在统一的相位基底
$|\theta\rangle=(2\pi)^{-1/2}\sum_\ell \ee^{\ii\ell\theta}|\ell\rangle$
中，
$\hat b|\theta\rangle=\ee^{\ii\theta}|\theta\rangle$、
$\hat b^\dagger|\theta\rangle=\ee^{-\ii\theta}|\theta\rangle$，所以
\begin{equation}
\boxed{
S_{\rm tr}(\mathcal Z)|\theta\rangle
=\exp[\mathcal Z \ee^{-\ii\theta}-\mathcal Z^*\ee^{\ii\theta}]|\theta\rangle.}
\label{eq:universal_translation_phase_solution}
\end{equation}
Bessel 边带只是这一相位函数的傅里叶系数。

具体 PINEM 近场确定 $\upsilon(t)$ 或 $\mathcal Z$ 的系统参数，而平移格上的闭式传播子由一般代数结构固定。

\begin{derivation}[从\eqnrefs{eq:bilateral_ladder_dynamical_error}到\eqnrefs{eq:universal_translation_phase_solution}]

正文已经把真实正能谱的下边界误差压缩为\eqnrefs{eq:bilateral_ladder_dynamical_error}。在该误差可忽略、且逐边耦合在实际占据窗口内近似相同时，主线模型只剩双边平移对称生成元
\begin{equation*}
H_I(t)=\ii\hbar[\upsilon(t)\hat b^\dagger-\upsilon^*(t)\hat b].
\end{equation*}
因为 $[\hat b,\hat b^\dagger]=0$，任意两个时刻的生成元满足
\begin{equation*}
[H_I(t),H_I(t')]=0.
\end{equation*}
因此时间排序指数化为普通指数。令
\begin{equation*}
\mathcal Z=\int_{t_i}^{t_f}\dd t\,\upsilon(t),
\end{equation*}
则
\begin{align*}
U_I(t_f,t_i)
&=\exp\!\left[-\frac{\ii}{\hbar}\int_{t_i}^{t_f}\dd t\,H_I(t)\right]\\
&=\exp\!\left[\int_{t_i}^{t_f}\dd t\,
(\upsilon(t)\hat b^\dagger-\upsilon^*(t)\hat b)\right]\\
&=\exp\!\left(\mathcal Z\hat b^\dagger-\mathcal Z^*\hat b\right),
\end{align*}
即正文\eqnrefs{eq:universal_translation_propagator}。

再使用正文给出的 Fourier 相位基
\begin{equation*}
|\theta\rangle=(2\pi)^{-1/2}\sum_{\ell\in\mathbb Z}
 \ee^{\ii\ell\theta}|\ell\rangle,
\end{equation*}
有
\begin{equation*}
\hat b|\theta\rangle=\ee^{\ii\theta}|\theta\rangle,
\qquad
\hat b^\dagger|\theta\rangle=\ee^{-\ii\theta}|\theta\rangle.
\end{equation*}
把这两个本征值代入指数，得到
\begin{equation*}
S_{\rm tr}(\mathcal Z)|\theta\rangle
=\exp[\mathcal Z \ee^{-\ii\theta}-\mathcal Z^*\ee^{\ii\theta}]|\theta\rangle,
\end{equation*}
即正文\eqnrefs{eq:universal_translation_phase_solution}。
\end{derivation}
\subsection{把场强、相位匹配和边带概率连成一次计算}
\label{sec:electron-pinem-field-calibration}

前面已把电子的调制写成沿轨迹的场积分 $\beta$。现在用一个给定近场求可测的边带概率。代入场幅之前，先把实场峰值与复振幅的关系写清楚：本书的 $E_{\omega,z}$ 表示正频分量，因此单频实场为
\begin{equation}
 E_z(z,t)=E_{\omega,z}(z)\ee^{-\ii\omega t}
           +E_{\omega,z}^*(z)\ee^{\ii\omega t}
 =2|E_{\omega,z}(z)|\cos[\omega t-\arg E_{\omega,z}(z)].
 \label{eq:electron-real-field-positive-frequency}
\end{equation}
因此实场的余弦峰值是 $E_{\rm pk}^{\rm(real)}=2|E_{\omega,z}|$，均方根值是 $\sqrt2|E_{\omega,z}|$。若仿真采用 $E_z=\Re[\mathcal E_{\rm ph}\ee^{-\ii\omega t}]$，应输入 $E_{\omega,z}=\mathcal E_{\rm ph}/2$。漏掉这个二分之一会使 $|\beta|$ 加倍，并使弱边带概率放大四倍。

从\eqnrefs{eq:pinem_standard_beta_pulsed} 出发，先令脉冲包络在电子通过期间近似恒定，再考虑轨迹上长度 $L_{\rm int}$ 内的一个纵向空间谐波：
\begin{equation*}
 E_{\omega,z}(z)=E_{\omega,0}\ee^{\ii\kappa_m z},\quad
 -L_{\rm int}/2<z<L_{\rm int}/2,
 \qquad \Delta\kappa=\kappa_m-\frac{\omega}{v_0}.
\end{equation*}
$E_{\omega,0}$ 是复正频振幅，$\kappa_m$ 是场在轨迹方向的波数；矩形包络是实际场采样的简化模型，其边缘需要由结构或源支持，不能解释为无界真空中的自由平面波。有限长度同步积分给出
\begin{equation}
 \beta=\frac{q_{\mathrm e} E_{\omega,0}L_{\rm int}}{\hbar\omega}
 \sinc\!\left(\frac{\Delta\kappa L_{\rm int}}2\right),\qquad
 P_\ell=J_\ell^2(2|\beta|).
 \label{eq:electron-rectangular-nearfield-beta}
\end{equation}
每一小段 $\dd z$ 都贡献一个复振幅；失配使这些复数在复平面转动，相加后部分抵消。因此更长的作用区只有在相位匹配时才持续增强耦合。取复数而不只取场强平方，是保留这种相消所必需的。

对 $v_0<c$ 的电子，同步波数 $\omega/v_0$ 大于真空光波数 $\omega/c$。均匀真空中的传播光无法提供这个纵向 Fourier 分量；结构近场、介质或倒格矢可以承担所需动量。对于真空侧沿表面传播、垂直表面衰减的单个谐波，Maxwell 色散给
\begin{equation*}
 \kappa_\perp^{\rm(ev)}=\sqrt{\kappa_m^2-\omega^2/c^2},\qquad
 \ell_{\rm ev}=1/\kappa_\perp^{\rm(ev)}.
\end{equation*}
实正的 $\kappa_\perp^{\rm(ev)}$ 表示\emph{倏逝}：场幅随离开表面的距离按指数衰减，而不是形成垂直传播的光线。这也解释电子为什么通常必须靠近纳米结构才能有效采样同步分量。这里的 $\ell_{\rm ev}$ 是\emph{场幅}衰减长度，强度衰减长度是其一半。

\begin{exbox}[由实场峰值解析求边带与能宽]
给定实场峰值 $E_{\rm pk}^{\rm(real)}$、作用长度 $L_{\rm int}$ 和失配 $\Delta\kappa$，求弱边带权重，并用完整能谱宽度反求场幅。令
\begin{equation*}
 b=|\beta|=\frac{eE_{\rm pk}^{\rm(real)}L_{\rm int}}{2\hbar\omega}
 \left|\sinc\!\left(\frac{\Delta\kappa L_{\rm int}}2\right)\right|.
\end{equation*}
用 Bessel 幂级数而不逐阶数值求值：
\begin{equation*}
 J_0(2b)=1-b^2+\frac{b^4}{4}+O(b^6),\quad
 J_1(2b)=b-\frac{b^3}{2}+O(b^5),\quad
 J_2(2b)=\frac{b^2}{2}+O(b^4).
\end{equation*}
平方后得到
\begin{equation*}
 P_0=1-2b^2+\frac32b^4+O(b^6),\quad
 P_{\pm1}=b^2-b^4+O(b^6),\quad
 P_{\pm2}=\frac{b^4}{4}+O(b^6).
\end{equation*}
检查归一化时必须把正负阶都算入：$P_0+2P_1+2P_2=1+O(b^6)$。四阶项恰好抵消，是检查展开是否遗漏系数的简便方法。

若不要求 $b\ll1$，仍可使用完整 Bessel 分布的矩恒等式 $\operatorname{Var}\ell=2b^2$。于是交换导致的能量标准差为
\begin{equation*}
 \sigma_{E,\rm ex}=\sqrt2\hbar\omega b
 =\frac{eE_{\rm pk}^{\rm(real)}L_{\rm int}}{\sqrt2}
 \left|\sinc\!\left(\frac{\Delta\kappa L_{\rm int}}2\right)\right|.
\end{equation*}
已知失配时可直接反解 $E_{\rm pk}^{\rm(real)}$。同步处展宽与场强、长度线性相关；接近 $\sinc$ 零点时反演分母很小，场幅误差被放大。这里使用的是相互作用增加的完整能谱方差；源能宽、仪器展宽与有限窗口的处理见\secref{sec:electron-detector-inversion-worked}。
\end{exbox}

\figref{fig:electron-field-calibration} 给出同一窗口在不同失配下的耦合及其边带概率。
\begin{figure}[htbp]
 \centering
 \includegraphics[width=0.92\linewidth]{electron_field_calibration.pdf}
 \caption{从有限长度场积分到电子谱。左图为复振幅相加形成的 $\sinc$ 失配响应；右图取 $\SI{200}{\kilo\electronvolt}$ 电子、$\SI{800}{\nano\metre}$ 光、$\SI{2}{\micro\metre}$ 同步作用区和 $\SI{1}{\mega\volt\per\metre}$ 实场峰值，此时 $|\beta|=0.6452$，分别显示增益和损失阶数。实场峰值与正频振幅的二倍关系已纳入计算。}
 \label{fig:electron-field-calibration}
\end{figure}

\begin{derivation}[自由单光子吸收的运动学检查]
对孤立自由电子吸收一个真实真空光子，若 $p_f=p_i+k$，则 $p_i^2=p_f^2=m_{\mathrm e}^2c^2$ 与 $k^2=0$ 要求 $2p_i\cdot k=0$。但正能电子与正能光子满足
\begin{equation*}
 p_i\cdot k=\frac{E_i\hbar\omega}{c^2}
 (1-\beta_{\rm rel}\cos\vartheta)>0,
\end{equation*}
其中 $\vartheta$ 是入射电子和光子的夹角，$\beta_{\rm rel}=v_i/c<1$。因此无其他参与者的在壳单光子吸收不能发生。近场同步积分没有违背这一点：结构或介质参与了动量守恒。单纯把真空光场人为截成有限空间窗口，不能替代对截断边缘所代表物理源的说明。
\end{derivation}


\section{有限脉冲与空间几何}

实际光脉冲有有限持续时间，近场也只分布在样品附近。不同到达时刻、不同横向位置的电子因此经历不同的场积分。本节把这些差别放回 $\beta$ 中，求有限电子波包的输出振幅，并说明什么时候应当相加振幅，什么时候应当平均概率。


\subsection{有限脉冲与相干平均}

若电子波包的持续时间与光脉冲相当，波包前后各部分受到的调制通常不同。前面得到的 $\beta(\zeta,\tau)$ 已保留这种依赖：$\zeta$ 标记共动波包内的位置，$\tau$ 表示电子与光的相对延迟。在程函近似与通道平移对称仍成立时，每个位置对应的散射相位为
\begin{equation}
\boxed{
\mathcal U_{\rm p}(\zeta,\tau)
=\exp[\beta(\zeta,\tau)\ee^{\ii\kappa_\omega\zeta}
-\beta^*(\zeta,\tau)\ee^{-\ii\kappa_\omega\zeta}],
\qquad \kappa_\omega=\omega/v_0.}
\label{eq:pinem_finite_pulse_local_phase}
\end{equation}
与连续波极限相比，改变的不是局域幺正结构，而是不同 $\zeta$ 对应不同的复耦合，因此实际电子输出必须把这些切片的相干振幅与入射纵向波包一起传播或积分。

多量子路径也可以在不指定脉冲形状时计数。令 $\ell_+=\max(\ell,0)$、$\ell_-=\max(-\ell,0)$。对净交换 $\ell$ 个量子的所有吸收/发射配对，用非负整数 $j$ 标记额外的往返对数，则局域路径系数为
\begin{equation}
\boxed{
C_j^{(\ell)}(\zeta,\tau)
=\frac{\beta(\zeta,\tau)^{j+\ell_+}
[-\beta^*(\zeta,\tau)]^{j+\ell_-}}
{(j+\ell_+)!(j+\ell_-)!}.}
\label{eq:park_Cjn_general}
\end{equation}
因此有限脉冲的普适变化是：同一净边带中的不同 $j$ 路径具有不同的局域权重，最终测得的边带不再只由一个全局复数 $\beta$ 决定。若入射电子纵向振幅记为 $\psi_i(\zeta)$，则对任意具体脉冲都应先在振幅层构造
\[
\psi_f(\zeta,\tau)
=\mathcal U_{\rm p}(\zeta,\tau)\psi_i(\zeta),
\]
再根据探测基计算概率；只有在 $\beta(\zeta,\tau)$ 在电子支持上近似常量时，才恢复单一 Bessel 边带分布。

\begin{exbox}[高斯光脉冲]
取
\[
\beta(\zeta,\tau)=\beta_0
\exp\!\left[-\frac{(\zeta+v_0\tau)^2}{4v_0^2\sigma_p^2}\right].
\]
第 $j$ 条、净交换数为 $\ell$ 的路径含总次数
$r+s=2j+|\ell|$，所以其包络权重额外带有
\[
\exp\!\left[-\frac{(2j+|\ell|)(\zeta+v_0\tau)^2}
{4v_0^2\sigma_p^2}\right].
\]
因此更高阶往返路径采样更窄的有效时窗。若电子自身也有有限纵向包络，实际边带还要对电子振幅作相干积分；短电子极限恢复局域 $J_\ell^2[2|\beta(\zeta_0,\tau)|]$，而长电子或未锁相平均会逐渐抹去载波相位敏感的交叉项。
\end{exbox}

边带归一化直接来自局域相位算符的幺正性，Bessel 恒等式只是这一幺正性的级数表示。有限脉冲下若 $\beta(\zeta,\tau)$ 随切片变化，不同切片的 $A_\ell$ 具有不同复相位；纵向态保持相干时，应先按传播与探测模式叠加振幅，再取模平方。

\begin{derivation}[从\eqnrefs{eq:pinem_finite_pulse_local_phase}到\eqnrefs{eq:park_Cjn_general}]
固定一个电子切片 $(\zeta,\tau)$，把 $\beta(\zeta,\tau)$ 暂记为 $\beta$，并令
\[
 \theta\equiv\kappa_\omega\zeta.
\]
局域散射算符是纯相位
\begin{equation*}
 U(\theta)=\exp\!\left(\beta \ee^{\ii\theta}-\beta^*\ee^{-\ii\theta}\right).
\end{equation*}
由于两个指数中的量只是复数，它们彼此对易，所以可精确分解为
\begin{equation*}
 U(\theta)
 =\ee^{\beta \ee^{\ii\theta}}
  \ee^{-\beta^*\ee^{-\ii\theta}}.
\end{equation*}
分别作 Taylor 展开，得到双重级数
\begin{align*}
 U(\theta)
 &=\left[\sum_{r=0}^{\infty}
 \frac{\beta^r}{r!}\ee^{\ii r\theta}\right]
 \left[\sum_{s=0}^{\infty}
 \frac{(-\beta^*)^s}{s!}\ee^{-\ii s\theta}\right]\\
 &=\sum_{r,s\ge0}
 \frac{\beta^r(-\beta^*)^s}{r!s!}
 \ee^{\ii(r-s)\theta}.
\end{align*}
这里 $r$ 与 $s$ 分别计数从两个指数中取出的正、负频因子数；真正由能量测量分辨的是它们的差
\[
 \ell=r-s,
\]
而不是 $r$、$s$ 各自。于是同一个最终边带 $\ell$ 包含无穷多条具有相同净交换量、但往返次数不同的振幅路径。

若 $\ell\ge0$，所有非负整数解为
\[
 r=j+\ell,\qquad s=j,
 \qquad j=0,1,2,\ldots;
\]
若 $\ell<0$，则
\[
 r=j,\qquad s=j-\ell.
\]
引入
\[
 \ell_+=\max(\ell,0),
 \qquad
 \ell_-=\max(-\ell,0),
\]
两种情况统一为
\[
 r=j+\ell_+,
 \qquad
 s=j+\ell_-.
\]
因此第 $j$ 条路径对净边带 $\ell$ 的复振幅贡献为
\begin{equation*}
 C_j^{(\ell)}
 =\frac{\beta^{j+\ell_+}(-\beta^*)^{j+\ell_-}}
 {(j+\ell_+)!(j+\ell_-)!},
\end{equation*}
恢复 $(\zeta,\tau)$ 依赖即得到正文\eqnrefs{eq:park_Cjn_general}。这里 $C_j^{(\ell)}$ 是振幅而不是概率；所有具有相同 $\ell$ 的路径必须先相干相加：
\begin{equation*}
 A_\ell(\zeta,\tau)
 =\sum_{j=0}^{\infty}C_j^{(\ell)}(\zeta,\tau).
\end{equation*}

这一求和可以显式完成。写
\[
 \beta=|\beta|\ee^{\ii\varphi_\beta}.
\]
对 $\ell\ge0$，有
\begin{align*}
 A_\ell
 &=\beta^\ell
 \sum_{j=0}^{\infty}
 \frac{(-|\beta|^2)^j}{j!(j+\ell)!}\\
 &=\ee^{\ii\ell\varphi_\beta}
 \sum_{j=0}^{\infty}
 \frac{(-1)^j|\beta|^{2j+\ell}}{j!(j+\ell)!}\\
 &=\ee^{\ii\ell\varphi_\beta}J_\ell(2|\beta|),
\end{align*}
其中最后一步使用 Bessel 函数的级数定义
\[
 J_\ell(2x)
 =\sum_{j=0}^{\infty}
 \frac{(-1)^j x^{2j+\ell}}{j!(j+\ell)!}.
\]
对 $\ell<0$ 使用 $J_{-n}(x)=(-1)^nJ_n(x)$ 得到同一统一结果
\begin{equation*}
 \boxed{
 A_\ell=\ee^{\ii\ell\varphi_\beta}J_\ell(2|\beta|).}
\end{equation*}
因此 Bessel 边带并不是额外假设，而是同一净交换通道内所有往返路径振幅的精确重求和。

还可以直接检验概率守恒。由于指数
$\beta \ee^{\ii\theta}-\beta^*\ee^{-\ii\theta}$
纯虚，任意 $\theta$ 都有 $|U(\theta)|=1$。Fourier 系数定义为
\begin{equation*}
 A_\ell=\frac{1}{2\pi}
 \int_0^{2\pi}\dd\theta\,
 U(\theta)\ee^{-\ii\ell\theta}.
\end{equation*}
Parseval 恒等式于是给
\begin{align*}
 \sum_{\ell\in\mathbb Z}|A_\ell|^2
 &=\frac{1}{2\pi}\int_0^{2\pi}\dd\theta\,|U(\theta)|^2\\
 &=1.
\end{align*}
\end{derivation}
\subsection{Green 递推与经典耦合}

也可以直接在能量通道中求同一个低反冲散射问题。电子每吸收或发射一个量子，振幅就转移到相邻通道；自由 Green 函数描述两次交换之间的传播。把这些路径逐次相加，可以得到边带振幅的递推关系，并检验它与轨迹相位法给出的 Bessel 解一致。

对中心速度 $v_0$ 的窄带电子，选中的单色近场只连接相差 $\hbar\omega$ 的通道。把电子能量格写成 $|\ell\rangle$，在平移对称极限中相邻边的复交换率相同，时间积分定义的无量纲耦合就是
\begin{equation}
\boxed{
\beta
=\frac{\qe}{\hbar\omega}
\int_{-\infty}^{\infty}\dd z\,
E_{\omega,z}(z)\ee^{-\ii\omega z/v_0}.}
\label{eq:beta_master}
\end{equation}
\eqnrefs{eq:beta_master} 是脉冲耦合~\eqnrefs{eq:pinem_standard_beta_pulsed} 在包络对电子近似恒定时的极限：$|\beta|$ 是累计相干交换强度，$\arg\beta$ 是电子相对于光学场的相位基准。

在相位坐标 $\theta$ 中，经典散射算符写成
\begin{equation*}
S_{\rm cl}(\beta)|\theta\rangle
=\exp[\beta \ee^{-\ii\theta}-\beta^*\ee^{\ii\theta}]|\theta\rangle.
\end{equation*}
在固定参考面，以中心电子通过时刻为 $t=0$，能量阶数 $\ell$ 的时间因子为 $\ee^{-\ii\ell\omega t}$。因此与 $\beta=|\beta|\ee^{\ii\varphi_\beta}$ 一致的时间相位为 $\varphi_{\rm cl}(t)=2|\beta|\sin(\varphi_\beta-\omega t)$，其瞬时能量偏移是相位的负时间梯度：
\begin{equation}
\boxed{
\delta E_{\rm inst}(t)
=-\hbar\partial_t\varphi_{\rm cl}(t)
=2\hbar\omega|\beta|\cos(\varphi_\beta-\omega t).}
\label{eq:classical_energy_mod_from_beta}
\end{equation}
而能量本征测量把同一个周期相位 Fourier 分解成离散边带。利用 Bessel 生成函数得到算符展开
\begin{equation}
\boxed{
S_{\rm cl}(\beta)
=\sum_{\ell=-\infty}^{\infty}
J_\ell(2|\beta|)\ee^{\ii\ell\varphi_\beta}(\hat b^\dagger)^\ell.}
\label{eq:pinem_fourier_expansion}
\end{equation}
所以“连续能量调制”和“离散 PINEM 边带”是同一散射相位在时间/能量两个共轭表象中的描述。对单能入射和纯相位调制，$J_{-\ell}^2=J_{\ell}^2$ 使增益与损失完全对称，因此平均边带阶数为
\begin{equation}
\boxed{\sum_{\ell=-\infty}^{\infty}\ell J_\ell^2(2|\beta|)=0.}
\label{eq:pinem_zero_mean}
\end{equation}
平均能量保持不变并不意味着没有能量展宽：方差随 $|\beta|$ 增大，而平均值因对称性保持零。

单能入射、双边均匀能量格与纯相位散射共同定义理想 Bessel 有效模型 $P_\ell=J_\ell^2(2|\beta|)$。在这一模型内，边带阶数的统计矩可由特征函数统一求得。定义
\begin{equation}
\boxed{
\chi_\ell(\vartheta)
\equiv\sum_{\ell=-\infty}^{\infty}P_\ell \ee^{\ii\ell\vartheta}
=J_0\!\left(4|\beta|\sin\frac{\vartheta}{2}\right).}
\label{eq:pinem-sideband-characteristic-dp83}
\end{equation}
在 $\vartheta=0$ 处求导得到边带阶数的偶数矩：
\begin{equation}
\boxed{
\langle \ell^2\rangle=2|\beta|^2,
\qquad
\langle \ell^4\rangle=2|\beta|^2+6|\beta|^4.}
\label{eq:pinem-sideband-moments-dp83}
\end{equation}
因此对单能入射，输出能量方差为
\begin{equation}
\boxed{
\operatorname{Var}(E_{\rm out})
=2|\beta|^2(\hbar\omega)^2,}
\label{eq:pinem-energy-variance-dp83}
\end{equation}
而不是一个只凭“边带数大约为 $2|\beta|$”估计的经验宽度。四阶矩还给出对非高斯尾部的直接检验；它与二阶矩一起可用于判断实验谱是否仍由单一纯相位 PINEM 核描述。把 $\ell\hbar\omega$ 解释为真实电子输出能量变化还要求均匀格近似在显著占据窗口内成立；一旦真实通道能量 $\mathcal E_\ell$ 的反冲失谐不可忽略，能量矩必须改用 $\mathcal E_\ell$ 加权，而不能继续把\eqnrefs{eq:pinem-energy-variance-dp83} 当作精确电子色散结果。

\begin{derivation}[从\eqnrefs{eq:pinem_fourier_expansion}到\eqnrefs{eq:pinem-energy-variance-dp83}]
令
\begin{equation*}
 x\equiv2|\beta|.
\end{equation*}
正文纯相位散射因子具有 Fourier 展开
\begin{equation*}
 \ee^{\ii x\sin\phi}
 =\sum_{\ell=-\infty}^{\infty}
 J_\ell(x)\ee^{\ii\ell\phi}.
\end{equation*}
对于理想单能入射，输出第 $\ell$ 条边带的振幅只差一个整体相位，因此
\begin{equation*}
 P_\ell=J_\ell^2(x).
\end{equation*}
为了同时获得归一化和所有矩，不逐项求和 $\sum_\ell \ell^nJ_\ell^2(x)$，而构造概率特征函数
\begin{equation*}
 \chi_\ell(\vartheta)
 =\sum_{\ell=-\infty}^{\infty}
 P_\ell \ee^{\ii\ell\vartheta}.
\end{equation*}

考虑同一个相位因子的角向自相关
\begin{equation*}
 C(\vartheta)
 \equiv\frac1{2\pi}
 \int_0^{2\pi}\dd\phi\,
 \ee^{\ii x\sin(\phi+\vartheta)}
 \ee^{-\ii x\sin\phi}.
\end{equation*}
先用 Fourier 展开计算。第一因子为
\begin{equation*}
 \ee^{\ii x\sin(\phi+\vartheta)}
 =\sum_{\ell}J_\ell(x)
 \ee^{\ii\ell\phi}\ee^{\ii\ell\vartheta},
\end{equation*}
第二因子可写成
\begin{equation*}
 \ee^{-\ii x\sin\phi}
 =\sum_mJ_m(x)\ee^{-\ii m\phi},
\end{equation*}
这里利用 $J_m(x)$ 对实 $x$ 为实数。于是
\begin{align*}
 C(\vartheta)
 &=\frac1{2\pi}
 \sum_{\ell,m}J_\ell(x)J_m(x)\ee^{\ii\ell\vartheta}
 \int_0^{2\pi}\dd\phi\,
 \ee^{\ii(\ell-m)\phi}\\
 &=\sum_{\ell,m}J_\ell(x)J_m(x)\ee^{\ii\ell\vartheta}
 \delta_{\ell m}\\
 &=\sum_\ell J_\ell^2(x)\ee^{\ii\ell\vartheta}
 =\chi_\ell(\vartheta).
\end{align*}
因此自相关本身就是边带阶数的特征函数。

另一方面，不使用 Fourier 级数也可以直接算同一个积分。三角恒等式给
\begin{align*}
 \sin(\phi+\vartheta)-\sin\phi
 &=2\sin\frac{\vartheta}{2}
 \cos\!\left(\phi+\frac{\vartheta}{2}\right).
\end{align*}
令
\begin{equation*}
 u=\phi+\frac{\vartheta}{2},
\end{equation*}
由于积分区间跨越完整 $2\pi$ 周期，平移不改变积分：
\begin{align*}
 C(\vartheta)
 &=\frac1{2\pi}\int_0^{2\pi}\dd u\,
 \exp\!\left[
 \ii\,2x\sin\frac{\vartheta}{2}\cos u
 \right].
\end{align*}
Bessel 函数的积分定义
\begin{equation*}
 J_0(y)
 =\frac1{2\pi}\int_0^{2\pi}\dd u\,\ee^{\ii y\cos u}
\end{equation*}
于是给出
\begin{align*}
 \chi_\ell(\vartheta)
 &=J_0\!\left(2x\sin\frac{\vartheta}{2}\right)\\
 &=J_0\!\left(4|\beta|\sin\frac{\vartheta}{2}\right),
\end{align*}
即正文\eqnrefs{eq:pinem-sideband-characteristic-dp83}。令 $\vartheta=0$：
\begin{equation*}
 \sum_\ell P_\ell
 =\chi_\ell(0)=J_0(0)=1,
\end{equation*}
所以边带分布的归一化也同时得到，而不是另行假定。

特征函数的 $n$ 阶导数满足
\begin{equation*}
 \chi_\ell^{(n)}(0)
 =\sum_\ell(\ii\ell)^nP_\ell
 =\ii^n\langle\ell^n\rangle.
\end{equation*}
由于 $P_{-\ell}=P_\ell$，所有奇数矩为零。为了得到二阶和四阶矩，先把复合函数在 $\vartheta=0$ 附近展开。令
\begin{equation*}
 y(\vartheta)
 =4|\beta|\sin\frac{\vartheta}{2}.
\end{equation*}
正弦展开为
\begin{align*}
 y(\vartheta)
 &=4|\beta|
 \left(\frac\vartheta2-\frac1{6}\frac{\vartheta^3}{8}
 +O(\vartheta^5)\right)\\
 &=2|\beta|\vartheta
 -\frac{|\beta|}{12}\vartheta^3
 +O(\vartheta^5).
\end{align*}
因此
\begin{align*}
 y^2
 &=4|\beta|^2\vartheta^2
 -\frac{|\beta|^2}{3}\vartheta^4
 +O(\vartheta^6),\\
 y^4
 &=16|\beta|^4\vartheta^4
 +O(\vartheta^6).
\end{align*}
再用
\begin{equation*}
 J_0(y)
 =1-\frac{y^2}{4}+\frac{y^4}{64}+O(y^6),
\end{equation*}
逐项代入：
\begin{align*}
 \chi_\ell(\vartheta)
 &=1-\frac14
 \left(4|\beta|^2\vartheta^2
 -\frac{|\beta|^2}{3}\vartheta^4\right)
 +\frac1{64}
 16|\beta|^4\vartheta^4
 +O(\vartheta^6)\\
 &=1-|\beta|^2\vartheta^2
 +\left(
 \frac{|\beta|^2}{12}
 +\frac{|\beta|^4}{4}
 \right)\vartheta^4
 +O(\vartheta^6).
\end{align*}
另一方面，按概率矩展开
\begin{align*}
 \chi_\ell(\vartheta)
 &=1
 +\ii\langle\ell\rangle\vartheta
 -\frac{\langle\ell^2\rangle}{2}\vartheta^2
 -\frac{\ii\langle\ell^3\rangle}{6}\vartheta^3
 +\frac{\langle\ell^4\rangle}{24}\vartheta^4
 +O(\vartheta^5).
\end{align*}
比较 $\vartheta^2$ 系数：
\begin{equation*}
 -\frac{\langle\ell^2\rangle}{2}=-|\beta|^2
 \quad\Rightarrow\quad
 \langle\ell^2\rangle=2|\beta|^2.
\end{equation*}
比较 $\vartheta^4$ 系数：
\begin{align*}
 \frac{\langle\ell^4\rangle}{24}
 &=\frac{|\beta|^2}{12}+\frac{|\beta|^4}{4},\\
 \langle\ell^4\rangle
 &=2|\beta|^2+6|\beta|^4.
\end{align*}
这得到正文\eqnrefs{eq:pinem-sideband-moments-dp83}。

最后把边带阶数映射成理想均匀能量格。在该有效模型内
\begin{equation*}
 E_{\rm out}=E_0+\ell\hbar\omega.
\end{equation*}
所以
\begin{align*}
 \langle E_{\rm out}\rangle
 &=E_0+\hbar\omega\langle\ell\rangle
 =E_0,\\
 \operatorname{Var}(E_{\rm out})
 &=\langle(E_{\rm out}-E_0)^2\rangle\\
 &=(\hbar\omega)^2\langle\ell^2\rangle\\
 &=2|\beta|^2(\hbar\omega)^2,
\end{align*}
即正文\eqnrefs{eq:pinem-energy-variance-dp83}。若真实反冲使通道能量不再满足 $\mathcal E_\ell=E_0+\ell\hbar\omega$，最后这一步必须改为直接对真实 $\mathcal E_\ell$ 做加权求和；前面的 Bessel 概率矩只描述理想均匀格模型内部的统计。
\end{derivation}



\subsection{相干参考与相位可见性}\label{subsec:pinem-coherent-reference-dp31}

复数 $\beta$ 的模和相位都进入电子振幅，而能谱仪只测各振幅的模平方。如果不同的 $\arg\beta$ 给出同一个能谱，单次能谱就无法确定这个相位。为使相位影响概率，可以让未知调制与一个已知参考调制干涉。下面先写出单次测量丢失了什么，再说明参考场怎样提供缺少的信息。

单次理想 PINEM 能谱只含
\begin{equation*}
 P_\ell=J_\ell^2(2|\beta|),
\end{equation*}
因此能量本征测量对 $\arg\beta$ 严格不敏感。这不是反演算法缺陷，而是当前前向模型本身存在相位盲区：把 $\beta$ 乘以任意整体相位都不会改变任何边带概率。要读取复耦合的相位，必须先用一个已知相位参考把相位差转换成可测的强度变化。

设未知相互作用为 $\beta=|\beta|\ee^{\ii\varphi_\beta}$，随后施加与其锁相的已知参考相互作用 $\beta_R=|\beta_R|\ee^{\ii\varphi_R}$。在双边平移近似下，两个散射算符具有同一生成元，因此精确合成为
\begin{equation}
 \boxed{
 S_{\rm cl}(\beta_R)S_{\rm cl}(\beta)
 =S_{\rm cl}(\beta_{\rm tot}),
 \qquad
 \beta_{\rm tot}=\beta+\beta_R.}
 \label{eq:pinem-reference-composition-dp31}
\end{equation}
若两次相互作用之间存在已知自由传播相位 $\varphi_{\rm prop}$，则第二次参考作用的总可控相位写成 $\varphi_R\to\varphi_R+\varphi_{\rm prop}$。定义信号耦合与该总参考相位之差 $\delta\equiv\varphi_\beta-\varphi_R$，则
\begin{equation*}
 B^2\equiv|\beta_{\rm tot}|^2
 =|\beta|^2+|\beta_R|^2
 +2|\beta||\beta_R|\cos\delta,
\end{equation*}
最终边带概率变成
\begin{equation}
 \boxed{P_\ell(\delta)=J_\ell^2(2B).}
 \label{eq:pinem-reference-sidebands-dp31}
\end{equation}
相位已经通过复振幅相干叠加进入 $B$，因此扫描参考相位就能把 $\arg\beta$ 转换成可观测的边带重分布。

对于理想的完整边带读出，全部边带关于 $\delta$ 的经典 Fisher 信息可解析求和：
\begin{equation}
 \boxed{
 F_\delta^{\rm(full)}
 =8\,
 \frac{|\beta|^2|\beta_R|^2\sin^2\delta}
 {|\beta+\beta_R|^2}.}
 \label{eq:pinem-reference-phase-fisher-dp31}
\end{equation}
相位以弧度计，因而 $F_\delta$ 无量纲。令 $r=|\beta|$、$R=|\beta_R|$。对固定的 $r$ 与 $R$，理想完整边带读出的最大局部相位信息为
\begin{equation}
 \boxed{
 F_{\delta,\max}^{\rm(full)}=8\min(r^2,R^2),
 \qquad
 \cos\delta_{\rm opt}
 =-\frac{\min(r,R)}{\max(r,R)}.}
 \label{eq:pinem-reference-phase-fisher-optimum-dp32}
\end{equation}
当 $r=R$ 时最优点逼近完全相消方向；这是一个非正则极限，应从边带概率本身连续取极限，而不能把含 $1/B$ 的中间表达式直接代入 $B=0$。正交参考 $\delta=\pi/2$ 更远离这一奇异参数化点，且在 $R=r$ 时仍有 $F_\delta=4r^2$，因此常是稳定而非绝对最优的实验工作点。若参考显著强于未知信号，正交工作点的信息趋近 $8r^2$。有限边带窗口、探测效率和能量展宽都属于对理想边带分布的经典粗粒化，因此只能降低实际可用的 Fisher 信息。


\subsection{反冲修正}

Bessel 解要求相邻边的完整耦合矩阵元近似相同，并且作用包络不能重新解析出被平均掉的快速支路。设 $\mathcal V_N$ 是实际第 $N\to N+1$ 条边的完整矩阵元，在显著占据窗口 $|N|\le N_{\rm occ}$ 内定义
\begin{equation}
\boxed{
\epsilon_{\rm eq}
=\max_{|N|\le N_{\rm occ}}
\left|\frac{\mathcal V_N}{\mathcal V_0}-1\right|\ll1.}
\label{eq:feist2015_equal_rabi_recoil_condition}
\end{equation}
这个定义只比较实际矩阵元，不预先把其变化归因于速度、场谱或自旋器重叠。

为了把\eqnrefs{eq:feist2015_equal_rabi_recoil_condition} 分解成可计算来源，随后代入一维纵向、$\phi=0$、固定自旋的正能 Dirac 通道。先在长度 $L_z$ 的周期盒中取
\begin{equation*}
 \langle z|p_N,s,+\rangle
 =\frac{1}{\sqrt{L_z}}\ee^{\ii p_Nz/\hbar}u_s(p_N),
 \qquad u_s^\dagger(p_N)u_s(p_N)=1,
\end{equation*}
并令 $\Delta\kappa_N\equiv(p_{N+1}-p_N)/\hbar$。若纵向矢势可分离为 $A_z(z,t)=a(t)A_z(z)$，定义盒归一 Fourier 分量
\begin{equation*}
 \widetilde A_z(\kappa)
 \equiv\frac1{L_z}\int_0^{L_z}\dd z\,A_z(z)\ee^{-\ii\kappa z}
\end{equation*}
以及正能 Dirac 电流矩阵元
\begin{equation*}
 \mathcal J_N
 \equiv c\,u_s^\dagger(p_{N+1})\alpha_z u_s(p_N).
\end{equation*}
把盒归一态代入 $\hat V_A=-\qe c\alpha_zA_z$，逐项积分得到
\begin{align*}
\mathcal V_N
&\equiv {}_V\!\langle p_{N+1},s,+|\hat V_A|p_N,s,+\rangle_V\\
&=-\frac{\qe c}{L_z}
\int_0^{L_z}\dd z\,
 u_s^\dagger(p_{N+1})\alpha_z u_s(p_N)
 A_z(z)\ee^{-\ii(p_{N+1}-p_N)z/\hbar}\\
&=-\qe\,\mathcal J_N\,\widetilde A_z(\Delta\kappa_N).
\end{align*}
因此
\begin{equation}
 \boxed{
 \mathcal V_N
 =-\qe\,\mathcal J_N\,\widetilde A_z(\Delta\kappa_N).}
 \label{eq:pinem-edge-matrix-element-dirac-dp116}
\end{equation}
于是定义
\begin{equation*}
 \epsilon_J
 \equiv\max_N\left|\frac{\mathcal J_N}{\mathcal J_0}-1\right|,
 \qquad
 \epsilon_A
 \equiv\max_N\left|
 \frac{\widetilde A_z(\Delta\kappa_N)}
 {\widetilde A_z(\Delta\kappa_0)}-1\right|,
\end{equation*}
便得到严格误差界
\begin{equation}
 \boxed{
 \epsilon_{\rm eq}
 \le\epsilon_J+\epsilon_A+\epsilon_J\epsilon_A.}
 \label{eq:pinem-equal-coupling-error-dp68}
\end{equation}
因此逐边变化有两个独立来源：电子正能电流矩阵元的变化和场的纵向空间谱变化。只有再取同自旋、小动量步长极限，才有
\begin{equation*}
 \lim_{p_{N+1}\to p_N}\mathcal J_N
 =c\,u_s^\dagger(p_N)\alpha_z u_s(p_N)
 =\frac{\partial E}{\partial p}\bigg|_{p_N}
 =v_N,
\end{equation*}
此时 $\epsilon_J$ 才进一步约化为常见的速度色散判据。

若相互作用具有归一化开关包络 $w(t)$，令 $\widetilde w(0)=1$。共振支路在旋转表象中是慢项，而被舍弃的反旋支路相对它多出 $2\omega$ 的快速相位，因此两条支路的一阶时间积分之比为
\begin{equation}
\boxed{
\frac{\mathcal A_{\rm fast}^{(1)}}{\mathcal A_{\rm slow}^{(1)}}
=\frac{\widetilde w(2\omega)}{\widetilde w(0)}
=\widetilde w(2\omega),
\qquad
\widetilde w(\Omega)=\int\dd t\,w(t)\ee^{\ii\Omega t}.}
\label{eq:pinem_fast_branch_envelope_ratio}
\end{equation}
因此反旋或不共振项的有限脉冲抑制应直接用
\begin{equation}
\boxed{
\epsilon_{2\omega}=|\widetilde w(2\omega)|\ll1.}
\label{eq:feist2015_long_interaction_condition}
\end{equation}
对矩形时间窗它退化成 sinc；对高斯脉冲、啁啾脉冲或脉冲列必须使用各自的 $\widetilde w$。

若等耦合条件失败，最自然的修正不是“改 Bessel 参数”，而是保留自由色散曲率。纵向相对论色散在 $p_0$ 处满足 $E''(p_0)=1/m_\parallel$，所以对 $p_\ell=p_0+\ell\hbar \kappa_z$，
\begin{equation}
\boxed{
\mathcal E_\ell
=E_0+\ell\hbar v_0\kappa_z
+\ell^2\hbar\omega_{\rm r}
+O\!\left[(\ell\hbar \kappa_z)^3E^{(3)}(p_0)\right],
\qquad
\omega_{\rm r}=\frac{\hbar \kappa_z^2}{2m_\parallel}.}
\label{eq:pinem_recoil_lattice_expansion}
\end{equation}
若同步条件取 $v_0\kappa_z=\omega$，第 $\ell\to\ell+1$ 条边相对于理想 $\hbar\omega$ 梯的失谐就是
\begin{equation}
\boxed{
\Delta_\ell
=\frac{\mathcal E_{\ell+1}-\mathcal E_\ell}{\hbar}-\omega
=(2\ell+1)\omega_{\rm r}+O[(\ell\hbar \kappa_z)^2\kappa_zE^{(3)}(p_0)].}
\label{eq:pinem_recoil_edge_detuning}
\end{equation}
因此占据窗 $|\ell|\le N_{\rm occ}$ 内，忽略二阶反冲至少要求
\begin{equation}
\boxed{
\epsilon_{\rm rec}^{(T)}
\equiv(2N_{\rm occ}+1)\omega_{\rm r}T\ll1.}
\label{eq:pinem_recoil_phase_condition}
\end{equation}
若需要把三阶及更高色散的截断误差也纳入同一个有限时间判据，定义
\[
\delta_\ell^{(3)}\equiv
\frac{R_{\ell+1}^{(3)}-R_\ell^{(3)}}{\hbar},
\qquad
\epsilon_{\rm rec}^{(3)}\equiv
T\max_{|\ell|\le N_{\rm occ}}|\delta_\ell^{(3)}|.
\]
逐边总失谐的累计相位满足
\begin{equation}
\boxed{
T\max_{|\ell|\le N_{\rm occ}}|\Delta_\ell|
\le
\epsilon_{\rm rec}^{(T)}+\epsilon_{\rm rec}^{(3)}.}
\label{eq:pinem_recoil_total_phase_bound}
\end{equation}
因此使用均匀能量梯时，不仅要检查二次曲率产生的 $\epsilon_{\rm rec}^{(T)}$，还要确认实际占据窗口中的高阶 Taylor 余项没有重新积累出 $O(1)$ 相位。
若还要把失谐相对于交换速率本身忽略，则需同时满足 $\epsilon_{\rm rec}^{(\Omega)}=(2N_{\rm occ}+1)\omega_{\rm r}/|\Omega|\ll1$。这两个量分别比较反冲与有限作用时间、反冲与耦合动力学；它们比“单光子反冲很小”给出更直接的 Bessel 适用判据。\chapref{chap:periodic-recoil}将从\eqnrefs{eq:pinem_recoil_lattice_expansion} 而不是从均匀梯重新开始。

\begin{derivation}[从\eqnrefs{eq:feist2015_equal_rabi_recoil_condition}到\eqnrefs{eq:pinem-equal-coupling-error-dp68}]
在选定的一维纵向模型中，线性相互作用算符为
\begin{equation*}
 \hat V_A(t)=-\qe c\alpha_z a(t)A_z(\hat z).
\end{equation*}
周期盒正能态满足
\begin{equation*}
 \langle z|p_N,s,+\rangle
 =L_z^{-1/2}\ee^{\ii p_Nz/\hbar}u_s(p_N).
\end{equation*}
因此去掉共同时间包络 $a(t)$ 后，第 $N\to N+1$ 条边的矩阵元为
\begin{align*}
 \mathcal V_N
 &=\langle p_{N+1},s,+|
 (-\qe c\alpha_zA_z)
 |p_N,s,+\rangle\\
 &=-\frac{\qe c}{L_z}
 u_s^\dagger(p_{N+1})\alpha_z u_s(p_N)
 \int_0^{L_z}\dd z\,
 A_z(z)\ee^{-\ii(p_{N+1}-p_N)z/\hbar}\\
 &=-\qe\,\mathcal J_N\,
 \widetilde A_z(\Delta\kappa_N),
\end{align*}
其中
\begin{equation*}
 \mathcal J_N=c\,u_s^\dagger(p_{N+1})\alpha_z u_s(p_N),
 \qquad
 \Delta\kappa_N=\frac{p_{N+1}-p_N}{\hbar}.
\end{equation*}
这得到正文\eqnrefs{eq:pinem-edge-matrix-element-dirac-dp116}。

假定 $\mathcal V_0\neq0$、$\mathcal J_0\neq0$ 且 $\widetilde A_z(\Delta\kappa_0)\neq0$，则
\begin{align*}
 \frac{\mathcal V_N}{\mathcal V_0}
 &=\frac{\mathcal J_N}{\mathcal J_0}
 \frac{\widetilde A_z(\Delta\kappa_N)}
 {\widetilde A_z(\Delta\kappa_0)}\\
 &=\left[1+\delta_J^{(N)}\right]
 \left[1+\delta_A^{(N)}\right]\\
 &=1+\delta_J^{(N)}+\delta_A^{(N)}
 +\delta_J^{(N)}\delta_A^{(N)},
\end{align*}
其中
\begin{equation*}
 \delta_J^{(N)}\equiv\frac{\mathcal J_N}{\mathcal J_0}-1,
 \qquad
 \delta_A^{(N)}\equiv
 \frac{\widetilde A_z(\Delta\kappa_N)}
 {\widetilde A_z(\Delta\kappa_0)}-1.
\end{equation*}
于是对 $|N|\le N_{\rm occ}$ 取最大值并使用三角不等式：
\begin{align*}
 \epsilon_{\rm eq}
 &=\max_N\left|
 \delta_J^{(N)}+\delta_A^{(N)}
 +\delta_J^{(N)}\delta_A^{(N)}\right|\\
 &\le\epsilon_J+\epsilon_A+\epsilon_J\epsilon_A,
\end{align*}
这就是\eqnrefs{eq:pinem-equal-coupling-error-dp68}。

最后说明速度因子的来源。对同一归一化正能本征态，Hellmann--Feynman 恒等式给
\begin{align*}
 \frac{\partial E(p)}{\partial p}
 &=u_s^\dagger(p)
 \frac{\partial}{\partial p}
 \left(c\alpha_z p+\beta_D\me c^2\right)
 u_s(p)\\
 &=c\,u_s^\dagger(p)\alpha_z u_s(p).
\end{align*}
在 $p_{N+1}\to p_N$ 的小步长极限，$\mathcal J_N\to v_N$。
\end{derivation}
速度色散是正能 Dirac 电流矩阵元的低反冲极限，而\eqnrefs{eq:pinem-edge-matrix-element-dirac-dp116} 保留有限步长下的自旋器重叠。

\begin{derivation}[从\eqnrefs{eq:pinem_fast_branch_envelope_ratio}到\eqnrefs{eq:feist2015_long_interaction_condition}]
把归一化开关包络取在有限区间 $[t_i,t_f]$，满足
\[
 \widetilde w(0)=\int_{t_i}^{t_f}\dd t\,w(t)=1,
 \qquad
 \widetilde w(\Omega)=\int_{t_i}^{t_f}\dd t\,w(t)\ee^{\ii\Omega t}.
\]
正文的一阶微扰已经把慢支路与快支路之比化为
$|\widetilde w(2\omega)|$。为了看清“作用时间长”为什么能够压低它，假设 $w$ 在作用区间绝对连续。对 $\Omega\ne0$ 分部积分：
\begin{align*}
 \widetilde w(\Omega)
 &=\left.\frac{w(t)\ee^{\ii\Omega t}}{\ii\Omega}\right|_{t_i}^{t_f}
 -\frac1{\ii\Omega}\int_{t_i}^{t_f}\dd t\,
 \dot w(t)\ee^{\ii\Omega t}.
\end{align*}
取模并用三角不等式得到严格上界
\begin{align*}
 |\widetilde w(\Omega)|
 \le\frac{|w(t_f)|+|w(t_i)|
 +\int_{t_i}^{t_f}\dd t\,|\dot w(t)|}{|\Omega|}.
\end{align*}
因此在 $\Omega=2\omega$ 处，若包络端点平滑降到零，便有
\begin{equation*}
 \left|\frac{\mathcal A_{\rm fast}^{(1)}}{\mathcal A_{\rm slow}^{(1)}}\right|
 =|\widetilde w(2\omega)|
 \le\frac{1}{2\omega}
 \int_{t_i}^{t_f}\dd t\,|\dot w(t)|.
\end{equation*}
这个式子比单纯说“$T$ 很长”更精确：快支路由包络在光学频率尺度上的变化率控制；同样持续时间的两个脉冲，只要边缘平滑程度不同，反旋泄漏就可以不同。

矩形窗口不满足端点为零的平滑条件，需要直接积分。对
$w(t)=T^{-1}\mathbf 1_{[-T/2,T/2]}(t)$，
\begin{align*}
 \widetilde w(2\omega)
 &=\frac1T\int_{-T/2}^{T/2}\dd t\,\ee^{2\ii\omega t}\\
 &=\frac{\sin(\omega T)}{\omega T}
 =\sinc(\omega T).
\end{align*}
所以正文\eqnrefs{eq:feist2015_long_interaction_condition} 的一般小参数是
$\epsilon_{2\omega}=|\widetilde w(2\omega)|$；只有指定矩形开关之后，它才退化为由 $\omega T$ 单独控制的 sinc 条件。
\end{derivation}

\begin{derivation}[从\eqnrefs{eq:pinem_recoil_lattice_expansion}到\eqnrefs{eq:pinem_recoil_total_phase_bound}]
正文\eqnrefs{eq:pinem_recoil_lattice_expansion} 给出
\[
\mathcal E_\ell
=E_0+\ell\hbar\omega+\ell^2\hbar\omega_{\rm r}+R_\ell^{(3)}.
\]
相邻两条真实电子通道之间的角频率为
\[
\omega_{\ell+1,\ell}
=\frac{\mathcal E_{\ell+1}-\mathcal E_\ell}{\hbar},
\]
故相对于同步频率 $\omega$ 的逐边失谐为
\[
\Delta_\ell
=\frac{\mathcal E_{\ell+1}-\mathcal E_\ell}{\hbar}-\omega.
\]
对离散能量差逐项计算：
\begin{align*}
\mathcal E_{\ell+1}-\mathcal E_\ell
={}&[(\ell+1)-\ell]\hbar\omega
+[(\ell+1)^2-\ell^2]\hbar\omega_{\rm r}
+R_{\ell+1}^{(3)}-R_\ell^{(3)}\\
={}&\hbar\omega+(2\ell+1)\hbar\omega_{\rm r}
+R_{\ell+1}^{(3)}-R_\ell^{(3)}.
\end{align*}
因此
\[
\Delta_\ell
=(2\ell+1)\omega_{\rm r}+\delta_\ell^{(3)},
\qquad
\delta_\ell^{(3)}
=\frac{R_{\ell+1}^{(3)}-R_\ell^{(3)}}{\hbar}.
\]
在显著占据窗口 $|\ell|\le N_{\rm occ}$ 内，二次色散贡献的最大累计相位为
\begin{align*}
T\max_{|\ell|\le N_{\rm occ}}
|(2\ell+1)\omega_{\rm r}|
&=T\omega_{\rm r}
\max_{|\ell|\le N_{\rm occ}}|2\ell+1|\\
&=(2N_{\rm occ}+1)\omega_{\rm r}T\\
&=\epsilon_{\rm rec}^{(T)},
\end{align*}
得到正文\eqnrefs{eq:pinem_recoil_phase_condition}。对完整逐边失谐使用三角不等式：
\begin{align*}
T\max_{|\ell|\le N_{\rm occ}}|\Delta_\ell|
&\le
T\max_{|\ell|\le N_{\rm occ}}|(2\ell+1)\omega_{\rm r}|
+T\max_{|\ell|\le N_{\rm occ}}|\delta_\ell^{(3)}|\\
&=\epsilon_{\rm rec}^{(T)}+\epsilon_{\rm rec}^{(3)},
\end{align*}
即正文\eqnrefs{eq:pinem_recoil_total_phase_bound}。
\end{derivation}
\section{量子交换：多模一般动力学方程与单模闭式}

经典场模型把光的作用写成给定复振幅，适合平均场主导的相干光。若要比较不同光子数分布，或求电子与光之间的纠缠，就需要保留光场的量子态。我们从上一章的多模联合 Hamiltonian 出发，先写出一般交换方程，再在单模及低反冲条件下求解析传播子。

多模连续谱层级~\eqnrefs{eq:universal_multiboson_amplitude_hierarchy} 给出量子交换的母动力学。先把有限组已谱隔离的环境模式写成多模联合格；只有当其它模式的有限时间跃迁幅度均低于目标精度时，才保留一个孤立模式。随后独立检查电子逐边失谐与耦合的平移对称性，满足条件后才能引入双边电子平移算符 $\hat b$。光场的相干态大占据极限发生在另一 Hilbert 空间方向，不替代前两项谱选择条件。


量子交换的母模型不应从单模开始。先假设环境谱已经由上一章的一般 Green/正常模分析选出 $M$ 个需要显式保留的正交玻色模式，但电子仍保留完整正能连续谱。第 $\nu$ 个模式的中心频率为 $\omega_\nu$，湮灭算符为 $\hat a_\nu$。在相互作用绘景中，最一般的线性单量子交换仍是连续核
\begin{equation}
\boxed{
\hat H_{{\rm Q},I}^{(M)}(t)
=
\hbar
\sum_{\nu=1}^{M}
\sum_{ss'}
\int\dd^3p\,\dd^3p'\,
\left[
\mathcal G_{\nu;s's}(\mathbf p',\mathbf p;t)
|\mathbf p',s'\rangle\langle\mathbf p,s|
\,\hat a_\nu
+\mathrm{H.c.}
\right].}
\label{eq:multimode_qpinem_continuous_kernel_dp101}
\end{equation}
$\mathcal G_{\nu;s's}$ 由 QED 电流--模式矩阵元或相应材料量子化给出；这一层还没有电子“格点”，也没有为每个光学模式预先指定唯一动量步长。

很多纳米光子学几何中，第 $\nu$ 个模式只在一个窄的纵向 Fourier 区间内与准直电子强耦合。若实际电子工作窗内的主导转移为 $\hbar\kappa_\nu$，并且自旋翻转及其它动量转移的累计传播修正很小，可以在通道子空间上写
\[
\hat{\mathcal G}_{\nu,I}(t)
=
\ii\upsilon_\nu(t)\hat b_\nu^\dagger
+\hat R_{\nu,\kappa}(t),
\qquad
\hat b_\nu^\dagger|p\rangle
=
|p+\hbar\kappa_\nu\rangle .
\]
这里 $\hat R_{\nu,\kappa}$ 收集非主导 Fourier 分量、逐边矩阵元变化和被舍弃自旋通道。若 $\Pi_{\rm in}$ 投影到实际入射工作窗，以只保留主导平移项的传播子记为 $\hat U_{\rm tr}$，则一个直接的有限时间充分条件是
\begin{equation}
\boxed{
\epsilon_{\kappa}^{(M)}
\equiv
\frac1\hbar
\int_{t_i}^{t_f}\dd t\,
\left\|
\hbar\sum_\nu
\left[
\hat R_{\nu,\kappa}(t)\hat a_\nu
+\hat R_{\nu,\kappa}^\dagger(t)\hat a_\nu^\dagger
\right]
\hat U_{\rm tr}(t,t_i)\Pi_{\rm in}
\right\|
\ll1 .}
\label{eq:multimode_momentum_step_projection_error_dp101}
\end{equation}
只有通过这一类通道投影后，才得到常用的多模平移 Hamiltonian
\begin{equation}
\boxed{
 \hat H_{\rm Q}^{(M)}(t)
 =\ii\hbar\sum_{\nu=1}^{M}
 \left[
 \upsilon_\nu(t)\hat b_\nu^\dagger\hat a_\nu
 -\upsilon_\nu^*(t)\hat b_\nu\hat a_\nu^\dagger
 \right].}
\label{eq:multimode_qpinem_hamiltonian_dp47}
\end{equation}
因此 $\upsilon_\nu$ 由连续耦合核在选定电子通道上的投影唯一确定。

重复交换以后，整数向量
\[
\boldsymbol\ell=(\ell_1,\ldots,\ell_M)\in\mathbb Z^M
\]
映射到真实电子动量
\begin{equation*}
\mathbf p_{\boldsymbol\ell}
=
\mathbf p_0
+\hbar\sum_\nu\ell_\nu\boldsymbol\kappa_\nu .
\end{equation*}
把 $|\boldsymbol\ell\rangle_e$ 作为正交的多维格点标签，还要求整数映射在实际占据窗口 $\mathcal W_\ell$ 内保持一一对应。若存在不同整数向量满足
\[
\sum_\nu
(\ell_\nu-\ell_\nu')\boldsymbol\kappa_\nu=0,
\]
它们代表同一个电子动量末态，必须先在振幅层面相干相加，而不能作为不同格点分别计数。对可唯一标记的通道，真实电子能量为
\begin{equation}
 \mathcal E_{\boldsymbol\ell}
 =E\!\left(
 \mathbf p_0+\hbar\sum_\nu
 \ell_\nu\boldsymbol\kappa_\nu
 \right),
\end{equation}
并相对于理想多频能量梯定义
\begin{equation}
\boxed{
 \delta_{\boldsymbol\ell}
 =
 \frac{
 \mathcal E_{\boldsymbol\ell}-E_0
 -\sum_\nu\ell_\nu\hbar\omega_\nu
 }{\hbar}.}
\label{eq:multimode_recoil_detuning_dp47}
\end{equation}

对初始 Fock 扇区
$\mathbf n=(n_1,\ldots,n_M)$，取守恒基底
\[
 |\boldsymbol\ell;\mathbf n\rangle
 =
 |\boldsymbol\ell\rangle_e
 \otimes
 \bigotimes_\nu
 |n_\nu-\ell_\nu\rangle_\nu ,
\]
其中每个光子占据数都必须非负。把联合态展开为
\[
|\Psi_{\mathbf n}(t)\rangle
=
\sum_{\boldsymbol\ell}
c_{\boldsymbol\ell}^{(\mathbf n)}(t)
|\boldsymbol\ell;\mathbf n\rangle ,
\]
逐项使用玻色梯矩阵元后得到
\begin{equation}
\boxed{
\begin{aligned}
 \ii\dot c_{\boldsymbol\ell}^{(\mathbf n)}
 =&\,\delta_{\boldsymbol\ell}
 c_{\boldsymbol\ell}^{(\mathbf n)}\\
 &+\ii\sum_\nu \upsilon_\nu(t)
 \sqrt{n_\nu-\ell_\nu+1}\,
 c_{\boldsymbol\ell-\mathbf e_\nu}^{(\mathbf n)}\\
 &-\ii\sum_\nu \upsilon_\nu^*(t)
 \sqrt{n_\nu-\ell_\nu}\,
 c_{\boldsymbol\ell+\mathbf e_\nu}^{(\mathbf n)} .
\end{aligned}}
\label{eq:multimode_qpinem_lattice_dp47}
\end{equation}
这才是有限组显式模式下的多模 QPINEM 数值一般动力学方程：真实电子色散、多个 Fock 下界、有限脉冲和模式间共享的电子反冲全部保留。单模模型只需令 $M=1$，但若进一步要得到双边位移闭式，还必须额外满足逐边失谐与逐边耦合近似均匀。

对多模相干态
$|\{\alpha_\nu\}\rangle$，经典场极限应逐模取
\[
|\alpha_\nu|\to\infty,
\qquad
\gqmode{\nu}\to0,
\qquad
\beta_\nu
\equiv
\alpha_\nu\gqmode{\nu}
\quad\text{保持有限}.
\]
若在相应工作窗中各电子平移算符可用完整连续谱上的平移表示，它们彼此对易，电子得到
\begin{equation}
\boxed{
 S_{\rm cl}^{(M)}
 =
 \exp\!\left[
 \sum_\nu
 \left(
 \beta_\nu\hat b_\nu^\dagger
 -\beta_\nu^*\hat b_\nu
 \right)
 \right].}
\label{eq:multicolor_classical_scattering_dp47}
\end{equation}
若各动量步长进一步满足
$\kappa_\nu=m_\nu\kappa_0$，不同颜色的整数路径可能到达同一个最终电子格点。把每个颜色的 Jacobi--Anger 展开相乘，并对所有满足
$\ell=\sum_\nu m_\nu r_\nu$
的路径作相干求和，得到
\begin{equation}
\boxed{
 A_\ell
 =\sum_{\{r_\nu\}}
 \delta_{\ell,\sum_\nu m_\nu r_\nu}
 \prod_\nu
 J_{r_\nu}(2|\beta_\nu|)
 \ee^{\ii r_\nu\varphi_\nu},
 \qquad
 \varphi_\nu=\arg\beta_\nu .}
\label{eq:multicolor_pinem_convolution_dp47}
\end{equation}
因此多色 PINEM 的相位敏感性来自“不同交换路径汇聚到同一电子末态”后的振幅干涉；若路径对应不同可分辨电子动量，则应保留多指标末态，而不能先把概率卷积到一条能量格上。


\subsection{单模交换哈密顿量}

设前面给出的单模选择和旋波条件成立，并且在实际占据的电子通道中，各条边的耦合与失谐可以用共同值代替。电子通道因此可近似编号为 $\ell\in\mathbb Z$。定义电子平移算符与光子湮灭算符的作用为
\[
\hat b^\dagger|\ell\rangle_e=|\ell+1\rangle_e,\qquad
\hat b|\ell\rangle_e=|\ell-1\rangle_e,\qquad
\hat a|n\rangle_{\rm ph}=\sqrt n\,|n-1\rangle_{\rm ph}.
\]
前两者改变电子通道编号，后者改变光子数。它们作用在不同的态空间上，因而彼此对易。电子吸收一个光子时，光子数减少一而电子阶数增加一，对应乘积 $\hat b^\dagger\hat a$；逆过程对应 $\hat b\hat a^\dagger$。

沿用真实反冲格的耦合率 $\mathcal G_0(t)$ 和共同失谐 $\Delta_0$，把相互作用绘景中的速率记为
\[
\upsilon(t)\equiv-\ii\mathcal G_0(t)\ee^{\ii\Delta_0t},
\qquad [\upsilon]={\rm s^{-1}}.
\]
因子 $-\ii$ 是相位约定，使传播子中的累积耦合可直接写为
$\gq=\int_{t_i}^{t_f}\dd t\,\upsilon(t)$。$\gq$ 无量纲。若光场处在振幅为 $\alpha_{\rm coh}(t)$ 的相干态，经典耦合为 $\beta=\int\dd t\,\alpha_{\rm coh}(t)\upsilon(t)$；相干幅度为常数时即 $\beta=\alpha_{\rm coh}\gq$。若以后把电子限制为一条二态交换边，常耦合记号 $g_c=\ii\upsilon$ 仍具有频率单位。由上述吸收与发射过程，联合 Hamiltonian 写成
\begin{equation}
\boxed{
\hat H_Q(t)=\ii\hbar[\upsilon(t)\hat b^\dagger\hat a-\upsilon^*(t)\hat b \hat a^\dagger].}
\label{eq:qpinem-single-mode-hamiltonian-dp68}
\end{equation}
其中 $[\upsilon]={\rm s^{-1}}$。电子双边格只是局部近似，而光子数仍有 $n=0$ 下界，因此量子边界误差比经典情况多出玻色增强因子。

若 $\Pi_{\min}$ 投影到电子最低格点附近，沿单边近似演化定义的量子边界传播误差为
\begin{equation}
\boxed{
\epsilon_{\rm edge,Q}^{\rm dyn}
=\int_{t_i}^{t_f}\dd t\,|\upsilon(t)|
\sqrt{\langle\Pi_{\min}(\hat n+1)\rangle_t}\ll1.}
\label{eq:quantum_bilateral_ladder_dynamical_error}
\end{equation}
若实际光子数扇区截断在 $n\le N$，则
$\epsilon_{\rm edge,Q}^{\rm dyn}\le\Lambda_I\sqrt{(N+1)P_{\rm edge}}$。因此经典与量子双边延拓不能共用同一个边界判据。

单模 QPINEM Hamiltonian 同时依赖环境单模隔离、电子单动量步长投影、格点平移对称和 RWA。任一条件失效时，应退回相应上一级完整 Hamiltonian。

\begin{derivation}[从\eqnrefs{eq:universal_free_electron_boson_hamiltonian}到\eqnrefs{eq:qpinem-single-mode-hamiltonian-dp68}]

一般线性电子--玻色相互作用在相互作用绘景中可写成
\[
\hat V_I(t)
=
\hbar\sum_\lambda
\left[
\hat{\mathcal G}_{\lambda,I}(t)\hat a_\lambda
+
\hat{\mathcal G}_{\lambda,I}^\dagger(t)\hat a_\lambda^\dagger
\right].
\]
这里将光模自由演化的 $\ee^{-\ii\omega_\lambda t}$ 也吸收到 $\hat{\mathcal G}_{\lambda,I}$ 中；$\hat a_\lambda$ 为固定的相互作用绘景模算符，不能再重复附加该相位。$\hat{\mathcal G}_{\lambda,I}$ 仍作用在完整正能电子连续谱上。先把电子自由本征态记为 $|p,s\rangle$，则
\[
\hat{\mathcal G}_{\lambda,I}(t)
=
\sum_{ss'}\int\dd p\,\dd p'\,
\mathcal G_{\lambda;s's}(p',p)\,
\ee^{+\ii\{[E(p')-E(p)]/\hbar-\omega_\lambda\}t}
|p',s'\rangle\langle p,s|.
\]
因此“单模 QPINEM”需要依次做三个投影，而不是一步把完整 Hamiltonian 改写成升降算符。

第一步是环境谱投影。设目标模式为 $\lambda_0$，其频率为 $\omega_0$；其余模式与目标模式的最小有效失谐记为
\[
\Delta_{\rm spec}
=
\inf_{\lambda\neq\lambda_0}
|\omega_\lambda-\omega_0|.
\]
若作用时间 $T$ 和实际耦合带宽满足
\[
\Delta_{\rm spec}T\gg1,
\qquad
\sup_{\lambda\neq\lambda_0}
\int_{t_i}^{t_f}\dd t\,
\|\hat{\mathcal G}_{\lambda,I}(t)\|
\ll1,
\]
则有限时间 Fourier 核抑制非目标模式，环境和式可约化为 $\lambda_0$。

第二步是电子通道投影。连续谱广义本征态满足 Dirac--$\delta$ 归一化，而离散通道模型需要一组真正 Kronecker 正交的归一态。先在长度 $L_z$ 的周期盒中正则化纵向连续谱，使允许动量间隔为 $2\pi\hbar/L_z$；再选择与主导动量转移 $\hbar\kappa_z$ 相容的一串正交盒归一态。记
\[
|\ell\rangle_V
\equiv
|p_\ell,s_0\rangle_V,
\qquad
p_\ell=p_0+\ell\hbar\kappa_z,
\qquad
{}_V\langle\ell|m\rangle_V=\delta_{\ell m}.
\]
若实际电子是有限宽波包，也可以用彼此近似正交、中心动量为 $p_\ell$ 的归一化通道包代替盒态；最终有效格点只要求相邻通道的重叠误差小于所保留动力学精度。记号上省略盒归一化下标 $V$，而离散化极限保持不变。投影后的电子耦合算符为
\[
\hat{\mathcal G}_{0,I}(t)
=
\sum_\ell
\mathcal G_\ell(t)
\ee^{+\ii\Delta_\ell t}
|\ell+1\rangle\langle\ell|
+
\hat R_{\rm off}(t),
\]
其中
\[
\Delta_\ell
=
\frac{E(p_{\ell+1})-E(p_\ell)}{\hbar}
-\omega_0
\]
是逐边真实失谐，$\hat R_{\rm off}$ 收集自旋翻转、非主导动量转移和其它离散通道。

第三步才是平移对称近似。若显著占据窗口 $|\ell|\le L_{\rm occ}$ 内
\[
\epsilon_{\mathcal G}
\equiv
\max_\ell
\left|
\frac{\mathcal G_\ell-\mathcal G_0}{\mathcal G_0}
\right|
\ll1,
\qquad
\epsilon_\Delta
\equiv
T\max_\ell|\Delta_\ell-\Delta_0|
\ll1,
\]
便可把共同边耦合写成
\[
\mathcal G_\ell(t)\ee^{+\ii\Delta_\ell t}
=
\mathcal G(t)
+
O\!\left(|\mathcal G_0|(\epsilon_{\mathcal G}+\epsilon_\Delta)\right),
\]
并定义双边电子平移算符
\[
\hat b^\dagger
=
\sum_{\ell\in\mathbb Z}
|\ell+1\rangle\langle\ell|,
\qquad
\hat b
=
\sum_{\ell\in\mathbb Z}
|\ell\rangle\langle\ell+1|.
\]
于是
\[
\hat{\mathcal G}_{0,I}(t)
=
\mathcal G(t)\hat b^\dagger
+\delta\hat{\mathcal G}(t).
\]

把这一结果代回单模完整 Hamiltonian，
\[
\hat V_I(t)
=
\hbar
\left[
\mathcal G(t)\hat b^\dagger\hat a
+
\mathcal G^*(t)\hat b\hat a^\dagger
\right]
+\hat R_{\rm RWA}(t)
+\hat R_{\rm tr}(t),
\]
其中 $\hat R_{\rm tr}$ 收集逐边不均匀和通道投影误差。对一般实线性场，完整相互作用还包含电子上移同时产生光子、电子下移同时湮灭光子的反旋支路；它们在相互作用绘景中带有约 $2\omega_0$ 的快速相位。这些快速支路能否舍弃，不能只由“$\omega_0$ 很大”判断。前章\eqnrefs{eq:rwa_projected_dynamical_error} 已把旋波余项的真实有限时间误差写成 Duhamel 范数；在脉冲包络缓变且反旋相位频率约为 $2\omega_0$ 的常见情形，分部积分才进一步给出\eqnrefs{eq:rwa-finite-time-error-scales-dp68} 那类充分尺度条件。只有该投影传播误差满足
\[
\epsilon_{\rm RWA}^{\rm dyn}\ll1
\]
时，反旋支路才可受控消去；$|\mathcal G|/\omega_0$ 与 $(\omega_0T)^{-1}$ 只是缓变脉冲下对这个更一般误差界的估计。

保留真实矩阵元的记号 $\mathcal G(t)$，并另定义 $\upsilon(t)=-\ii\mathcal G(t)$，即 $\mathcal G(t)=\ii\upsilon(t)$。这样 $\gq=\int\upsilon(t)\dd t$ 与经典 $\beta$ 的相位对应是显式的。代入上式，
\[
\boxed{
\hat H_Q(t)
=
\ii\hbar
\left[
\upsilon(t)\hat b^\dagger\hat a
-
\upsilon^*(t)\hat b\hat a^\dagger
\right].
}
\]
这就是正文\eqnrefs{eq:qpinem-single-mode-hamiltonian-dp68}。
\end{derivation}

\subsection{Magnus 闭合与守恒量}

含时 Hamiltonian 在不同时刻一般不对易，因此传播子需要保留时间排序。为了求解析表达式，可把传播子写成一个指数 $\exp(\Omega_1+\Omega_2+\cdots)$：第一项 $\Omega_1$ 是相互作用的时间积分，第二项 $\Omega_2$ 修正不同时刻的作用次序，更高项含嵌套对易子。这称为 Magnus 展开。

对\eqnrefs{eq:qpinem-single-mode-hamiltonian-dp68}，电子双边平移满足 $\hat b^\dagger\hat b=\hat b\hat b^\dagger=I$，光子算符满足 $[\hat a,\hat a^\dagger]=I$。两时刻 Hamiltonian 的对易子于是正比于恒等算符，再与 Hamiltonian 对易就为零。因此本模型只需保留前两项。把对应时间积分分别记为累积耦合与整体相位，
\begin{equation}
\boxed{
\gq=\int_{t_i}^{t_f}\dd t\,\upsilon(t),
\qquad
\Phi_M
=\frac{1}{2\ii}
\int_{t_i}^{t_f}\dd t_1
\int_{t_i}^{t_1}\dd t_2\,
\bigl[\upsilon^*(t_1)\upsilon(t_2)-\upsilon(t_1)\upsilon^*(t_2)\bigr].}
\label{eq:qpinem-magnus-pulse-parameters-dp68}
\end{equation}
在\eqnrefs{eq:qpinem-single-mode-hamiltonian-dp68} 的模型定义域内，完整时间有序传播子为
\begin{equation}
\boxed{
\hat U_Q(t_f,t_i)
=\ee^{\ii\Phi_M}
\exp[\gq\,\hat b^\dagger\hat a-\gq^*\hat b\hat a^\dagger].}
\label{eq:qpinem-magnus-exact-dp68}
\end{equation}
$\ee^{\ii\Phi_M}$ 是与电子和光子量子数都无关的整体相位，因此所有单次概率和约化密度算符都不受它影响。若 $\upsilon(t)$ 在整个脉冲中具有固定复相位，则被积的反对称组合恒为零，$\Phi_M=0$；此时正文常用的交换算符可简写为
\begin{equation}
\boxed{
S_Q(\gq)=\exp[\gq\,\hat b^\dagger\hat a-\gq^*\hat b\hat a^\dagger].}
\label{eq:qpinem_S}
\end{equation}
电子平移算符与光子升降算符分属不同的 Hilbert 空间；$\hat b$ 是幺正平移算符，而 $\hat a$ 是玻色湮灭算符。两者分别实现电子平移代数与玻色 Fock 代数。


在同一单模双边交换模型内，生成元具有守恒量
\begin{equation}
\boxed{
\mathcal N_{\rm ex}=\Lambda_e+N_{\rm ph},
\qquad
\Lambda_e=\sum_{\ell\in\mathbb Z}\ell|\ell\rangle\langle\ell|,
\qquad
N_{\rm ph}=\hat a^\dagger\hat a.}
\label{eq:qpinem-exchange-number-dp68}
\end{equation}
因为每次电子上移一个能量步长都伴随一个光子被湮灭，反过程则伴随光子产生。于是
\begin{equation}
\boxed{[\mathcal N_{\rm ex},S_Q]=0.}
\label{eq:qpinem_conservation}
\end{equation}
若初态为 $|0,n\rangle$，所有末态必须满足 $\ell+m=n$。这条选择规则来自交换数守恒，与对 $\gq$ 的微扰展开无关；在单模交换 Hamiltonian 的适用域内，对任意有限 $\gq$ 都成立。

若 $\upsilon(t)=\ee^{\ii\varphi_0}u(t)$ 且 $u(t)$ 为实包络，则 $\upsilon^*(t_1)\upsilon(t_2)$ 为实数，从而 $\Phi_M=0$；此时传播子退化为正文常用的交换算符 $S_Q(\gq)$。

\begin{derivation}[从\eqnrefs{eq:qpinem-single-mode-hamiltonian-dp68}到\eqnrefs{eq:qpinem-magnus-exact-dp68}]
先把“时间有序指数为什么能够求闭式”化成一个纯算符代数问题。定义
\[
\hat B\equiv\hat b^\dagger\hat a,
\qquad
\hat B^\dagger=\hat b\hat a^\dagger.
\]
电子平移算符与光子算符作用在张量积 Hilbert 空间的不同因子上，故彼此对易；理想双边电子格又满足 $\hat b^\dagger\hat b=\hat b\hat b^\dagger=I_e$。于是
\begin{align*}
\hat B\hat B^\dagger
&=\hat b^\dagger\hat b\,\hat a\hat a^\dagger
=I_e(\hat N_{\rm ph}+1),\\
\hat B^\dagger\hat B
&=\hat b\hat b^\dagger\,\hat a^\dagger\hat a
=I_e\hat N_{\rm ph},
\end{align*}
从而
\begin{equation*}
[\hat B,\hat B^\dagger]=I_e\otimes I_{\rm ph}\equiv I.
\end{equation*}
因此复合交换算符 $\hat B$、$\hat B^\dagger$ 在这个有效 Hilbert 空间中形成与单个玻色模相同的 Heisenberg--Weyl 对易代数；这正是闭式存在的根源。

传播子满足 $\dot U=\hat A(t)U$，其中
\begin{equation*}
\hat A(t)=-\frac{\ii}{\hbar}\hat H_Q(t)
=\upsilon(t)\hat B-\upsilon^*(t)\hat B^\dagger.
\end{equation*}
Magnus 展开定义
\begin{align*}
\Omega_1&=\int_{t_i}^{t_f}\dd t_1\,\hat A(t_1),\\
\Omega_2&=\frac12\int_{t_i}^{t_f}\dd t_1
\int_{t_i}^{t_1}\dd t_2\,[\hat A(t_1),\hat A(t_2)],
\end{align*}
而 $\Omega_3$ 及更高项只含形如 $[\hat A(t_1),[\hat A(t_2),\hat A(t_3)]]$ 的嵌套对易子。先显式计算二时刻对易子：
\begin{align*}
[\hat A_1,\hat A_2]
={}&-\upsilon_1\upsilon_2^*[\hat B,\hat B^\dagger]
-\upsilon_1^*\upsilon_2[\hat B^\dagger,\hat B]\\
={}&\upsilon_1^*\upsilon_2-\upsilon_1\upsilon_2^*.
\end{align*}
右端乘的是单位算符，所以对任意 $t_3$ 都有
\[
[\hat A(t_3),[\hat A(t_1),\hat A(t_2)]]=0.
\]
所有含三重及更高嵌套对易子的 Magnus 项因此同时消失，而不是逐阶近似舍去。

第一项直接积分为
\begin{align*}
\Omega_1
&=\left(\int_{t_i}^{t_f}\dd t\,\upsilon(t)\right)\hat B
-\left(\int_{t_i}^{t_f}\dd t\,\upsilon^*(t)\right)\hat B^\dagger\\
&=\gq\hat B-\gq^*\hat B^\dagger.
\end{align*}
第二项则为
\begin{align*}
\Omega_2
&=\frac12\int_{t_i}^{t_f}\dd t_1
\int_{t_i}^{t_1}\dd t_2\,
[\upsilon^*(t_1)\upsilon(t_2)-\upsilon(t_1)\upsilon^*(t_2)]I\\
&=\ii\Phi_M I.
\end{align*}
被积括号等于 $2\ii\,\operatorname{Im}[\upsilon^*(t_1)\upsilon(t_2)]$，所以 $\Phi_M$ 为实数，$\ee^{\ii\Phi_M}$ 确为纯整体相位。由于 $[\Omega_1,\Omega_2]=0$，
\begin{align*}
U_Q(t_f,t_i)
&=\exp(\Omega_1+\Omega_2)\\
&=\ee^{\ii\Phi_M}
\exp[\gq\hat b^\dagger\hat a-\gq^*\hat b\hat a^\dagger],
\end{align*}
即正文\eqnrefs{eq:qpinem-magnus-exact-dp68}。
\end{derivation}

\subsection{Fock 响应}

守恒量已经把联合态限制在一条 $\ell+m=n$ 的格线上。对初始电子中心格点和光子 Fock 态 $|n\rangle$，该有效模型内的完整输出写成
\begin{equation}
\boxed{
|0,n\rangle\xrightarrow{S_Q}
\sum_{\ell=-\infty}^{n}
\mathcal A_{\ell|n}(\gq)|\ell,n-\ell\rangle.}
\label{eq:qpinem_joint_fock}
\end{equation}
由于电子相位基底把 $\hat b$ 对角化，$S_Q$ 在每个电子相位扇区上等价于一个光子位移算符。结果可写成单模位移算符矩阵元
\begin{equation}
\boxed{
\mathcal A_{\ell|n}(\gq)
=\langle n-\ell|D(-\gq^*)|n\rangle.}
\label{eq:A_l_n_displacement}
\end{equation}
把位移算符矩阵元展开后，令 $r=|\gq|^2$，可把正文\eqnrefs{eq:A_l_n_displacement} 写成显式的分段闭式
\begin{equation}
\boxed{
\mathcal A_{\ell|n}=\ee^{-r/2}
\begin{cases}
\displaystyle
\sqrt{\frac{(n-\ell)!}{n!}}\,\gq^{\ell}
L_{n-\ell}^{(\ell)}(r), & 0\le\ell\le n,\\[8pt]
\displaystyle
\sqrt{\frac{n!}{(n-\ell)!}}\,(-\gq^*)^{-\ell}
L_n^{(-\ell)}(r), & \ell<0,
\end{cases}}
\label{eq:qpinem-fock-laguerre-dp68}
\end{equation}
其中 $L_j^{(k)}$ 是缔合 Laguerre 多项式。第一支有上界 $\ell\le n$，因为电子净增益 $\ell$ 个量子必须从初始 Fock 态中湮灭同样数目的光子；电子净损失 $\ell<0$ 则可以任意地产生光子。因而在该单模双边交换模型内，完整 Fock 响应由 Laguerre 多项式控制；Bessel 分布只在相干态经典场极限中恢复。



当孤立玻色模式取为近场光子模式，电子动量步长由该模式的主导纵向 Fourier 分量选出，并且电子双边格误差满足\eqnrefs{eq:quantum_bilateral_ladder_dynamical_error} 时，\eqnrange{eq:qpinem-single-mode-hamiltonian-dp68}{eq:qpinem-fock-laguerre-dp68} 就构成受控的自由电子--单光子模式 Fock 交换模型。它的实验输入是模式频率、归一化场分布、有限作用时间和初始光子统计；模式重叠、逐边反冲可分辨或电子谱边界被占据时，应返回多模连续谱或有限反冲联合格。



弱耦合时，令 $r_{\rm ph}\in\mathbb N_+$ 表示电子净增益阶数；第 $r_{\rm ph}$ 阶增益需要从光场中湮灭 $r_{\rm ph}$ 个光子，故最低阶概率正比于同阶光子数下降阶乘：
\begin{equation}
\boxed{
P_{+r_{\rm ph}}
=\frac{|\gq|^{2r_{\rm ph}}}{(r_{\rm ph}!)^2}
\left\langle (\hat a^\dagger)^{r_{\rm ph}} \hat a^{r_{\rm ph}}\right\rangle
+O(|\gq|^{2r_{\rm ph}+2}).}
\label{eq:qpinem-weak-factorial-moment-dp19}
\end{equation}


同一展开在 $r_{\rm ph}=1$ 时给
$P_{+1}=|\gq|^2\langle \hat a^\dagger\hat a\rangle+O(|\gq|^4)$，所以
\begin{equation*}
\frac{(r_{\rm ph}!)^2P_{+r_{\rm ph}}}{P_{+1}^{r_{\rm ph}}}
=\frac{\langle(\hat a^\dagger)^{r_{\rm ph}}\hat a^{r_{\rm ph}}\rangle}
{\langle \hat a^\dagger\hat a\rangle^{r_{\rm ph}}}+O(|\gq|^2)
=g_{\rm ph}^{(r_{\rm ph})}(0)+O(|\gq|^2).
\end{equation*}
这建立了“电子边带作为光子统计探测器”的概率关系：光子数分布先经条件响应核映射为理想电子边带概率；有限探测窗口、效率与读出噪声随后由测量章定义的 POVM 或量子仪器作用在这些理想边带上。

% -----------------------------------------------------------------------------
% -----------------------------------------------------------------------------

\begin{exbox}[QPINEM 真空损失]
量子交换传播子~\eqnrefs{eq:qpinem_S} 对 Fock 输入的弱交换条件为 $|g_{\mathrm Q}|^2(n+1)\ll1$。在实际占据范围内满足该条件后，展开为
\begin{equation*}
 S_Q
 =1+g_{\mathrm Q}\hat b^\dagger\hat a
 -g_{\mathrm Q}^*\hat b\hat a^\dagger+O(|g_{\mathrm Q}|^2).
\end{equation*}
对初态 $|0,n\rangle$，第一项保持电子能量不变；第二项湮灭一个光子并把电子推到 $\ell=+1$；第三项产生一个光子并把电子推到 $\ell=-1$。因此
\begin{align}
 P_{+1|n}&=|g_{\mathrm Q}|^2n+O(|g_{\mathrm Q}|^4),\notag\\
 P_{-1|n}&=|g_{\mathrm Q}|^2(n+1)+O(|g_{\mathrm Q}|^4).
 \label{eq:qpinem-gain-loss-n-nplus1-dp65}
\end{align}
这里的 $+1$ 是电子\emph{增益}一个光子能量，$-1$ 是电子\emph{损失}一个光子能量。两者的差
\begin{equation}
 \boxed{P_{-1|n}-P_{+1|n}=|g_{\mathrm Q}|^2+O(|g_{\mathrm Q}|^4)}
 \label{eq:qpinem-spontaneous-difference-dp65}
\end{equation}
与 $n$ 无关，正是玻色对易关系 $[\hat a,\hat a^\dagger]=1$ 留下的真空项。

例如取 $|g_{\mathrm Q}|=0.10$：真空 $n=0$ 时
\begin{equation*}
 P_{+1|0}=0,
 \qquad
 P_{-1|0}\simeq0.010.
\end{equation*}
所以“没有入射光子”并不意味着电子不能损失能量；电子仍可向真空模式自发发射一个量子。对 $n=1$，增益和损失概率分别约为 $0.010$ 与 $0.020$；对 $n=10$，则约为 $0.10$ 与 $0.11$。只有同时维持 $|g_{\mathrm Q}|^2(n+1)\ll1$ 时，增大 $n$ 才使这两条弱交换率的相对差异按 $1/n$ 减小。这只检验增益--损失对称，不能据此把 Fock 输入的电子约化态视为经典相干调制后的纯态。


\medskip

\eqnrefs{eq:qpinem_S} 给出单模 QPINEM 的完整联合散射算符。在 $|g_{\mathrm Q}|\ll1$ 的弱交换极限，散射算符对初态 $|0,n\rangle$ 的一阶振幅只含一个光子升降算符，因此 Fock 梯矩阵元直接给出 $\sqrt n$ 与 $\sqrt{n+1}$ 的玻色增强。概率保留到 $O(|g_{\mathrm Q}|^2)$ 即得到\eqnrefs{eq:qpinem-gain-loss-n-nplus1-dp65}；电子双边平移仍要求实际占据通道远离谱下界。


由指数展开
\begin{equation*}
 S_Q=1+X+O(X^2),
 \qquad
 X=g_{\mathrm Q}\hat b^\dagger\hat a-g_{\mathrm Q}^*\hat b\hat a^\dagger.
\end{equation*}
作用在 $|0,n\rangle$ 上：
\begin{align*}
 \hat b^\dagger\hat a|0,n\rangle
 &=\sqrt n\,|+1,n-1\rangle,\\
 \hat b\hat a^\dagger|0,n\rangle
 &=\sqrt{n+1}\,|-1,n+1\rangle.
\end{align*}
因此
\begin{equation*}
 S_Q|0,n\rangle
 =|0,n\rangle
 +g_{\mathrm Q}\sqrt n\,|+1,n-1\rangle
 -g_{\mathrm Q}^*\sqrt{n+1}\,|-1,n+1\rangle
 +O(|g_{\mathrm Q}|^2).
\end{equation*}
两个末态的电子边带和光子数都互相正交，故到这一阶不存在交叉项。振幅取模平方即得
\begin{equation*}
 P_{+1|n}=|g_{\mathrm Q}|^2n,
 \qquad
 P_{-1|n}=|g_{\mathrm Q}|^2(n+1).
\end{equation*}
相减后只剩 $|g_{\mathrm Q}|^2$。从算符层面看，这个 $1$ 来自
\begin{equation*}
 \langle n|\hat a\hat a^\dagger|n\rangle
 =\langle n|(\hat a^\dagger\hat a+1)|n\rangle=n+1,
\end{equation*}
所以它不能由经典复振幅替代；经典场极限只是在 $n\gg1$ 时把相对真空修正压到 $1/n$。
\end{exbox}



\begin{exbox}[相同平均光子数的不同电子谱]
现在固定三种光场都具有同一平均光子数
\begin{equation*}
 \bar n=\langle\hat n\rangle=2,
\end{equation*}
但分别取 Fock 态 $|2\rangle$、平均数为 2 的相干态和平均数为 2 的单模热态。它们的一阶平均完全相同，因此在 $|g_{\mathrm Q}|=0.10$ 的弱耦合下都有
\begin{equation*}
 P_{+1}\simeq|g_{\mathrm Q}|^2\bar n=0.020.
\end{equation*}
但二阶增益边带由\eqnrefs{eq:qpinem-weak-factorial-moment-dp19} 控制：
\begin{equation*}
 P_{+2}
 \simeq\frac{|g_{\mathrm Q}|^4}{4}
 \langle\hat n(\hat n-1)\rangle.
\end{equation*}
三种光场的结果是
\begin{center}
\begin{tabular}{c|c|c|c}
 光场 & $\langle\hat n(\hat n-1)\rangle$ & $P_{+2}$ & $g^{(2)}(0)$\\\hline
 $|2\rangle$ & 2 & $5.0\times10^{-5}$ & $1/2$\\
 相干态，$|\alpha|^2=2$ & 4 & $1.0\times10^{-4}$ & $1$\\
 单模热态，$\bar n=2$ & 8 & $2.0\times10^{-4}$ & $2$
\end{tabular}
\end{center}
也就是说，三种场的平均光子数完全相同，电子的一阶增益概率也相同，但二阶边带相差四倍。用只包含电子计数的组合
\begin{equation}
 \boxed{
 \frac{4P_{+2}}{P_{+1}^2}
 =g^{(2)}(0)+O(|g_{\mathrm Q}|^2)}
 \label{eq:qpinem-g2-electron-ratio-dp65}
\end{equation}
即可在弱耦合极限中区分亚 Poisson、Poisson 和 bunching 光子统计。这个算例说明 QPINEM 的真正新增信息不是“又出现一套边带”，而是高阶电子边带把光场的下降阶乘矩逐阶映射到电子能谱。
\end{exbox}


下降阶乘矩来自净增益 $r_{\rm ph}$ 格所要求的 $r_{\rm ph}$ 次光子湮灭；弱耦合只用于截断更高阶的成对往返交换。

\begin{derivation}[从\eqnrefs{eq:qpinem_S}到\eqnrefs{eq:qpinem-fock-laguerre-dp68}]
先解释为什么电子平移问题会变成光子的位移算符问题。定义电子相位态
\begin{equation*}
 |\theta\rangle=\frac{1}{\sqrt{2\pi}}
 \sum_{\ell\in\mathbb Z}\ee^{\ii\ell\theta}|\ell\rangle,
 \qquad
 \hat b|\theta\rangle=\ee^{\ii\theta}|\theta\rangle .
\end{equation*}
于是固定 $\theta$ 后，交换算符中的电子算符退化为复数：
\begin{align*}
 S_Q(\gq)|\theta\rangle|n\rangle
 &=|\theta\rangle
 \exp\!\left[
 \gq \ee^{-\ii\theta}\hat a
 -\gq^*\ee^{+\ii\theta}\hat a^\dagger
 \right]|n\rangle\\
 &=|\theta\rangle
 D\!\left(-\gq^*\ee^{+\ii\theta}\right)|n\rangle ,
\end{align*}
其中 $D(\alpha)=\exp(\alpha\hat a^\dagger-\alpha^*\hat a)$。将末态电子投影到 $|\ell\rangle$，相位积分只提取位移矩阵元中 Fourier 次数为 $\ell$ 的分量，得到正文\eqnrefs{eq:A_l_n_displacement}。真正需要计算的是一般的 $\langle m|D(\alpha)|n\rangle$。

因为 $[\alpha\hat a^\dagger,-\alpha^*\hat a]=|\alpha|^2$ 是 $c$ 数，Baker--Campbell--Hausdorff 公式在这里精确给出
\begin{equation*}
 D(\alpha)
 =\ee^{-|\alpha|^2/2}
 \ee^{\alpha\hat a^\dagger}\ee^{-\alpha^*\hat a}.
\end{equation*}
先让右边的湮灭指数作用在 $|n\rangle$：
\begin{equation*}
 \ee^{-\alpha^*\hat a}|n\rangle
 =\sum_{k=0}^{n}
 \frac{(-\alpha^*)^k}{k!}
 \sqrt{\frac{n!}{(n-k)!}}\,|n-k\rangle .
\end{equation*}
若 $m\ge n$，随后产生指数必须再产生
$j=m-n+k$ 个光子才可能投影到 $\langle m|$，因此
\begin{align*}
 \langle m|D(\alpha)|n\rangle
 &=\ee^{-|\alpha|^2/2}\sqrt{m!n!}
 \sum_{k=0}^{n}
 \frac{(-\alpha^*)^k\alpha^{m-n+k}}
 {k!(m-n+k)!(n-k)!}\\
 &=\ee^{-|\alpha|^2/2}
 \sqrt{\frac{n!}{m!}}\,
 \alpha^{m-n}L_n^{(m-n)}(|\alpha|^2).
\end{align*}
最后一步只用了缔合 Laguerre 多项式的有限和定义
\begin{equation*}
 L_n^{(r)}(x)
 =\sum_{k=0}^{n}
 \frac{(n+r)!}{(n-k)!(r+k)!k!}\,
 \frac{(-x)^k}{n!}.
\end{equation*}
对 $m\le n$ 可重复同一计算，或利用
$D^\dagger(\alpha)=D(-\alpha)$ 取厄米共轭，得到
\begin{equation*}
 \langle m|D(\alpha)|n\rangle
 =\ee^{-|\alpha|^2/2}
 \sqrt{\frac{m!}{n!}}\,
 (-\alpha^*)^{n-m}L_m^{(n-m)}(|\alpha|^2).
\end{equation*}
现在代入 $m=n-\ell$、$\alpha=-\gq^*$。当 $0\le\ell\le n$ 时属于 $m\le n$ 支，且 $-\alpha^*=\gq$；当 $\ell<0$ 时属于 $m>n$ 支，且 $\alpha=-\gq^*$。两支分别给出正文\eqnrefs{eq:qpinem-fock-laguerre-dp68}。
\end{derivation}
电子净增益满足 $\ell\le n$，因为净吸收 $\ell$ 个光子要求 $\hat a^\ell|n\rangle\neq0$；净损失对应产生光子，没有相同的上界。

\begin{derivation}[从\eqnrefs{eq:qpinem_S}到\eqnrefs{eq:qpinem-weak-factorial-moment-dp19}]
令
\[
\hat B=\hat b^\dagger\hat a,
\qquad [\hat B,\hat B^\dagger]=I,
\qquad r\equiv|\gq|^2.
\]
由 BCH 公式，交换算符可精确正规排序为
\begin{equation*}
S_Q(\gq)
=\ee^{-r/2}
\exp(-\gq^*\hat B^\dagger)
\exp(\gq\hat B).
\end{equation*}
这比直接把整个指数按总阶数展开更适合判断“电子净移动了多少格”。对电子初态 $|0\rangle$ 与末态 $|r_{\rm ph}\rangle$，定义只作用在光场上的条件振幅算符
\[
\hat M_{+r_{\rm ph}}
\equiv{}_e\!\langle r_{\rm ph}|S_Q|0\rangle_e.
\]
把两个指数分别展开：
\begin{align*}
\hat M_{+r_{\rm ph}}
=\ee^{-r/2}
\sum_{j,k\ge0}
\frac{(-\gq^*)^j\gq^k}{j!k!}
{}_e\!\langle r_{\rm ph}|
\hat b^j(\hat b^\dagger)^k|0\rangle_e
(\hat a^\dagger)^j\hat a^k.
\end{align*}
双边格上 $(\hat b^\dagger)^k$ 使电子上移 $k$ 格，而 $\hat b^j$ 再使它下移 $j$ 格，因此电子矩阵元严格给出整数选择条件
\[
k-j=r_{\rm ph}.
\]
令 $k=r_{\rm ph}+j$，得到一个对任意 $\gq$ 都成立的条件算符级数
\begin{equation*}
\hat M_{+r_{\rm ph}}
=\ee^{-r/2}
\sum_{j=0}^{\infty}
\frac{(-\gq^*)^j\gq^{r_{\rm ph}+j}}
{j!(r_{\rm ph}+j)!}
(\hat a^\dagger)^j\hat a^{r_{\rm ph}+j}.
\end{equation*}
最低项 $j=0$ 恰好要求连续湮灭 $r_{\rm ph}$ 个光子；下一项必须同时多一次电子下移与一次电子上移，所以总耦合阶数增加二，而不是一。再把 $\ee^{-r/2}=1+O(|\gq|^2)$ 展开，得到
\begin{equation*}
\hat M_{+r_{\rm ph}}
=\frac{\gq^{r_{\rm ph}}}{r_{\rm ph}!}\hat a^{r_{\rm ph}}
+O(|\gq|^{r_{\rm ph}+2}).
\end{equation*}
因此对任意光场初态 $\rho_{\rm ph}$，
\begin{align*}
P_{+r_{\rm ph}}
&=\Tr\!\left[
\hat M_{+r_{\rm ph}}\rho_{\rm ph}
\hat M_{+r_{\rm ph}}^\dagger\right]\\
&=\frac{|\gq|^{2r_{\rm ph}}}{(r_{\rm ph}!)^2}
\Tr\!\left[
\hat a^{r_{\rm ph}}\rho_{\rm ph}
(\hat a^\dagger)^{r_{\rm ph}}\right]
+O(|\gq|^{2r_{\rm ph}+2})\\
&=\frac{|\gq|^{2r_{\rm ph}}}{(r_{\rm ph}!)^2}
\left\langle(\hat a^\dagger)^{r_{\rm ph}}\hat a^{r_{\rm ph}}\right\rangle
+O(|\gq|^{2r_{\rm ph}+2}),
\end{align*}
即正文\eqnrefs{eq:qpinem-weak-factorial-moment-dp19}。
\end{derivation}
\subsection{相干光照下，电子何时仍可视为纯态}
\label{sec:electron-quantum-exchange-coherence}

经典 PINEM 给出一个纯电子态。若把同一束相干光完全量子化，电子输出是否仍等于这个纯态，需要对光场求偏迹才能判断。相干态满足 $\hat a|\alpha_{\rm coh}\rangle=\alpha_{\rm coh}|\alpha_{\rm coh}\rangle$，吸收项因此产生熟悉的经典振幅；发射项含有 $\hat a^\dagger$，还会产生不能用这个数代替的分量。下面保留两项，计算电子的密度矩阵、能谱和纯度。光子统计在电子谱中的表现由同一联合态描述。\cite{Dahan2021Science}

在已建立的单模、双边平移、低反冲交换模型中，去掉不影响观测的整体 Magnus 相位，令
\begin{equation*}
 S_Q=\exp(\gq\hat b^\dagger\hat a-\gq^*\hat b\hat a^\dagger),
 \qquad r_Q=|\gq|^2,\qquad \beta=\alpha_{\rm coh}\gq.
\end{equation*}
$r_Q$ 是无量纲单电子真空交换强度；它不同于由宏观场增强的 $|\beta|^2$。对电子中心通道 $|0\rangle_e$ 和相干光输入，定义经典 PINEM 纯态
\begin{equation*}
 |\psi_\beta\rangle=S_{\rm tr}(\beta)|0\rangle
 =\sum_\ell J_\ell(2|\beta|)\ee^{\ii\ell\arg\beta}|\ell\rangle.
\end{equation*}
精确约化电子态可写成一组向损失方向平移的纯态混合：
\begin{equation}
 \rho_e^{\rm out}=\sum_{m=0}^{\infty}w_m\,
 \hat b^m|\psi_\beta\rangle\langle\psi_\beta|(\hat b^\dagger)^m,
 \qquad w_m=\ee^{-r_Q}\frac{r_Q^m}{m!}.
 \label{eq:electron-coherent-light-reduced-state}
\end{equation}
每个 $\hat b^m|\psi_\beta\rangle$ 都是把经典输出向能量损失方向移动 $m$ 阶后的归一态，$w_m$ 是它的概率权重。局部推导中会把光的平均相干幅度移去，余下的不同光子数对应这些不同的电子平移。只观测电子时，对正交的光子末态求和，便得到上式的混合。

能谱、平均能量交换和方差随即成为
\begin{equation}
 \begin{aligned}
 P_\ell&=\sum_{m=0}^{\infty}\ee^{-r_Q}\frac{r_Q^m}{m!}
 J_{\ell+m}^2(2|\beta|),\\
 \langle\ell\rangle&=-r_Q,\qquad
 \operatorname{Var}\ell=2|\beta|^2+r_Q.
 \end{aligned}
 \label{eq:electron-coherent-light-spectrum-moments}
\end{equation}
负的均值表示电子向光场净提供能量 $r_Q\hbar\omega$。经典场贡献给出对称展宽，真空项给出额外的损失偏置；这两项在一个概率卷积中同时出现。

为了量化混合程度，使用\emph{纯度} $\operatorname{Tr}\rho^2$。若归一化密度矩阵的本征值为 $w_m$，纯度就是 $\sum_mw_m^2$；只有一个本征值等于一时为纯态，几个非零权重共同存在时纯度小于一。这里的平移态彼此正交：$S_{\rm tr}$ 与 $\hat b$ 对易，而 $\langle0|\hat b^{m-n}|0\rangle=\delta_{mn}$。因此已经找到了本征值 $w_m$，无需另求矩阵对角化，得到
\begin{equation}
 \operatorname{Tr}[(\rho_e^{\rm out})^2]
 =\sum_{m=0}^{\infty}w_m^2
 =\ee^{-2r_Q}I_0(2r_Q),
 \label{eq:electron-coherent-light-purity}
\end{equation}
其中 $I_0(x)=\sum_{m\ge0}(x^2/4)^m/(m!)^2$ 为零阶第一类修正 Bessel 函数。它与振荡型的 $J_0$ 不同。这个纯度在本模型中只依赖 $r_Q$；增大相干光幅度会增加经典调制，却不会自动消除电子留下的真空量子记录。

\begin{exbox}[经典纯态近似的误差怎样随耦合缩小]
以 $|\psi_\beta\rangle$ 为经典参考态，求量子电子态偏离它的大小，并比较“能谱接近”与“量子态接近”两种判据。

前面已证明各平移态彼此正交；未平移项的本征值为 $w_0=\ee^{-r_Q}$。因此与参考纯态的重叠概率为
\begin{equation*}
 F_\beta\equiv\langle\psi_\beta|\rho_e^{\rm out}|\psi_\beta\rangle
 =\ee^{-r_Q}.
\end{equation*}
采用\emph{迹距离} $\mathcal D_{\rm tr}(\rho,\sigma)=\tfrac12\|\rho-\sigma\|_1$，其中迹范数是奇异值之和；对厄米差矩阵，就是本征值绝对值之和。这里差矩阵的本征值是 $w_0-1,w_1,w_2,\ldots$，故无需做无限维对角化：
\begin{equation*}
 \mathcal D_{\rm tr}
 =\frac12\bigl[(1-w_0)+\sum_{m\ge1}w_m\bigr]
 =1-\ee^{-r_Q}.
\end{equation*}
迹距离控制任意一个测量事件的概率偏差。要求该误差不超过 $0<\varepsilon<1$，等价于
\begin{equation*}
 r_Q\le-\ln(1-\varepsilon).
\end{equation*}
小耦合时，$\mathcal D_{\rm tr}=r_Q+O(r_Q^2)$，纯度则为
$\operatorname{Tr}[(\rho_e^{\rm out})^2]=1-2r_Q+3r_Q^2+O(r_Q^3)$。

相反，量子方差相对于经典方差的增量在 $\beta\ne0$ 时为
\begin{equation*}
 \frac{\operatorname{Var}\ell-2|\beta|^2}{2|\beta|^2}
 =\frac{r_Q}{2|\beta|^2}
 =\frac1{2|\alpha_{\rm coh}|^2}.
\end{equation*}
增大平均光子数会让这个\emph{相对方差误差}缩小，却不改变固定 $r_Q$ 下的迹距离。要同时恢复完整的经典电子纯态，应取 $\gq\to0$、$|\alpha_{\rm coh}|\to\infty$ 且 $\beta$ 固定。在这个极限中，经典调制保持不变，而真空交换权重趋于零。
\end{exbox}

弱 Fock 交换还需要单独检查玻色增强。对确定光子数 $n$，截断单量子概率需要 $r_Q(n+1)\ll1$；对有宽尾的光子分布，应检查相应高阶光子矩及实际占据尾。于是 $P_{+1}\simeq r_Q\langle n\rangle$ 与高阶阶乘矩公式都不是只由 $r_Q\ll1$ 保证。经典场和量子场两种近似条件应分别报告。

\begin{derivation}[移去光场平均值后，剩下真空记录]
用位移算符 $D(\alpha)=\exp(\alpha\hat a^\dagger-\alpha^*\hat a)$，有 $|\alpha\rangle=D(\alpha)|0\rangle$ 与 $D^\dagger\hat aD=\hat a+\alpha$。因此
\begin{equation*}
 D^\dagger S_QD
 =\exp[\gq\hat b^\dagger\hat a-\gq^*\hat b\hat a^\dagger
       +\beta\hat b^\dagger-\beta^*\hat b]
 =S_Q S_{\rm tr}(\beta).
\end{equation*}
最后一步使用双边平移的 $[\hat b,\hat b^\dagger]=0$，使量子部分与经典部分对易。对光真空作用时，令 $A=-\gq^*\hat b\hat a^\dagger$、$B=\gq\hat b^\dagger\hat a$，则 $[A,B]=r_Q I$，正规排序给
\begin{equation*}
 S_Q|0,0\rangle=\ee^{-r_Q/2}\sum_{m=0}^{\infty}
 \frac{(-\gq^*)^m}{\sqrt{m!}}|-m,m\rangle.
\end{equation*}
移回的 $D(\alpha)$ 只作用于被取迹的光场，不改变电子约化态；不同 $m$ 的光记录正交，因而得到\eqnrefs{eq:electron-coherent-light-reduced-state}。谱中随机整数可写作 $\ell=L_\beta-M$，其中 $P(L_\beta=j)=J_j^2(2|\beta|)$、$M$ 服从均值 $r_Q$ 的 Poisson 分布且二者独立。利用 $\langle L_\beta\rangle=0$、$\operatorname{Var}L_\beta=2|\beta|^2$，就得到\eqnrefs{eq:electron-coherent-light-spectrum-moments}。
\end{derivation}

\subsection{热光输入：由生成函数求完整交换分布}
\label{sec:electron-thermal-exchange-analytic}

相干光的平均光子数确定以后，仍有许多不同光态具有同样的平均能量。热光就是其中一种。它没有固定光学相位，光子数也遵循与相干态不同的分布。为了求这些差别怎样传给电子，保持上一节的单模低反冲交换算符 $S_Q$ 不变，仍取电子初态 $|0\rangle_e$，只把光场输入换为热态
\begin{equation*}
 \rho_{\rm th}=\sum_{n=0}^\infty
 \frac{\bar n_{\rm th}^{\,n}}{(1+\bar n_{\rm th})^{n+1}}|n\rangle\langle n|,
 \qquad \bar n_{\rm th}=\frac1{\exp[\hbar\omega/(k_{\rm B}T_{\rm ph})]-1}.
\end{equation*}
$T_{\rm ph}$ 是该模式的温度，$\bar n_{\rm th}$ 是无量纲平均光子数。热态在光子数基中为对角混合，没有确定的光学参考相位。热光与相干光对自由电子谱的不同作用可由同一 QPINEM 模型比较\cite{Dahan2021Science}。

出射电子可以占据许多边带。求这组概率时，可以利用 Fourier 级数把所有 $P_\ell$ 合成一个函数，称为概率分布的\emph{特征函数}：
\begin{equation*}
 \chi_e(\vartheta)=\sum_{\ell\in\mathbb Z}P_\ell\ee^{\ii\ell\vartheta}
 =\operatorname{Tr}[\rho_e^{\rm out}\ee^{\ii\vartheta\hat\Lambda_e}],
 \qquad \hat\Lambda_e|\ell\rangle=\ell|\ell\rangle.
\end{equation*}
其中 $\hat\Lambda_e$ 是阶数算符，$\vartheta$ 是无量纲 Fourier 变量。它只是用来整理概率，不代表额外的电子动力学。因为 $\partial_\vartheta\ee^{\ii\ell\vartheta}=\ii\ell\ee^{\ii\ell\vartheta}$，在原点求导便有
\[
\langle\ell\rangle=-\ii\chi_e'(0),\qquad
\langle\ell^2\rangle=-\chi_e''(0),\qquad \chi_e(0)=1.
\]
对数的导数则自动去掉低阶矩的乘积：
\[
-\ii(\ln\chi_e)'(0)=\langle\ell\rangle,\qquad
-(\ln\chi_e)''(0)=\langle\ell^2\rangle-\langle\ell\rangle^2.
\]
这些对数导数称为\emph{累积量}，前两阶就是均值和方差。用本节末尾的算符代数求热态平均，得到
\begin{equation}
 \chi_e(\vartheta)=
 \exp\!\left\{\lambda_+(\ee^{\ii\vartheta}-1)
             +\lambda_-(\ee^{-\ii\vartheta}-1)\right\},
 \qquad
 \lambda_+=r_Q\bar n_{\rm th},\quad
 \lambda_-=r_Q(1+\bar n_{\rm th}).
 \label{eq:electron-thermal-characteristic}
\end{equation}
该结果对本交换模型中任意有限 $r_Q=|\gq|^2$ 成立，不是只保留单量子的微扰式。$\lambda_\pm$ 是无量纲统计参数；形式上等同于两个独立 Poisson 变量 $N_\pm$ 的均值，净电子阶数为 $\ell=N_+-N_-$。这是一种\emph{末态分布的等价表示}，不自动意味着真实连续时间动力学是两个无记忆跳跃过程。

直接对 $\ln\chi_e$ 求导得
\begin{equation}
 \langle\ell\rangle=\lambda_+-\lambda_-=-r_Q,\qquad
 \operatorname{Var}\ell=\lambda_++\lambda_-=r_Q(2\bar n_{\rm th}+1).
 \label{eq:electron-thermal-moments}
\end{equation}
热光增强增益和损失两侧的涨落，但净损失偏置仍来自真空项。取 $\bar n_{\rm th}=|\alpha_{\rm coh}|^2$，热光与前节相干光的均值和方差完全相同；要区分两者，必须读更高统计矩、谱形或相干，不能只读展宽。

\begin{exbox}[解析恢复热交换概率并比较四阶累积量]
从特征函数求 $P_\ell$。对 $\ell\ge0$，条件 $N_+-N_-=\ell$ 等价于 $N_+=m+\ell,N_-=m$，故
\begin{equation*}
 P_\ell=\ee^{-(\lambda_++\lambda_-)}
 \sum_{m=0}^\infty\frac{\lambda_+^{m+\ell}\lambda_-^m}{(m+\ell)!m!}.
\end{equation*}
用第一类修正 Bessel 函数的定义
$I_j(2x)=\sum_{m\ge0}x^{2m+j}/[m!(m+j)!]$，并对负 $\ell$ 交换两参数，得到
\begin{equation*}
 P_\ell=\ee^{-(\lambda_++\lambda_-)}
 \left(\frac{\lambda_+}{\lambda_-}\right)^{\ell/2}
 I_{|\ell|}(2\sqrt{\lambda_+\lambda_-}),\qquad \bar n_{\rm th}>0.
\end{equation*}
它称为两个 Poisson 变量之差的 Skellam 分布。零温不能直接把零代入负幂；应回到求和式或特征函数，得到
$P_{-m}=\ee^{-r_Q}r_Q^m/m!$（$m\ge0$），其余增益概率为零。对任意正温度还可直接相除：
\begin{equation*}
 \frac{P_{+\ell}}{P_{-\ell}}
 =\left(\frac{\bar n_{\rm th}}{1+\bar n_{\rm th}}\right)^\ell
 =\ee^{-\ell\hbar\omega/(k_{\rm B}T_{\rm ph})},\qquad\ell>0.
\end{equation*}
这给出了不依赖总体耦合强度的增益--损失比。

再定义四阶累积量 $\mathcal C_4=\langle(\ell-\langle\ell\rangle)^4\rangle-3(\operatorname{Var}\ell)^2$。它是对高斯统计之外结构的一个度量。由\eqnrefs{eq:electron-thermal-characteristic}，
\begin{equation*}
 \mathcal C_4^{\rm(th)}=r_Q(2\bar n_{\rm th}+1).
\end{equation*}
相干光的经典 Bessel 部分具有特征函数 $J_0[4|\beta|\sin(\vartheta/2)]$，真空 Poisson 损失与它独立。将 $\ln J_0$ 展开至 $\vartheta^4$，得到
\begin{equation*}
 \mathcal C_4^{\rm(coh)}=r_Q+2|\beta|^2-6|\beta|^4.
\end{equation*}
在相同平均光子数、$|\beta|^2=r_Q\bar n_{\rm th}$ 下，两种方差相同，四阶累积量却相差 $6r_Q^2\bar n_{\rm th}^2$。这说明为什么只拟合一个高斯宽度会丢失光子统计信息。
\end{exbox}

\begin{derivation}[用交换算符代数完成热态偏迹]
定义 $B=\hat b^\dagger\hat a$。双边平移给 $[B,B^\dagger]=I$，因此可定义位移形式
$D_B(\zeta)=\exp(\zeta B^\dagger-\zeta^*B)$，并写 $S_Q=D_B(-\gq^*)$。
电子计数旋转满足
\begin{equation*}
 \ee^{\ii\vartheta\hat\Lambda_e}S_Q\ee^{-\ii\vartheta\hat\Lambda_e}
 =D_B(-\gq^*\ee^{-\ii\vartheta}).
\end{equation*}
由于初态 $|0\rangle_e$ 的阶数为零，$\chi_e$ 可改写为初态上 $S_Q^\dagger S_Q(\vartheta)$ 的期望。使用
\begin{equation*}
 D_B(\alpha)D_B(\beta)
 =\ee^{(\alpha\beta^*-\alpha^*\beta)/2}D_B(\alpha+\beta)
\end{equation*}
得到
\begin{equation*}
 \chi_e(\vartheta)=\ee^{-\ii r_Q\sin\vartheta}
 \left\langle D_B[-\gq^*(\ee^{-\ii\vartheta}-1)]\right\rangle_{0,\rm th}.
\end{equation*}
接下来显式求平均。正规排序给
$D_B(\zeta)=\ee^{-|\zeta|^2/2}\ee^{\zeta B^\dagger}\ee^{-\zeta^*B}$。
在 $|0,n\rangle$ 上，只有同次数的 $B^\dagger$ 与 $B$ 能返回原电子阶数，因此
\begin{equation*}
 \langle0,n|D_B(\zeta)|0,n\rangle
 =\ee^{-|\zeta|^2/2}\sum_{k=0}^n
 \frac{(-|\zeta|^2)^k}{(k!)^2}\frac{n!}{(n-k)!}
 =\ee^{-|\zeta|^2/2}L_n(|\zeta|^2).
\end{equation*}
$L_n$ 是普通 Laguerre 多项式，上一式的有限和也可作为其定义。令 $q_{\rm th}=\bar n_{\rm th}/(1+\bar n_{\rm th})$，用几何级数加权后的恒等式
\begin{equation*}
 (1-q_{\rm th})\sum_{n=0}^\infty q_{\rm th}^nL_n(x)
 =\exp[-\bar n_{\rm th}x]
\end{equation*}
这个恒等式可由几何级数逐项求导验证：对 $\sum_{n\ge0}q^n=(1-q)^{-1}$ 求 $k$ 阶导数并乘 $q^k$，有
\begin{equation*}
 (1-q)\sum_{n=k}^{\infty}q^n\frac{n!}{(n-k)!}
 =k!\left(\frac{q}{1-q}\right)^k.
\end{equation*}
将它代入 $L_n$ 的有限和，剩下 $\sum_k[-xq/(1-q)]^k/k!$，正是指数级数。于是得到 $\langle D_B(\zeta)\rangle_{0,\rm th}
=\exp[-(\bar n_{\rm th}+1/2)|\zeta|^2]$。
最后代入 $|\zeta|^2=2r_Q(1-\cos\vartheta)$，
\begin{equation*}
 \ln\chi_e=-\ii r_Q\sin\vartheta
 -r_Q(2\bar n_{\rm th}+1)(1-\cos\vartheta),
\end{equation*}
按 $\ee^{\pm\ii\vartheta}$ 重组即得\eqnrefs{eq:electron-thermal-characteristic}。热态对光学相位旋转不变，初始电子又在单一阶数上；交换数守恒遂使电子约化态在阶数基中保持对角。这里的 $P_\ell$ 因而给出了完整约化电子态，而非只给其一部分矩阵元。
\end{derivation}


\subsection{有限反冲下的联合量子格}\label{subsec:qpinem-finite-recoil-lattice-dp46}

\eqnrefs{eq:qpinem_S} 的位移算符闭式还使用了电子能量梯的平移对称性。若电子色散曲率已经使相邻交换边可分辨，则不能把反冲只作为输出相位补在散射之后，而必须把它放回联合电子--光子动力学。对初始光子数为 $n$ 的守恒扇区，取基底
\begin{equation*}
 |\ell;n\rangle\equiv|\ell\rangle_e\otimes|n-\ell\rangle_{\rm ph},
 \qquad \ell\le n,
\end{equation*}
并定义相对于理想均匀能量梯的通道失谐
\begin{equation}
\boxed{
 \delta_\ell
 \equiv\frac{\mathcal E_\ell-E_0-\ell\hbar\omega}{\hbar}.}
\label{eq:qpinem_recoil_channel_detuning_dp46}
\end{equation}
这里 $\mathcal E_\ell=E(p_0+\ell\hbar \kappa_z)$ 始终取真实电子色散；在二阶窄带展开下，若中心同步条件 $v_0\kappa_z=\omega$ 成立，则
\begin{equation}
 \delta_\ell=\ell^2\omega_{\rm r}+O(\ell^3),
 \qquad
 \delta_{\ell+1}-\delta_\ell=(2\ell+1)\omega_{\rm r}+O(\ell^2).
\label{eq:qpinem_recoil_detuning_quadratic_dp46}
\end{equation}

把联合态写成
\begin{equation*}
 |\Psi_n(t)\rangle=\sum_{\ell\le n}c_\ell^{(n)}(t)|\ell;n\rangle,
 \qquad
 \sum_{\ell\le n}|c_\ell^{(n)}|^2=1,
\end{equation*}
则单模交换与真实电子色散共同给出离散 Schr\"odinger 方程
\begin{equation}
\boxed{
 \ii\dot c_\ell^{(n)}
 =\delta_\ell c_\ell^{(n)}
 +\ii\upsilon(t)\sqrt{n-\ell+1}\,c_{\ell-1}^{(n)}
 -\ii\upsilon^*(t)\sqrt{n-\ell}\,c_{\ell+1}^{(n)}.}
\label{eq:qpinem_finite_recoil_coupled_lattice_dp46}
\end{equation}
这是一维非均匀紧束缚链：$\delta_\ell$ 是电子反冲产生的在位能，$\sqrt{n-\ell}$ 是光子 Fock 下界产生的非均匀跃迁矩阵元。二者来源完全不同，因此不能合并成一个经验“有效耦合”。\eqnrefs{eq:qpinem_finite_recoil_coupled_lattice_dp46} 也是数值计算有限反冲 QPINEM 的最小闭合模型；给定 $\upsilon(t)$、$p_0$、$\kappa_z$ 与初始光子统计后即可直接积分得到联合末态。

对实际显著占据的电子窗口 $|\ell|\le N_{\rm occ}$，反冲是否能在作用时间 $T$ 内被分辨可用
\begin{equation}
\boxed{
 \epsilon_{\rm Qrec}^{(T)}
 \equiv\max_{|\ell|\le N_{\rm occ}}
 |\delta_{\ell+1}-\delta_\ell|T
 \simeq(2N_{\rm occ}+1)\omega_{\rm r}T.}
\label{eq:qpinem_recoil_time_control_dp46}
\end{equation}
若还要求不同交换边的失谐小于其局域 Rabi 速率，则对给定光子数扇区定义
\begin{equation}
\boxed{
 \epsilon_{\rm Qrec}^{(\Omega)}(n)
 \equiv
 \max_{|\ell|\le N_{\rm occ}}
 \frac{|\delta_{\ell+1}-\delta_\ell|}
 {|\upsilon|\sqrt{n-\ell}}
 \ll1.}
\label{eq:qpinem_recoil_rabi_control_dp46}
\end{equation}
因此位移算符闭\eqnrefs{eq:qpinem_S} 需要的不只是“单光子反冲能量很小”，而是所有被实际占据的交换边在有限时间和有限耦合速率下都不可分辨。若\eqnrefs{eq:qpinem_recoil_time_control_dp46} 或~\eqnrefs{eq:qpinem_recoil_rabi_control_dp46} 失效，应直接积分\eqnrefs{eq:qpinem_finite_recoil_coupled_lattice_dp46}，而不是继续使用 Laguerre/Bessel 闭式再乘经验修正因子。

在单模、固定纵向动量步长且只保留正能电子通道的有效模型内，对一般光子数混合态 $\rho_{\rm ph}=\sum_n p_n|n\rangle\langle n|$，电子能谱的完整前向模型为
\begin{equation}
\boxed{
 P_\ell
 =\sum_{n=0}^{\infty}p_n\,
 \left|c_\ell^{(n)}(t_f)\right|^2,}
\label{eq:qpinem_finite_recoil_number_forward_dp46}
\end{equation}
其中每一列 $|c_\ell^{(n)}|^2$ 都由同一组物理参数通过\eqnrefs{eq:qpinem_finite_recoil_coupled_lattice_dp46} 计算。低反冲极限下
$c_\ell^{(n)}(t_f)\to\mathcal A_{\ell|n}(\gq)$，\eqnrefs{eq:qpinem_finite_recoil_number_forward_dp46} 随即退化为\secref{subsec:single-mode-number-instrument-dp19}使用的条件响应核。

相干态经典极限也可在同一方程中定量辨认。令光场初态为 $|\alpha_{\rm coh}\rangle$，并取
\begin{equation}
 |\alpha_{\rm coh}|\to\infty,
 \qquad |\gq|\to0,
 \qquad \beta=\alpha_{\rm coh}\gq
 \quad\text{保持有限},
\label{eq:qpinem_coherent_classical_scaling_dp46}
\end{equation}
则 $\hat a=\alpha_{\rm coh}+\delta \hat a$ 使交换生成元分成经典项
$\beta \hat b^\dagger-\beta^*\hat b$ 与涨落项
$\gq \hat b^\dagger\delta \hat a-\gq^*\hat b\delta \hat a^\dagger$。后者的算符系数按 $|\gq|\sim|\alpha_{\rm coh}|^{-1}$ 消失，于是电子约化散射趋向经典 PINEM 算符 $S_{\rm cl}(\beta)$。因此
\begin{equation}
\boxed{
 P_\ell\longrightarrow J_\ell^2(2|\beta|)}
\label{eq:qpinem_to_pinem_probability_limit_dp46}
\end{equation}
这正是固定经典场幅 $\beta$ 时 $|\alpha_{\rm coh}|\to\infty$、$|\gq|\to0$ 的强相干态极限。

电子初态的有限纵向带宽则通过输入密度矩阵进入。令
\begin{equation*}
 \rho_\ee^{\rm in}
 =\int\!\dd p\,\dd p'\,
 \rho_e(p,p')|p\rangle\langle p'|,
\end{equation*}
并用总散射算符 $S$ 定义电子约化输出
\begin{equation}
\boxed{
 \rho_\ee^{\rm out}
 =\operatorname{Tr}_{\rm ph}
 \left[S(\rho_\ee^{\rm in}\otimes\rho_{\rm ph})S^\dagger\right].}
\label{eq:finite_envelope_reduced_output_dp47}
\end{equation}
若探测器只分辨相对于局部中心动量的边带阶数，有限能量分辨函数 $R_E(\Delta E)$ 进一步给出实验概率
\begin{equation}
\boxed{
 P_j^{\rm meas}
 =\int\!\dd E\,R_E(E-E_j)
 \langle E|\rho_\ee^{\rm out}|E\rangle,
 \qquad
 \int\!\dd E\,R_E(E)=1 .}
\label{eq:pinem_detector_convolution_dp47}
\end{equation}
卷积映射把理论边带概率与谱仪像素计数区分开；后续统计反演应使用\eqnrefs{eq:pinem_detector_convolution_dp47} 的卷积核，而不是直接把理想 $P_\ell$ 与实验直方图比较。

有限输入能散 $\sigma_E$ 是否会洗掉相邻边带取决于两个不同条件。分辨边带需要
\begin{equation}
 \epsilon_{\rm res}\equiv\frac{\sigma_E}{\hbar\omega}\ll1,
\end{equation}
而保持输入相干的相位条件由输入密度矩阵非对角宽度 $\sigma_{E,{\rm coh}}$ 决定。二者一般不相等；能谱窄并不自动意味着长相干时间。


选择规则 $\ell+m=n$ 是联合 $U(1)$ 对称性的守恒扇区条件，并非某一微扰阶中的偶然零点。

\begin{derivation}[从\eqnrefs{eq:qpinem_recoil_channel_detuning_dp46}到\eqnrefs{eq:qpinem_finite_recoil_coupled_lattice_dp46}]
先把“固定交换量扇区”写成真正的张量积基底。正文定义的守恒算符为
\begin{equation*}
\mathcal N_{\rm ex}=\Lambda_e+N_{\rm ph},
\end{equation*}
其中 \(\Lambda_e|\ell\rangle_e=\ell|\ell\rangle_e\)，
\(N_{\rm ph}|m\rangle_{\rm ph}=m|m\rangle_{\rm ph}\)。固定本征值 \(n\) 后定义
\begin{equation*}
|\ell;n\rangle
\equiv
|\ell\rangle_e\otimes|n-\ell\rangle_{\rm ph},
\qquad n-\ell\ge0.
\end{equation*}
于是
\begin{equation*}
\mathcal N_{\rm ex}|\ell;n\rangle
=[\ell+(n-\ell)]|\ell;n\rangle
=n|\ell;n\rangle.
\end{equation*}
这个定义把电子“上升一格”和光子“减少一个”写进了同一个基矢，而不是事后再猜平方根因子。

在相对于线性同步梯 \(E_0+\ell\hbar\omega\) 的旋转表象中，真实电子色散留下
\begin{equation*}
H_{\rm rec}
=\hbar\sum_\ell\delta_\ell|\ell\rangle_e\langle\ell|,
\end{equation*}
交换项为
\begin{equation*}
H_Q(t)
=\ii\hbar\left[
\upsilon(t)\hat b^\dagger\hat a
-\upsilon^*(t)\hat b\hat a^\dagger
\right].
\end{equation*}
现在逐项作用。对左邻通道，
\begin{align*}
\hat b^\dagger\hat a|\ell-1;n\rangle
&=\hat b^\dagger\hat a
\left(|\ell-1\rangle_e\otimes|n-\ell+1\rangle_{\rm ph}\right)\\
&=\sqrt{n-\ell+1}\,
|\ell\rangle_e\otimes|n-\ell\rangle_{\rm ph}\\
&=\sqrt{n-\ell+1}\,|\ell;n\rangle.
\end{align*}
对右邻通道，
\begin{align*}
\hat b\hat a^\dagger|\ell+1;n\rangle
&=\hat b\hat a^\dagger
\left(|\ell+1\rangle_e\otimes|n-\ell-1\rangle_{\rm ph}\right)\\
&=\sqrt{n-\ell}\,|\ell;n\rangle.
\end{align*}
因此 \(\sqrt{n-\ell+1}\) 与 \(\sqrt{n-\ell}\) 分别是玻色湮灭、产生算符在实际光子数上的矩阵元；它们不是人为放进耦合常数的修正。

展开固定扇区态
\begin{equation*}
|\Psi_n(t)\rangle
=\sum_{\ell\le n}c_\ell^{(n)}(t)|\ell;n\rangle
\end{equation*}
并代入
\(\ii\hbar\partial_t|\Psi_n\rangle=(H_{\rm rec}+H_Q)|\Psi_n\rangle\)。左乘 \(\langle\ell;n|\) 后，正交性逐项给出
\begin{align*}
\ii\dot c_\ell^{(n)}
={}&\delta_\ell c_\ell^{(n)}
+\ii\upsilon(t)\sqrt{n-\ell+1}\,c_{\ell-1}^{(n)}\\
&-\ii\upsilon^*(t)\sqrt{n-\ell}\,c_{\ell+1}^{(n)},
\end{align*}
即正文\eqnrefs{eq:qpinem_finite_recoil_coupled_lattice_dp46}。

最后把失谐本身接回电子色散。若
\begin{equation*}
\mathcal E_\ell
=E_0+\ell\hbar v_0\kappa_z
+\ell^2\hbar\omega_{\rm r}+R_\ell^{(3)},
\end{equation*}
且中心同步 \(v_0\kappa_z=\omega\)，则
\begin{align*}
\delta_\ell
&\equiv\frac{\mathcal E_\ell-E_0-\ell\hbar\omega}{\hbar}\\
&=\ell^2\omega_{\rm r}+\frac{R_\ell^{(3)}}{\hbar},
\end{align*}
从而得到正文\eqnrefs{eq:qpinem_recoil_detuning_quadratic_dp46}。
\end{derivation}
有限反冲模型中的对角失谐来自电子真实色散，非均匀玻色跃迁系数来自光子数矩阵元；两者是独立的结构来源。

\begin{derivation}[从\eqnrefs{eq:qpinem_coherent_classical_scaling_dp46}到\eqnrefs{eq:qpinem_to_pinem_probability_limit_dp46}]
令光场初态为 \(|\alpha_{\rm coh}\rangle\)。用位移算符把大相干幅度移到真空涨落上，
\begin{equation*}
D_{\rm ph}^\dagger(\alpha_{\rm coh})
\hat a
D_{\rm ph}(\alpha_{\rm coh})
=\hat a+\alpha_{\rm coh}.
\end{equation*}
联合散射指数中的生成元因此精确分成
\begin{align*}
X
={}&\gq\hat b^\dagger(\hat a+\alpha_{\rm coh})
-\gq^*\hat b(\hat a^\dagger+\alpha_{\rm coh}^*)\\
={}&X_{\rm cl}+X_{\rm fl},
\end{align*}
其中
\begin{equation*}
X_{\rm cl}=\beta\hat b^\dagger-\beta^*\hat b,
\qquad
X_{\rm fl}=\gq\hat b^\dagger\hat a-\gq^*\hat b\hat a^\dagger,
\qquad
\beta\equiv\alpha_{\rm coh}\gq.
\end{equation*}
在理想双边电子能量格上
\([\hat b,\hat b^\dagger]=0\)，且电子平移与光子算符彼此对易，因此
\begin{equation*}
[X_{\rm cl},X_{\rm fl}]=0.
\end{equation*}
所以大相干背景与量子涨落不是近似地、而是精确地因子化：
\begin{equation*}
\ee^X
=\ee^{X_{\rm cl}}\ee^{X_{\rm fl}}.
\end{equation*}

现在才取经典场极限
\begin{equation*}
|\alpha_{\rm coh}|\to\infty,
\qquad
|\gq|\to0,
\qquad
\beta=\alpha_{\rm coh}\gq\ \text{固定}.
\end{equation*}
为说明涨落部分确实消失，而不是只比较两个形式上的系数，取任意固定光子数截断 \(P_N\)。双边平移满足 \(\|\hat b\|=\|\hat b^\dagger\|=1\)，而在有限 Fock 子空间上
\begin{equation*}
\|\hat aP_N\|=\sqrt N,
\qquad
\|\hat a^\dagger P_N\|=\sqrt{N+1}.
\end{equation*}
故
\begin{equation*}
\|X_{\rm fl}P_N\|
\le |\gq|\bigl(\sqrt N+\sqrt{N+1}\bigr)
\longrightarrow0.
\end{equation*}
对每个固定的有限涨落扇区，\(\ee^{X_{\rm fl}}\) 因而强收敛到恒等算符；保留下来的电子传播子是
\begin{equation*}
S_{\rm cl}(\beta)
=\exp(\beta\hat b^\dagger-\beta^*\hat b).
\end{equation*}

最后求电子边带振幅。利用正文已经建立的双边平移 Fourier 表示或 Jacobi--Anger 生成函数，
\begin{equation*}
\langle\ell|S_{\rm cl}(\beta)|0\rangle
=J_\ell(2|\beta|)\ee^{\ii\ell\arg\beta}.
\end{equation*}
因此
\begin{equation*}
P_\ell
=\left|\langle\ell|S_{\rm cl}(\beta)|0\rangle\right|^2
=J_\ell^2(2|\beta|),
\end{equation*}
即正文\eqnrefs{eq:qpinem_to_pinem_probability_limit_dp46}。在这一极限中，“经典 PINEM”的数学含义是：相干平均场 \(\beta\) 保持有限，而固定涨落扇区上的量子生成元按 \(|\gq|\) 消失。
\end{derivation}


\begin{derivation}[从\eqnrefs{eq:qpinem-single-mode-hamiltonian-dp68}到\eqnrefs{eq:quantum_bilateral_ladder_dynamical_error}]
真实电子能量格存在最低格点 $|\ell_{\min}\rangle$。 令 $s$ 表示物理单边平移算符、$\hat b$ 表示双边延拓后的幺正平移算符，则

\begin{align*}
s|\ell\rangle&=|\ell-1\rangle,
&&\ell>\ell_{\min},\\
s|\ell_{\min}\rangle&=0,\\
\hat b|\ell\rangle&=|\ell-1\rangle,
&&\ell\in\mathbb Z.
\end{align*}
因此在物理单边子空间内
\begin{equation*}
(\hat b-s)\Pi_{\rm phys}
=|\ell_{\min}-1\rangle\langle\ell_{\min}|,
\qquad
(\hat b^\dagger-s^\dagger)\Pi_{\rm phys}=0.
\end{equation*}
双边模型与单边模型的交换哈密顿量分别记为 $H_b(t)$ 与 $H_s(t)$。由
\begin{equation*}
H_Q(t)=\ii\hbar[\upsilon(t)\hat b^\dagger\hat a-\upsilon^*(t)\hat b\hat a^\dagger]
\end{equation*}
可得它们在物理子空间上的差
\begin{align*}
R_{\rm edge}(t)
&\equiv[H_b(t)-H_s(t)]\Pi_{\rm phys}\\
&=-\ii\hbar\upsilon^*(t)
|\ell_{\min}-1\rangle\langle\ell_{\min}|\,\hat a^\dagger .
\end{align*}
整体相位符号不影响范数。对单边参考演化态 $|\Psi_s(t)\rangle$，令
\begin{equation*}
\Pi_{\min}=|\ell_{\min}\rangle\langle\ell_{\min}|,
\qquad
\hat n=\hat a^\dagger\hat a.
\end{equation*}
则
\begin{align*}
\|R_{\rm edge}(t)|\Psi_s(t)\rangle\|^2
&=\hbar^2|\upsilon(t)|^2
\langle\Psi_s|
\Pi_{\min}\hat a\hat a^\dagger\Pi_{\min}
|\Psi_s\rangle\\
&=\hbar^2|\upsilon(t)|^2
\langle\Pi_{\min}(\hat n+1)\rangle_t,
\end{align*}
其中第二步使用 $\hat a\hat a^\dagger=\hat n+1$，并利用电子投影与光子算符作用在不同 Hilbert 空间而彼此对易。两种传播子满足 Duhamel 恒等式
\begin{equation*}
U_b(t_f,t_i)-U_s(t_f,t_i)
=-\frac{\ii}{\hbar}
\int_{t_i}^{t_f}\dd t\,
U_b(t_f,t)R_{\rm edge}(t)U_s(t,t_i).
\end{equation*}
左右传播子均为幺正算符，因此对给定初态取范数并使用三角不等式得到
\begin{align*}
\|(U_b-U_s)|\Psi_i\rangle\|
&\le
\frac1\hbar\int_{t_i}^{t_f}\dd t\,
\|R_{\rm edge}(t)|\Psi_s(t)\rangle\|\\
&=
\int_{t_i}^{t_f}\dd t\,|\upsilon(t)|
\sqrt{\langle\Pi_{\min}(\hat n+1)\rangle_t}.
\end{align*}
这就是\eqnrefs{eq:quantum_bilateral_ladder_dynamical_error}。若参考演化始终限制在 $0\le n\le N$，则算符不等式
\begin{equation*}
\Pi_{\min}(\hat n+1)
\le (N+1)\Pi_{\min}
\end{equation*}
给出
\begin{equation*}
\epsilon_{\rm edge,Q}^{\rm dyn}
\le
\left(\int_{t_i}^{t_f}\dd t\,|\upsilon(t)|\right)
\sqrt{(N+1)P_{\rm edge}},
\end{equation*}
即正文给出的有限光子数上界。
\end{derivation}
\begin{derivation}[从\eqnrefs{eq:qpinem-exchange-number-dp68}到\eqnrefs{eq:qpinem_conservation}]
先从电子格基底本身推出平移算符的对易关系，而不把它当作已知公式。由
\[
\Lambda_e|\ell\rangle=\ell|\ell\rangle,
\qquad
\hat b^\dagger|\ell\rangle=|\ell+1\rangle,
\qquad
\hat b|\ell\rangle=|\ell-1\rangle,
\]
有
\begin{align*}
[\Lambda_e,\hat b^\dagger]|\ell\rangle
&=(\ell+1)|\ell+1\rangle-\ell|\ell+1\rangle
=\hat b^\dagger|\ell\rangle,\\
[\Lambda_e,\hat b]|\ell\rangle
&=(\ell-1)|\ell-1\rangle-\ell|\ell-1\rangle
=-\hat b|\ell\rangle.
\end{align*}
因为这对所有整数 $\ell$ 成立，故
\[
[\Lambda_e,\hat b^\dagger]=\hat b^\dagger,
\qquad [\Lambda_e,\hat b]=-\hat b.
\]
同理，由 $N_{\rm ph}|n\rangle=n|n\rangle$ 与玻色升降关系得到
\[
[N_{\rm ph},\hat a]=-\hat a,
\qquad
[N_{\rm ph},\hat a^\dagger]=\hat a^\dagger.
\]
于是联合交换数
\[
\mathcal N_{\rm ex}=\Lambda_e+N_{\rm ph}
\]
对两个基本交换顶角分别给出
\begin{align*}
[\mathcal N_{\rm ex},\hat b^\dagger\hat a]
&=[\Lambda_e,\hat b^\dagger]\hat a
+\hat b^\dagger[N_{\rm ph},\hat a]
=0,\\
[\mathcal N_{\rm ex},\hat b\hat a^\dagger]
&=[\Lambda_e,\hat b]\hat a^\dagger
+\hat b[N_{\rm ph},\hat a^\dagger]
=0.
\end{align*}
因此对任意时变包络 $\upsilon(t)$ 都有
\[
[\mathcal N_{\rm ex},\hat H_Q(t)]=0.
\]
为了把“瞬时对易”严格传到有限时间传播子，定义 Heisenberg 型量
\[
\mathcal N_U(t)=U_Q^\dagger(t,t_i)\mathcal N_{\rm ex}U_Q(t,t_i).
\]
利用
$\ii\hbar\,\partial_tU_Q=\hat H_QU_Q$
和
$\partial_tU_Q^\dagger=(\ii/\hbar)U_Q^\dagger\hat H_Q$
可得
\begin{align*}
\frac{\dd\mathcal N_U}{\dd t}
&=\frac{\ii}{\hbar}U_Q^\dagger\hat H_Q\mathcal N_{\rm ex}U_Q
-\frac{\ii}{\hbar}U_Q^\dagger\mathcal N_{\rm ex}\hat H_QU_Q\\
&=\frac{\ii}{\hbar}U_Q^\dagger[\hat H_Q,\mathcal N_{\rm ex}]U_Q=0.
\end{align*}
所以 $\mathcal N_U(t)=\mathcal N_{\rm ex}$，等价于
\[
[\mathcal N_{\rm ex},U_Q(t_f,t_i)]=0.
\]
固定相位时 $U_Q$ 除整体相位外就是 $S_Q$，从而得到正文\eqnrefs{eq:qpinem_conservation}。

守恒扇区也可直接从谱投影看出。若初态 $|0,n\rangle$ 满足
$\mathcal N_{\rm ex}|0,n\rangle=n|0,n\rangle$，则
\begin{align*}
\mathcal N_{\rm ex}U_Q|0,n\rangle
=U_Q\mathcal N_{\rm ex}|0,n\rangle
=nU_Q|0,n\rangle.
\end{align*}
所以输出态仍完全位于本征值为 $n$ 的子空间；基底态 $|\ell,m\rangle$ 的必要条件为 $\ell+m=n$。
\end{derivation}



\begin{derivation}[从\eqnrefs{eq:multimode_qpinem_continuous_kernel_dp101}到\eqnrefs{eq:multimode_qpinem_lattice_dp47}]

先把第 $\nu$ 个连续电子耦合算符记为
\[
\hat{\mathcal G}_{\nu,I}(t)
=
\sum_{ss'}
\int\dd^3p\,\dd^3p'\,
\mathcal G_{\nu;s's}(\mathbf p',\mathbf p;t)
|\mathbf p',s'\rangle\langle\mathbf p,s|.
\]
在实际窄带工作窗内，假定主导矩阵元集中在一个自旋保持、近固定动量转移通道
\[
(\mathbf p,s_0)
\longrightarrow
(\mathbf p+\hbar\boldsymbol\kappa_\nu,s_0).
\]
定义参考平移算符
\[
\hat b_\nu^\dagger
=
\int\dd^3p\,
|\mathbf p+\hbar\boldsymbol\kappa_\nu,s_0\rangle
\langle\mathbf p,s_0|
\]
以及投影后的共同速率 $\upsilon_\nu(t)$。于是精确耦合算符在该工作窗上分解成
\[
\hat{\mathcal G}_{\nu,I}(t)
=
\ii\upsilon_\nu(t)\hat b_\nu^\dagger
+
\hat R_{\nu,\kappa}(t).
\]
把所有余项收集为
\[
\hat R_\kappa(t)
=
\hbar
\sum_\nu
\left[
\hat R_{\nu,\kappa}(t)\hat a_\nu
+
\hat R_{\nu,\kappa}^\dagger(t)\hat a_\nu^\dagger
\right].
\]
令 $\hat H_{\rm tr}$ 只保留主导平移项，并由它生成 $\hat U_{\rm tr}$。完整传播子 $\hat U$ 与 $\hat U_{\rm tr}$ 的 Duhamel 恒等式为
\[
(\hat U-\hat U_{\rm tr})\Pi_{\rm in}
=
-\frac{\ii}{\hbar}
\int_{t_i}^{t_f}\dd t\,
\hat U(t_f,t)
\hat R_\kappa(t)
\hat U_{\rm tr}(t,t_i)\Pi_{\rm in}.
\]
完整传播子幺正，因此
\[
\|(\hat U-\hat U_{\rm tr})\Pi_{\rm in}\|
\le
\frac1\hbar
\int_{t_i}^{t_f}\dd t\,
\|
\hat R_\kappa(t)
\hat U_{\rm tr}(t,t_i)\Pi_{\rm in}
\|,
\]
这正是正文\eqnrefs{eq:multimode_momentum_step_projection_error_dp101}。只有该量足够小时，连续核才受控退化为
\[
\hat H_{\rm Q}^{(M)}
=
\ii\hbar\sum_\nu
\left[
\upsilon_\nu\hat b_\nu^\dagger\hat a_\nu
-\upsilon_\nu^*\hat b_\nu\hat a_\nu^\dagger
\right].
\]

设整数向量 $\boldsymbol\ell$ 在实际工作窗中唯一对应电子动量
\[
\mathbf p_{\boldsymbol\ell}
=
\mathbf p_0
+\hbar\sum_\nu\ell_\nu\boldsymbol\kappa_\nu,
\]
并定义
\[
|\boldsymbol\ell;\mathbf n\rangle
=
|\mathbf p_{\boldsymbol\ell}\rangle_e
\otimes
\bigotimes_\nu
|n_\nu-\ell_\nu\rangle_\nu .
\]
对第 $\nu$ 个吸收顶角，
\begin{align*}
\hat b_\nu^\dagger\hat a_\nu
|\boldsymbol\ell;\mathbf n\rangle
&=
\hat b_\nu^\dagger
|\mathbf p_{\boldsymbol\ell}\rangle
\otimes
\sqrt{n_\nu-\ell_\nu}\,
|n_\nu-\ell_\nu-1\rangle_\nu
\otimes_{\mu\neq\nu}
|n_\mu-\ell_\mu\rangle_\mu\\
&=
\sqrt{n_\nu-\ell_\nu}\,
|\boldsymbol\ell+\mathbf e_\nu;\mathbf n\rangle .
\end{align*}
同理，
\[
\hat b_\nu\hat a_\nu^\dagger
|\boldsymbol\ell;\mathbf n\rangle
=
\sqrt{n_\nu-\ell_\nu+1}\,
|\boldsymbol\ell-\mathbf e_\nu;\mathbf n\rangle .
\]

在旋去理想多频能量
$E_0+\sum_\nu\ell_\nu\hbar\omega_\nu$
的表象中，真实电子色散留下对角项
\[
\hat H_{\rm rec}
=
\hbar\sum_{\boldsymbol\ell}
\delta_{\boldsymbol\ell}
|\boldsymbol\ell\rangle\langle\boldsymbol\ell|.
\]
展开
\[
|\Psi_{\mathbf n}(t)\rangle
=
\sum_{\boldsymbol\ell}
c_{\boldsymbol\ell}^{(\mathbf n)}(t)
|\boldsymbol\ell;\mathbf n\rangle
\]
并代入
\[
\ii\hbar\partial_t|\Psi_{\mathbf n}\rangle
=
(\hat H_{\rm rec}+\hat H_{\rm Q}^{(M)})
|\Psi_{\mathbf n}\rangle.
\]
左乘 $\langle\boldsymbol\ell;\mathbf n|$ 后，来自相邻格点的两类项分别为
\[
+\ii\hbar
\upsilon_\nu
\sqrt{n_\nu-\ell_\nu+1}\,
c_{\boldsymbol\ell-\mathbf e_\nu}^{(\mathbf n)}
\]
和
\[
-\ii\hbar
\upsilon_\nu^*
\sqrt{n_\nu-\ell_\nu}\,
c_{\boldsymbol\ell+\mathbf e_\nu}^{(\mathbf n)}.
\]
与对角反冲项合并并除以 $\hbar$，得到正文\eqnrefs{eq:multimode_qpinem_lattice_dp47}。

最后说明多指标标签为什么需要一一对应。若存在
$\boldsymbol\ell\neq\boldsymbol\ell'$ 却有
\[
\sum_\nu
(\ell_\nu-\ell_\nu')
\boldsymbol\kappa_\nu=0,
\]
则
$|\mathbf p_{\boldsymbol\ell}\rangle=
|\mathbf p_{\boldsymbol\ell'}\rangle$，
这两个“格点”并非正交电子态。此时必须先把所有到达同一物理动量的路径振幅相加，再定义独立基底；直接把 $\boldsymbol\ell,\boldsymbol\ell'$ 当成两个末态会重复计算 Hilbert 空间维数。
\end{derivation}

\begin{derivation}[从\eqnrefs{eq:multicolor_classical_scattering_dp47}到\eqnrefs{eq:multicolor_pinem_convolution_dp47}]

在多模相干态经典极限中，每个量子模式的算符被分解成平均幅与量子涨落；固定
\[
\beta_\nu=\alpha_\nu g_{\\mathrm Q,\nu}
\]
并令相对涨落修正消失后，电子散射算符为
\[
S_{\rm cl}^{(M)}
=
\exp\!\left[
\sum_\nu
(\beta_\nu\hat b_\nu^\dagger-\beta_\nu^*\hat b_\nu)
\right].
\]
完整连续动量空间中的平移算符彼此对易，
\[
[\hat b_\mu,\hat b_\nu]=
[\hat b_\mu,\hat b_\nu^\dagger]=0,
\]
所以指数可逐模分解：
\[
S_{\rm cl}^{(M)}
=
\prod_\nu
\exp(
\beta_\nu\hat b_\nu^\dagger
-\beta_\nu^*\hat b_\nu).
\]

对单个模式，记
\[
\beta_\nu=|\beta_\nu|\ee^{\ii\varphi_\nu},
\qquad
\hat T_\nu\equiv\hat b_\nu^\dagger.
\]
在理想双边平移空间中 $\hat T_\nu$ 是幺正算符，
\[
\hat T_\nu^{-1}=\hat b_\nu.
\]
因此任意整数幂统一写作
\[
\hat T_\nu^{r}
=
\begin{cases}
(\hat b_\nu^\dagger)^r,&r\ge0,\\
(\hat b_\nu)^{-r},&r<0.
\end{cases}
\]
取 $\hat T_\nu$ 的相位本征态
\[
\hat T_\nu|\theta_\nu\rangle
=
\ee^{-\ii\theta_\nu}|\theta_\nu\rangle,
\]
则单模生成元在该表象中成为普通复数：
\begin{align*}
\beta_\nu\hat T_\nu-\beta_\nu^*\hat T_\nu^{-1}
&\longrightarrow
|\beta_\nu|
\left[
\ee^{\ii(\varphi_\nu-\theta_\nu)}
-\ee^{-\ii(\varphi_\nu-\theta_\nu)}
\right]\\
&=
2\ii|\beta_\nu|
\sin(\varphi_\nu-\theta_\nu).
\end{align*}
Jacobi--Anger 恒等式给
\[
\exp[
2\ii|\beta_\nu|\sin(\varphi_\nu-\theta_\nu)
]
=
\sum_{r_\nu\in\mathbb Z}
J_{r_\nu}(2|\beta_\nu|)
\ee^{\ii r_\nu\varphi_\nu}
\ee^{-\ii r_\nu\theta_\nu}.
\]
把
$\ee^{-\ii r_\nu\theta_\nu}$
重新提升为算符 $\hat T_\nu^{r_\nu}$，得到严格的整数平移展开
\[
\exp(
\beta_\nu\hat b_\nu^\dagger
-\beta_\nu^*\hat b_\nu)
=
\sum_{r_\nu\in\mathbb Z}
J_{r_\nu}(2|\beta_\nu|)
\ee^{\ii r_\nu\varphi_\nu}
\hat T_\nu^{r_\nu}.
\]

对所有模式相乘：
\[
S_{\rm cl}^{(M)}
=
\sum_{\{r_\nu\}}
\left[
\prod_\nu
J_{r_\nu}(2|\beta_\nu|)
\ee^{\ii r_\nu\varphi_\nu}
\right]
\prod_\nu\hat T_\nu^{r_\nu}.
\]
若
\[
\kappa_\nu=m_\nu\kappa_0,
\qquad
m_\nu\in\mathbb Z,
\]
则
\[
\prod_\nu\hat T_\nu^{r_\nu}
\]
把电子总共平移
\[
\Delta p
=
\hbar\kappa_0
\sum_\nu m_\nu r_\nu.
\]
因此所有满足
\[
\ell=\sum_\nu m_\nu r_\nu
\]
的整数路径都到达同一个电子末态 $|\ell\rangle$，其振幅必须相干求和。于是
\[
A_\ell
=
\sum_{\{r_\nu\}}
\delta_{\ell,\sum_\nu m_\nu r_\nu}
\prod_\nu
J_{r_\nu}(2|\beta_\nu|)
\ee^{\ii r_\nu\varphi_\nu},
\]
即正文\eqnrefs{eq:multicolor_pinem_convolution_dp47}。

若不同颜色的动量步长不满足这种整数公度关系，则一般不存在唯一的一维总阶数 $\ell$；此时末态应继续用真实电子动量或多指标通道标记，不能把不同可分辨末态强行卷积到一条 Bessel 能量格。
\end{derivation}



